\documentclass[11pt,a4paper]{article}
\usepackage[margin=27mm,headheight=14pt]{geometry}
\usepackage[T1]{fontenc}
\usepackage{lmodern,amsmath,amssymb,amsthm,mathtools,mathrsfs}
\usepackage{microtype,booktabs,longtable,array,enumitem,xcolor,graphicx}
\usepackage{xurl}
\usepackage[colorlinks=true,allcolors=blue!45!black]{hyperref}
\usepackage{bookmark}
\setlength{\emergencystretch}{2em}
\setlength{\parskip}{3pt}
\setlist{itemsep=2pt,topsep=4pt}
\allowdisplaybreaks[2]
\newtheorem{theorem}{Theorem}[section]
\newtheorem{proposition}[theorem]{Proposition}
\newtheorem{lemma}[theorem]{Lemma}
\newtheorem{corollary}[theorem]{Corollary}
\newtheorem{counterexample}[theorem]{Counterexample}
\newtheorem{example}[theorem]{Example}
\theoremstyle{definition}
\newtheorem{definition}[theorem]{Definition}
\newtheorem{remark}[theorem]{Remark}
\newcommand{\Sh}{\operatorname{Sh}}
\newcommand{\Hom}{\operatorname{Hom}}
\newcommand{\Map}{\operatorname{Map}}
\newcommand{\Pro}{\operatorname{Pro}}
\newcommand{\colim}{\operatorname*{colim}}
\newcommand{\limm}{\operatorname*{lim}}
\newcommand{\id}{\operatorname{id}}
\newcommand{\Spec}{\operatorname{Spec}}
\newcommand{\im}{\operatorname{im}}
\newcommand{\cl}{\operatorname{cl}}
\newcommand{\R}{\mathbb R}
\newcommand{\C}{\mathbb C}
\newcommand{\Z}{\mathbb Z}
\newcommand{\Q}{\mathbb Q}
\newcommand{\F}{\mathbb F}
\newcommand{\recordurl}{https://github.com/KokunoYumeto/mathematics-workbenches}
\newcommand{\workbenchurl}{https://github.com/KokunoYumeto/homotopy-and-categories-workbench}
\title{Gluing Grothendieck's tubes by their intersections}
\author{}
\date{23 November 2026}
\begin{document}
\maketitle
\begin{abstract}
Grothendieck's manuscript on stratified spaces proposes reconstruction
from tubes around strata and prescribed incidence pieces. A three-stratum
linear example shows why the genuine overlaps are essential. We prove
ordinary sheaf descent for a synchronized open-tube cover and characterize
arbitrary open descent diagrams by their pointwise membership groupoids.
We also classify the stable-graph contraction posets with conditional
binary meets. Higher sheaves require additional coherence: we give exact
shape and hypercover comparisons, finite-set presentation corrections,
and an explicit separation of the information retained by sheaves,
stacks and full exit categories. Under recursive compatible collars and
proper source bundles, a deployment reconstructs the exit category;
under tame hypotheses its neighborhood germs compare with homotopy links.
Explicit TOP, PL and smooth counterexamples show which manifold and
regularity assumptions cannot be omitted. The mathematics and text were produced by GPT-6 Astra (ChatGPT Pro
mode), with independent re-proofs in the stated scopes by separate
instances of the same model and earlier checks by Claude Opus 5.5 at
effort max; no human has checked the work.
\end{abstract}

\setcounter{tocdepth}{1}
\begingroup\small\setlength{\parskip}{0pt}
\makeatletter\renewcommand{\l@section}{\@dottedtocline{1}{0em}{2em}}\makeatother
\tableofcontents\endgroup
\newpage
\section{The reconstruction problem}
This section specifies the reconstruction problem and the information that its solution must preserve. In cote~151, scans~57--61 and~69, Grothendieck describes tubes and incidence pieces associated with closed strata. On PDF page~69 (archive printed footer~68) he states:
\begin{quote}
``Alors $X$ s'identifie au topos $\varinjlim$ de ces topos.''
\end{quote}
Thus the proposed diagram is meant to reconstruct the topos of $X$ as its colimit. The choice of pieces and arrows is decisive. His phrase ``des c\^otes mal taill\'ees'' on PDF page~59 (archive printed footer~58) describes the difficulty: pieces can be individually well chosen while their mutual intersections carry information absent from the prescribed diagram~\cite{Groth}. We use the archive's PDF scan numbering throughout.

Let $I$ be a poset and $(X_i)_{i\in I}$ a locally finite family of closed subsets of a space $X$. Assume
\begin{equation}\label{eq:closed-family}
 X=\bigcup_iX_i,\qquad i\le j\Longrightarrow X_i\subset X_j,
 \qquad X_i\cap X_j=\bigcup_{k\le i,j}X_k.
\end{equation}
Local finiteness means that every point has a neighbourhood meeting only finitely many members of the family. Empty members are allowed. Define
\[
 S_i=X_i\setminus\bigcup_{k<i}X_k,
 \qquad R_i=\bigcup_{k:\ i\nleq k}X_k.
\]
For every simultaneous choice $\alpha=(N_i)_{i\in I}$ of open neighbourhoods $N_i\supset X_i$, put $W_i^\alpha=N_i\setminus R_i$. The use of the order in this definition matters when distinct indices have the same closed subset.

We write $\Sh(Z)$ for sheaves of sets on $Z$, with geometric morphisms pointing in the direction of continuous maps. All colimits of ordinary topoi below are bicategorical colimits with invertible comparison cells. A pro-topos is a cofiltered diagram of topoi; the defining mapping category is
\begin{equation}\label{eq:pro-maps}
 \Hom_{2\Pro}((E_a)_a,(T_b)_b)
 =\mathop{\mathrm{bilim}}_b\mathop{\mathrm{bicolim}}_a\Hom(E_a,T_b).
\end{equation}
This convention specifies the category in which our pro-object assertion is made.

\begin{theorem}[Open-tube descent]\label{thm:tubes}
Assume \eqref{eq:closed-family}. The $W_i^\alpha$ form an open cover of $X$. For $q=0,1,2$ let
\[
 C_q^\alpha=\Sh\!\left(\coprod_{(i_0,\ldots,i_q)\in I^{q+1}}
 W_{i_0}^\alpha\cap\cdots\cap W_{i_q}^\alpha\right).
\]
For each $\alpha$ the augmented truncated \v Cech diagram $C_\bullet^\alpha$ has bicolimit $\Sh(X)$. If each degree is regarded as the single pro-object $C_q=(C_q^\alpha)_\alpha$, then
\[
 \mathop{\mathrm{bicolim}}_{\Delta_{\le2}^{\mathrm{op}}}C_q
 \simeq c\Sh(X).
\]
Here $c$ denotes the constant pro-object. For finite $I$ the same conclusion uses the individual intersections indexed by nonempty subsets of $I$ of cardinality at most three. For infinite $I$ it is the displayed synchronized pro-object that occurs in the statement.
\end{theorem}

In words, one must glue tubes using their intersections, with the usual compatibility on triple intersections. The ordinary descent theorem is the equivalence for sheaves on an open cover with overlap isomorphisms and a triple cocycle \cite[Lemmas~6.33.2 and~6.33.4, Tags~00AL and~00AN]{StacksGluing}. Theorem~\ref{thm:tubes} applies it to a particular open cover and gives the precise pro-object formulation. No Whitney, conical or smooth structure is required for this statement.

For the next assertion fix $d\in\{1,2\}$ and the closed chain
$X_0=\{0\}\subset X_1=\R^d\subset X_2=\R^{2d}$.
Write $\operatorname{germ}_A Z$ for the pro-system of open
neighbourhoods of $A$ in $Z$, and put
\[
\begin{gathered}
 S_0=X_0,\quad S_1=X_1\setminus X_0,\quad S_2=X_2\setminus X_1,\\
 V_{01}=\operatorname{germ}_{X_0}X_1,\quad
 V_{12}=(\operatorname{germ}_{X_1}X_2)\setminus X_0,\quad
 V_{02}=\operatorname{germ}_{X_0}X_2,\\
 V_{01}^*=V_{01}\setminus X_0,\quad
 V_{12}^*=V_{12}\setminus X_1,\quad
 V_{02}^*=V_{02}\setminus X_1.
\end{gathered}
\]
The incidence diagram has the arrows
$S_i\to V_{ij}\leftarrow V_{ij}^*\to S_j$ for $i<j$,
together with $V_{01}\to V_{02}$ and $V_{02}^*\to V_{12}^*$,
and all their composites. Its triple pieces repeat
$V_{02}^*,V_{02}^*,V_{02}$ and add no comparison.
The successor diagram keeps only the pairs $(0,1)$ and $(1,2)$
and their stratum arrows.

\begin{theorem}[Insufficiency of the prescribed incidence pieces]\label{thm:stones}
For the closed chain $\{0\}\subset\R\subset\C$, the incidence diagram just specified admits compatible sheaves which are not the restrictions of any sheaf on $\C$. The same failure occurs for $\{0\}\subset\R^2\subset\R^4$, whose open strata are connected smooth manifolds. The successor-pair diagram already fails for the first chain.
\end{theorem}

The assertion is stronger than failure for badly behaved stratifications: these chains consist of linear subspaces and have the standard conical local models. The failure concerns the diagram itself. In particular, regularity assumptions on the strata cannot supply an intersection which the diagram omitted.

\begin{theorem}[An exact criterion for an open diagram]\label{thm:membership}
Let $X$ be $T_1$, let $C$ be a small category, and let $D:C\to\operatorname{Open}(X)$ be a diagram of inclusions. Write $C_x$ for the full subcategory on objects $c$ such that $x\in D(c)$. The canonical cocone of topoi $\Sh(D(c))\to\Sh(X)$ is a bicolimit if and only if every groupoidification $C_x[\operatorname{Mor}(C_x)^{-1}]$ is equivalent to the terminal groupoid. Equivalently, every nerve $|N(C_x)|$ is nonempty, connected and simply connected.
\end{theorem}

This gives a local test for deleting objects or arrows from any specified open gluing diagram. It is a criterion for ordinary topoi, so it asks for fundamental groupoids, not contractibility of all higher homotopy groups.

\paragraph{Stable graph types.}
Let $g,n$ be nonnegative integers. A stable graph of type $(g,n)$ is a connected finite graph with loops
and multiple edges allowed, a nonnegative genus $g_v$ at each vertex,
and $n$ individually labelled legs. Its valence counts both ends of
a loop and all legs. Stability and total genus mean
\[
 2g_v-2+\operatorname{val}(v)>0,
 \qquad g=\sum_v g_v+|E|-|V|+1.
\]
Contracting a nonloop edge merges its two vertices and adds their
genera; contracting a loop raises its vertex genus by one. Order graph
isomorphism classes by reverse contraction: $H\leq G$ when $H$
contracts to $G$. The labels of the legs are always fixed.

\begin{theorem}[Conditional meets of stable-graph types]\label{thm:graph-meets}
Assume $2g-2+n>0$. Every pair of stable-graph types having a common
refinement has a greatest common refinement in the reverse-contraction
order if and only if
\[
 g=0\ (n\geq3),\qquad
 g=1\ (1\leq n\leq3),\qquad (g,n)=(2,0).
\]
Here ``greatest'' means the common refinement to which every common
refinement contracts. The statement concerns graph types; it does
not identify the incidence category of specified contractions and
automorphisms with an orbit poset.
\end{theorem}

The positive finite cases are proved by complete contraction and meet tables; genus zero has a uniform tree argument, and explicit families cover the negative cases.
The geometric contraction order and its automorphism data are described
by Abramovich, Caporaso and Payne \cite[Sections~3.2--4.3,
pp.~775--780]{ACP}. Wang and Zhou develop M\"obius inversion for a
related stable-graph poset \cite[Section~3.2 and Appendix,
preprint pp.~8--12 and~20--23]{WangZhou}; their legs are unlabelled,
whereas our labels are fixed. Tosteson's genus-zero reduction uses a
coarsened order and need not preserve the original contraction order
\cite[Propositions~3.5 and~3.7, Warning~3.6,
preprint pp.~8--10]{Tosteson}. Theorem~\ref{thm:graph-meets}
asks which of the original labelled posets have conditional meets.

\paragraph{The amount of exit information.}
\begin{theorem}[The truncation level of exodromy]\label{thm:truncation}
Let $X$ be paracompact and locally of singular shape, with a conical
stratification by a poset satisfying the ascending chain condition.
Write $\mathcal E=\operatorname{Exit}(X)$ and let $\operatorname{Cons}_{\le n}(X)$ denote constructible sheaves of $n$-truncated spaces. For $n\geq0$,
\[
 \operatorname{Cons}_{\leq n}(X)
 \simeq\operatorname{Fun}(\mathcal E,\mathcal S_{\leq n})
 \simeq\operatorname{Fun}(h_{n+1}\mathcal E,\mathcal S_{\leq n}).
\]
Here $h_{n+1}$ retains the $n$-truncations of all mapping spaces and
their coherent compositions. Thus ordinary constructible sheaves use
$\pi_0$ of exit mapping spaces; groupoid stacks use their fundamental
groupoids.
\end{theorem}


\paragraph{The higher comparison data.}
Let $A$ be a small ordinary category, $\mathcal X$ a Grothendieck
topos, and $\varphi:A\to\mathcal X$ a functor, with $U_a=\varphi(a)$.
Write $\widehat A=\operatorname{Fun}(A^{\mathrm{op}},\mathbf{Set})$ and
\[
 L:\widehat A\longrightarrow\mathcal X,\qquad
 L(F)=\int^{a\in A}F(a)\cdot U_a ,
\]
where the dot denotes copower. For every finite category $J$ and
$F:J\to\widehat A$, let $c_F:L(\lim_JF)\to\lim_JL(F)$ be the comparison.
Let $\mathscr X$ be either the associated $1$-localic infinity-topos or
its hypercompletion, with that choice fixed in each formula.
Use $\Delta K$ for its constant object associated to a space $K$ and
$\Pi_\infty$ for shape as a prospace. Put
\[
 B_{\varphi,n}=\coprod_{a_0\to\cdots\to a_n}U_{a_0},
 \qquad H_\varphi=|B_\varphi|=\operatorname{colim}_{a\in A}U_a .
\]
Faces compose arrows or apply $\varphi$ at the initial vertex;
degeneracies insert identities. For a finite simplicial set $K$,
$[K,B_\varphi]$ is the object of simplicial maps, computed by its finite
nondegenerate-simplex limit in $\mathcal X$. Write $\operatorname{Sd}$
for barycentric subdivision, with the last-vertex transition maps.

\begin{theorem}[Finite comparisons, shape, and local filling]
\label{thm:higher-comparisons}
With exactly the preceding data, the following assertions hold.
\begin{enumerate}
\item Every $c_F$ is epic if and only if $L$ preserves finite limits.
When these equivalent conditions hold, $H_\varphi\to1$ is an equivalence
both before and after hypercompletion.
\item The canonical shape comparison
\[
 \operatorname{colim}_{a\in A}\Pi_\infty(\mathscr X/U_a)
 \longrightarrow\Pi_\infty(\mathscr X)                         \tag{HS1}
\]
is an equivalence if and only if, for every space $K$,
\[
 \Map_{\mathscr X}(1,\Delta K)\longrightarrow
 \lim_{a\in A^{\mathrm{op}}}\Map_{\mathscr X}(U_a,\Delta K)       \tag{HS2}
\]
is an equivalence.
\item The augmentation $H_\varphi\to1$ becomes an equivalence after
hypercompletion if and only if, for every integer $n\ge0$, the map
\[
 \operatorname*{colim}_{r\ge0}
 [\operatorname{Sd}^{r}\Delta^n,B_\varphi]\longrightarrow
 \operatorname*{colim}_{r\ge0}
 [\operatorname{Sd}^{r}\partial\Delta^n,B_\varphi]               \tag{HF}
\]
is epic in $\mathcal X$. For $n=0$ the target is terminal.
The required number of subdivisions may vary on a covering family.
\end{enumerate}
\end{theorem}
The first clause is stronger than the all-space mapping test for shape.
Ordinary forest comparisons alone do not imply either form of homotopy
descent; Section~\ref{b5.2.-a-finite-counterexample} gives a finite example.
The proofs of the three clauses occur in
Sections~\ref{b5.1.-the-forest-statement-and-the-full-condition},
\ref{b5.6.-exact-general-homotopy-colimit-criterion}, and
\ref{b5.10.-an-explicit-all-index-local-filling-criterion}.

\paragraph{The geometric deployment data.}
A finite deployment over a poset $I$ consists of spaces $\Sigma_i$,
incoming faces $\Sigma_{ij}\subset\Sigma_j$ for $i<j$, and source maps
$s_{ij}:\Sigma_{ij}\to\Sigma_i$. A point of $\Sigma_{ij}$ is an arrow
from its source to its image in $\Sigma_j$. For $i<j<k$, the two
conditions that make these arrows compose are
\begin{equation}\label{eq:control}
 \Sigma_{ijk}=s_{jk}^{-1}(\Sigma_{ij})
       =\Sigma_{ik}\cap\Sigma_{jk},\qquad
 s_{ik}=s_{ij}s_{jk}\quad\hbox{on }\Sigma_{ijk}.
\end{equation}
The first equation is a condition on subsets; the second is an equation
of maps. Neither includes a radial control equation. Incoming faces
with incomparable labels are required to be disjoint. Higher flag spaces are the iterated arrow fiber products; the deployment category has only identity arrows within a label and composition given by the control equations.

\begin{theorem}[The exit comparison]\label{thm:exit-deployment}
Let $X$ be realized by a deployment over a finite poset with the data and equations specified above. Suppose all incoming target faces are closed and collared, their product collars commute at intersections, and the source maps are proper conically smooth stratified locally
trivial bundles, and the recursive collars commute with all source
maps and retained boundary collars. Require that cutting a minimal
stratum $A=\Sigma_i$ gives a collared complement $W$ and retained link $H$ with a specified identification $X=W\cup_H C(H\to A)$, where $C$ is the mapping cone and $H\to A$ is the induced source bundle. Require that $W$ and $H$ have deployments of the
same class on fewer labels. Here the complement carries the remaining-label
deployment, and the link at label $i$ has pieces $\Sigma_{ij}$ and flag
spaces $\Sigma_{i,j_0,\ldots,j_n}$, with source and attachment maps induced
from the given deployment; all these identifications retain the earlier
outer faces. Then the complete Segal category of the
deployment is naturally equivalent to $\operatorname{Exit}(X)$.
\end{theorem}

\paragraph{Tame pairs and regularity limits.}
A closed subset $Y$ of a metric space $X$ is \emph{forward tame}
if an open neighborhood $N$ has a homotopy $h:N\times[0,1]\to X$
with $h_0$ the inclusion, $h_1(N)\subset Y$, $h_t|_Y=\id$, and
$h_t(N\setminus Y)\subset X\setminus Y$ for $t<1$.
Neighborhood prospaces are indexed by shrinking these open neighborhoods.
\begin{corollary}[The global homotopy-link pushout]\label{cor:holink-pushout}
Let $Y$ be a closed forward-tame subset of a metric space $X$, and
$U=X\setminus Y$. For the ordinary homotopy link
$H=\{\gamma:\gamma(0)\in Y,\ \gamma(t)\in U\ (t>0)\}$,
with the compact-open topology and its initial and terminal endpoint maps, the prospaces of open neighborhoods $N\supset Y$ and of their punctures $N\setminus Y$ are equivalent to the constant prospaces $Y$ and $H$, respectively, compatibly with the endpoint maps up to their path homotopy. There is a natural equivalence
$X\simeq Y\cup^h_H U$.
\end{corollary}

\begin{theorem}[Proper Hurewicz fibrations with singular fibers]\label{thm:singular-fibers}
There exist a closed PL six-manifold $E_{\mathrm{PL}}$ and a proper
surjective PL map
\[
 p_{\mathrm{PL}}:E_{\mathrm{PL}}\longrightarrow S^1
\]
which is a Hurewicz fibration and has a fibre which is not a topological
manifold.  There also exist a closed smooth six-manifold $E_{\mathrm{sm}}$
and a smooth map $p_{\mathrm{sm}}:E_{\mathrm{sm}}\to S^1$ with the same
properties.  The smooth example has just one critical point.  The two
examples are constructed separately; the smoothing homeomorphism is not
asserted to be PL.
\end{theorem}

\paragraph{Relation to stratified homotopy theory.}
Lurie's exodromy theorem identifies constructible sheaves of spaces with exit-path functors under its geometric hypotheses \cite[Theorem~A.9.3]{Lurie}. Haine proves the stratified homotopy hypothesis after inverting equivalences on strata and links \cite[Theorem~3.2.4, pp.~25--26]{Haine}. Douteau and Waas compare the global topological and simplicial homotopy theories \cite{DW}. Ayala, Francis and Tanaka provide geometric conical neighbourhoods and resolutions \cite[Theorem~7.3.8]{AFT}. These results establish the homotopy-theoretic framework; the theorems above identify and test the sheaf-gluing diagrams proposed in the manuscript. A sheaf of sets remembers less coherence than an arbitrary sheaf of spaces, so truncating an ordinary-topos descent diagram is not a replacement for a higher descent theorem.

Hua's recent exposition already gives the successor-pair obstruction
for the linear three-stratum plane and proves topological gluing in
depth one \cite[Counterexample~70.3 and the following theorem,
p.~340]{Hua}. Theorem~\ref{thm:stones} proves a sheaf obstruction
even after all the prescribed pair and triple pieces are included,
and gives an example with connected strata. Hua also discusses the
arbitrary-index hypercover formulation on p.~339. Theorem~\ref{thm:higher-comparisons}
separates the finite-comparison condition, the shape test, and
hypercomplete local filling, with the relevant established inputs
identified in their proofs.

\paragraph{Critical points and the topology of fibres.}
The isolated critical point in Theorem~\ref{thm:singular-fibers}
has a classical context. Farber and Sch\"utz prove that every nonzero
real degree-one cohomology class of a closed connected smooth manifold
has a closed one-form representative with at most one zero
\cite[Theorem~1, author-version pp.~1--2]{FarberSchutz}.
For an integral class, integration gives a circle-valued map with at
most one critical point. Funar, Pintea and Zhang obtain another such
construction by composing a smooth homeomorphism with a bundle
\cite[Proposition~4.1, p.~363]{FunarPinteaZhang}; its fibres retain
their manifold topology. Theorem~\ref{thm:singular-fibers} establishes
the full Hurewicz lifting property together with a nonmanifold fibre.
Cacciatori and Pigola instead give obstructions for Hurewicz fibrations
with aspherical total space and locally contractible fibre
\cite[Theorems~1.1--1.2, author-version pp.~1--2]{CacciatoriPigola}.
Their hypotheses do not apply to the total spaces constructed here.

\paragraph{Why the proof works.}
The intersection axiom in \eqref{eq:closed-family} gives each point a least stratum index. That index places the point in one of the open tubes. Descent then compares sections on every genuine overlap. The counterexample assigns two different values to the same generization map in two tubes; the named incidence pieces see its source and target but never compare that map. Finally, a general open diagram can be analysed stalk by stalk: connectedness identifies copies of a stalk and simple connectivity removes monodromy among the identifications.

\section{Neighbourhoods and the open cover}
We first prove that the neighbourhood systems retain information which an ordinary inverse limit would discard.

\enlargethispage{\baselineskip}
\begin{proposition}\label{prop:collapse}
Let $Y$ be closed in a regular Hausdorff space $Z$. The inverse limit of the open subtopoi $\Sh(N)$, over all open $N\supset Y$, is $\Sh(Y)$. The inverse limit of $\Sh(N\setminus Y)$ is the empty topos.
\end{proposition}
\begin{proof}
A geometric morphism $f:E\to\Sh(Z)$ factors through $\Sh(N)$ precisely when the inverse image of the subterminal $N$ is terminal. Regularity gives
\[
 Z\setminus Y=\bigcup_{N\supset Y}\operatorname{int}(Z\setminus N).
\]
Indeed a point outside $Y$ has an open neighbourhood whose closure misses $Y$; its complement is an admissible $N$. A morphism factoring through every $N$ therefore sends each displayed complementary open to the initial object. Inverse images preserve unions, so it sends $Z\setminus Y$ to the initial object. This is exactly the condition for factorization through the complementary closed subtopos $\Sh(Y)$. Such factorizations are fully faithful on mapping categories, proving the bicategorical universal property. Applying the same argument inside $Z\setminus Y$ proves the punctured assertion.
\end{proof}
For example, shrinking intervals about the origin retains a two-sided punctured neighbourhood as a pro-object. The ordinary limit of the punctured open subtopoi is empty. This is the reason for keeping the system of neighbourhoods.

\begin{lemma}\label{lem:cover}
Under \eqref{eq:closed-family}, each point $x$ lies in exactly one $S_i$, and $X_i\setminus R_i=S_i$. Every $R_i$ is closed and every $(W_i^\alpha)_i$ covers $X$.
\end{lemma}
\begin{proof}
The set $I_x=\{i:x\in X_i\}$ is finite, nonempty, and downward directed: for $i,j\in I_x$, the intersection formula supplies $k\in I_x$ with $k\le i,j$. Successively taking lower bounds in this finite set produces its least element $i(x)$. A point belongs to $S_i$ exactly when $i=i(x)$. If $x\in X_i$ lies in some $X_k$ with $i\nleq k$, an intersection lower bound $h\le i,k$ must satisfy $h<i$, so $x\notin S_i$. Conversely a lower index $h<i$ satisfies $i\nleq h$ and contributes to $R_i$. Thus $X_i\setminus R_i=S_i$. A locally finite union of closed subsets is closed, proving the assertion about $R_i$. Finally $x\in X_{i(x)}\subset N_{i(x)}$ and $x\notin R_{i(x)}$, so $x\in W_{i(x)}^\alpha$.
\end{proof}

Local finiteness is used at two separate points: it makes $I_x$ finite and makes $R_i$ closed. The family of all singleton subsets of the usual real line satisfies the covering and intersection identities but is not locally finite. Its purported tubes are singletons and are not open. Thus the hypothesis cannot be removed from this construction.

\section{Descent, including the synchronized pro-object}
We now prove Theorem~\ref{thm:tubes} directly, including its universal property on geometric morphisms.

For an open cover $(W_i)$ a descent datum consists of sheaves $F_i$ on $W_i$ and isomorphisms $g_{ij}:F_i|_{W_i\cap W_j}\to F_j|_{W_i\cap W_j}$ satisfying $g_{ii}=1$ and $g_{jk}g_{ij}=g_{ik}$. Define
\[
 F(O)=\left\{(s_i)\in\prod_iF_i(O\cap W_i):
 g_{ij}(s_i|_{O\cap W_i\cap W_j})=s_j|_{O\cap W_i\cap W_j}\right\}.
\]
The sheaf axiom in each $W_i$ proves the sheaf axiom for $F$: compatible sections glue componentwise, and compatibility on overlaps is preserved because it can be tested locally. The cocycle identifies $F|_{W_i}$ with $F_i$. Compatible maps of descent data induce and are uniquely recovered from maps of $F$, giving an equivalence of categories.

Given compatible geometric morphisms $f_i:\Sh(W_i)\to T$, apply this construction to $f_i^*A$ for each object $A$ of $T$ and to all its arrows. It produces $f^*:T\to\Sh(X)$. Restrictions to the cover preserve finite limits and all colimits and are jointly conservative, so $f^*$ has the same preservation properties. A right adjoint can also be written explicitly: for $F\in\Sh(X)$ take the equalizer in $T$ of
\[
 \prod_i f_{i*}(F|_{W_i})
 \rightrightarrows
 \prod_{i,j} f_{ij*}(F|_{W_i\cap W_j}),
\]
where $f_{ij}$ is either of the identified restrictions. The section-gluing formula proves the adjunction. Natural transformations between inverse-image functors glue by the same argument. Hence $\Sh(X)$ is the bicategorical colimit, not merely a reconstruction of its objects. Objects, overlap identifications and cocycles occupy degrees $0,1,2$ respectively.

For the pro-object statement, order tuples of neighbourhoods by simultaneous shrinking. This is cofiltered: componentwise intersections refine any two choices. The indexing category $\Delta_{\le2}$ is finite. A filtered bicolimit of categories commutes with a finite bilimit: objects, comparison isomorphisms and equations in a finite cone can all be represented at one stage; two such finite systems agree after a common further refinement precisely when their images agree. This proves essential surjectivity and full faithfulness of the comparison. Therefore \eqref{eq:pro-maps} gives
\begin{align*}
 \mathop{\mathrm{bilim}}_{q\le2}\Hom_{2\Pro}(C_q,(T_b)_b)
 &\simeq \mathop{\mathrm{bilim}}_b\mathop{\mathrm{bicolim}}_\alpha
       \mathop{\mathrm{bilim}}_{q\le2}\Hom(C_q^\alpha,T_b)\\
 &\simeq\mathop{\mathrm{bilim}}_b\Hom(\Sh(X),T_b).
\end{align*}
The augmentations commute with shrinking because they are induced by the identity of $X$. This is the desired universal property. For finite $I$, the subset diagram with at most three indices encodes exactly the same descent data. For infinite $I$ the proof uses finitely many \emph{degrees} after forming the coproducts; it does not commute a filtered colimit with an infinite product of independently indexed neighbourhood systems. This proves Theorem~\ref{thm:tubes}.

\section{Two counterexamples to incidence-piece gluing}
We now exhibit the missing information in the smallest linear chain. Let $o$ be the origin, let $p,n$ be the positive and negative real rays, and let $u,d$ be the upper and lower half-planes. The cell poset has relations
\[
 o<p<u,\quad o<n<u,\quad o<p<d,\quad o<n<d.
\]
A covariant set-valued functor on this poset gives a sheaf with the indicated stalks and generization maps: use the associated sheaf on its upward Alexandrov topology and pull back by the continuous cell map. This gives concrete topological sheaves for the following finite diagrams.

Write $\operatorname{germ}_A Z$ for the pro-system of open neighbourhoods of $A$ in $Z$. The nine distinct pieces for the chain are
\[
\begin{array}{c|l}
 S_0&\{0\}\\
 S_1&\R\setminus\{0\}\\
 S_2&\C\setminus\R\\
 V_{01}&\operatorname{germ}_0\R\\
 V_{01}^*&\operatorname{germ}_0\R\setminus\{0\}\\
 V_{12}&\operatorname{germ}_{\R}\C\setminus\{0\}\\
 V_{12}^*&\operatorname{germ}_{\R}\C\setminus\R\\
 V_{02}&\operatorname{germ}_0\C\\
 V_{02}^*&\operatorname{germ}_0\C\setminus\R.
\end{array}
\]
For every pair there are the usual inclusions from the smaller stratum and from the punctured tube into the tube, and from the punctured tube into the larger stratum. The extra generating inclusions are $V_{01}\to V_{02}$ and $V_{02}^*\to V_{12}^*$. Include all composites. The triple-piece formula of the manuscript, for its only strict triple $0<1<2$, repeats $V_{02}^*,V_{02}^*,V_{02}$, so it adds no comparison. This explicitly describes the diagram used here.

Give $o,p,n,d$ a singleton stalk and $u$ a two-element stalk $\{0,1\}$. On $V_{02}$ let both ray-to-$u$ maps and the origin-to-$u$ map select $0$. On $V_{12}$ let both ray-to-$u$ maps select $1$. Maps with singleton targets are unique. All restrictions along the displayed inclusions agree: any common piece either contains only real cells or only nonreal cells, and therefore sees neither of the discrepant ray-to-half-plane maps. Choose the overlap identifications to be identities.

A hypothetical global sheaf would induce equal generization maps on the punctured disk lying in both tubes. The prescribed identifications on the rays and on the upper half-plane force those maps to select both $0$ and $1$. Shrinking either tube does not remove the rays near zero, so the contradiction persists at every stage. This disproves the all-pair and all-triple assertion.

For the successor diagram alone, put the value $0$ on $p\to u$ and $1$ on $n\to u$ in $V_{12}$. Give the real tube its singleton sheaf. The successor pieces again agree. A section representing the singleton germ at the origin would be constant on a sufficiently small connected upper half-disk, but its restrictions near the two rays force different values. Thus it cannot exist.

These failures of sheaf descent test the required universal property of topoi. The topos $\mathbf{Set}^{\mathrm{FinSet}}$ classifies an object: its geometric morphisms out of $E$ are equivalent to objects of $E$. Indeed a left-exact functor $\mathrm{FinSet}^{\mathrm{op}}\to E$ is the functor $K\mapsto A^K$ for its value $A$ at a singleton. Hence the compatible local objects above give a cocone of geometric morphisms with no extension. The same argument applies to germs through \eqref{eq:pro-maps}.

For connected strata replace the chain by $\{0\}\subset\R^2\subset\R^4$. On the closed plane put a singleton; on its complement put the constant two-element sheaf. Closed--open recollement realizes either constant generization choice $0$ or $1$. Put the first choice on the origin tube and the second on the punctured plane tube. Their named restrictions coincide, while their maps from the middle stratum to the top stratum differ on every punctured-origin intersection. The complement of a linear plane in $\R^4$ is connected. The three strata therefore satisfy the connectedness assertion, completing the proof of Theorem~\ref{thm:stones}.

\section{The pointwise criterion and minimal arities}
This section proves the criterion for an arbitrary open diagram and explains why triple compatibility is indispensable.

A descent datum on $D:C\to\operatorname{Open}(X)$ is a family of sheaves $F_c$ together with compatible isomorphisms along every arrow. At a point $x$ these stalks form a functor from the groupoidification of $C_x$ to sets. If that groupoid is terminal, all stalks identify uniquely. Take the disjoint union of the \emph{\'etale spaces} of the $F_c$ and quotient by the identifications along arrows. The quotient map is open: the saturation of an open set is the union of images under finite strings of open-subspace homeomorphisms. Connectedness gives a bijection from the stalk of any $F_c$ at $x$ to the quotient fibre, and simple connectivity makes that identification independent of the chosen string. Consequently the \emph{\'etale space} of $F_c$ maps homeomorphically onto the quotient above $D(c)$. Since these opens cover $X$, the quotient is an \emph{\'etale space} over $X$ and gives the descended sheaf. Maps descend too. Applying this construction to inverse-image functors proves the topos universal property as in the previous section.

Conversely, fix $x$. If $C_x$ is empty, two distinct skyscraper sheaves at the closed point $x$ have identical restrictions. If $C_x$ has two components, assign different cardinalities to stalks on these components. If it is connected with nontrivial isotropy, take a regular action of that isotropy group and transport it to its groupoid. Extend these set-valued data to skyscraper sheaves at $x$ on opens containing $x$ and terminal sheaves on the other opens. Since $x$ is closed, restricting a skyscraper to an open missing $x$ is terminal. Thus these are descent data on the whole diagram. A global sheaf would give a constant stalk datum with trivial monodromy, a contradiction. The object classifier again converts this to a topos obstruction. The nerve has the specified groupoidification as its fundamental groupoid because edges give generators and triangles impose composition relations. Higher simplices impose no additional fundamental-groupoid relations. This proves Theorem~\ref{thm:membership}.

\begin{corollary}\label{cor:minimal}
Vertices, pair overlaps and triple cocycles are the minimal universal arities for open-cover descent of ordinary topoi. In a given cover some pieces can be redundant, but intersections of four or more members are never needed for this descent problem.
\end{corollary}
\begin{proof}
Sufficiency was proved above. A point lying only in a deleted member witnesses the need for a vertex. Two members covering a point without their pair intersection leave two unrelated stalks. Finally take three copies of a point. Keep their three pair intersections but delete the triple intersection. The incidence category has nerve a cycle, and a nontrivial automorphism of a two-element set supplies monodromy around it. The criterion then disproves descent. Adding the triple fills the cycle. Thus no smaller universal list works.

For a binary cover $X=N\cup U$, these are exactly the two vertices
and their overlap in $\Sh(N)\leftarrow\Sh(N\cap U)\to\Sh(U)$.
Deleting either vertex or the overlap gives the corresponding
one-point obstruction above. An additional object $\Sh(Y)\to\Sh(N)$
with no other incidence is redundant: a pseudocone on the old diagram
already specifies its value and comparison on this new leaf by
composition, and all extensions form a contractible category. Thus
adjoining this soul does not change the bicolimit.
\end{proof}

\section{Base change and the order-defined deleted set}
We determine exactly when an inverse image of a neighbourhood system is cofinal in the neighbourhoods of its inverse image.

\begin{proposition}\label{prop:basechange}
Let $f:Y\to X$ be continuous and let $Z\subset X$. The sets $f^{-1}(N)$, for open $N\supset Z$, are cofinal among the open neighbourhoods of $f^{-1}Z$ if and only if
\[
 \overline{f(F)}\cap Z=\varnothing
 \quad\text{for every closed }F\subset Y\text{ with }F\cap f^{-1}Z=\varnothing.
\]
In particular every closed map satisfies this criterion. Arbitrary open immersions do not.
\end{proposition}
\begin{proof}
For an open $W\supset f^{-1}Z$ put $F=Y\setminus W$. The required containment $f^{-1}N\subset W$ is equivalent to $N\cap f(F)=\varnothing$. Such an open $N$ containing $Z$ exists exactly when $Z$ misses $\overline{f(F)}$. This proves both directions. If $f$ is closed, $f(F)$ is closed and disjoint from $Z$, so the criterion holds. For the open immersion $(0,1]\hookrightarrow[0,1]$ and $Z=\{0\}$, the inverse image is empty but every neighbourhood of $0$ has nonempty inverse image. Taking $W=\varnothing$ disproves cofinality.
\end{proof}

For the closed family of Section~1 the deleted sets commute with every inverse image:
\[
 f^{-1}R_i=\bigcup_{k:\ i\nleq k}f^{-1}X_k.
\]
This equality uses the indexing order, not a reconstruction of that order from the supports after pullback. Coincident or empty inverse-image strata may erase support information without changing the index relation. This is the precise reason for the order-defined convention.


For example, let $X_0=\{0\}\subset X_1=\R$ and pull back along
$\{0\}\hookrightarrow\R$. The support-defined deleted set for index
$1$ was $\{0\}$, but recomputing strict containment after pullback makes
it empty because the two supports coincide. The order-defined deleted
set remains $\{0\}$. More generally
$R_i\cap X_i=\bigcup_{h<i}X_h$: intersect each term of $R_i$ with
$X_i$ and use \eqref{eq:closed-family}; its lower bound cannot equal
$i$. Conversely every $h<i$ occurs in the union.

Descent for a specified open-cover diagram is local on $X$.
The identity $(W_i\cap W_j)\cap O=(W_i\cap O)\cap(W_j\cap O)$,
and its triple version, show that the explicit gluing construction
commutes with restriction to $O$. The local inverses agree on overlaps
by uniqueness of sheaf gluing. This holds stagewise for a specified
pulled-back synchronized system. It does not justify replacing that
system by independently recomputed neighborhoods, as the open-immersion
counterexample shows.

\section{Closed reconstruction and the stable-graph meet condition}\label{sec:closed-graphs}
Open descent and closed reconstruction have different mechanisms. The
closed-family axioms recover a space as a quotient, but a diagram of
embeddings need not embed its own members in that quotient. This section
identifies the missing condition and determines when it holds at the
level of stable-graph types.

\begin{proposition}[Closed reconstruction and its embedding obstruction]\label{prop:closed-reconstruction}
For a locally finite closed family satisfying \eqref{eq:closed-family},
the colimit of the inclusion diagram $i\mapsto X_i$ is $X$ as a
topological space. Empty strata may be discarded.
For an arbitrary finite diagram of closed embeddings over a finite
poset, its members embed as closed subspaces of its colimit if and
only if no zigzag of transition identifications identifies distinct
points in one member. If the indexing poset is a finite
meet-semilattice, the cartesian-intersection condition
\[
 X_{i\wedge j}\xrightarrow{\sim}X_i\times_{X_k}X_j
 \qquad(i,j\leq k)
\]
is sufficient and makes the intersection of the images of $X_i,X_j$
exactly the image of $X_{i\wedge j}$.
\end{proposition}
\begin{proof}
The least containing index of a point, constructed in
Lemma~\ref{lem:cover}, identifies all representatives of that point
through a common lower member. Hence the set colimit is $X$. The
surjection $\coprod_iX_i\to X$ is closed: the union of a family of
subsets closed in the closed $X_i$ is closed by local finiteness.
It is therefore a quotient map. Moreover
$X_k=\bigcup_{i\leq k,\,S_i\ne\varnothing}X_i$, since the least
containing index of any point of $X_k$ occurs in that union. This
identity preserves the intersection axiom after deleting empty strata.

For the second assertion the zigzag condition is precisely injectivity
of each set map to the colimit. Under it a zigzag revisiting an index
revisits the same point, so its intervening loop can be removed.
There are finitely many simple index paths. Transporting a closed
subset along one step takes either a closed image under a closed
embedding or an inverse image. Its saturation therefore meets each
member in a finite union of closed subsets. The quotient topology
makes the map of every member closed, and the same reasoning applies
to each closed subset of that member. Injectivity now gives a closed
embedding. Necessity follows from injectivity.

Under the cartesian condition, follow a zigzag while maintaining a
point of a common lower member mapping to the starting and current
representatives. An upward step preserves it. A downward step to
$X_j$ lifts it uniquely through the pullback with the previous lower
member, replacing that lower index by its meet with $j$. At the end,
representatives in $X_i,X_j$ come from a common lower member and hence
from $X_{i\wedge j}$. Taking $i=j$ proves the zigzag condition; taking
arbitrary $i,j$ proves the intersection statement, including its
subspace topology because the maps are closed embeddings.
\end{proof}

\begin{counterexample}[A meet in the index poset does not suffice]\label{ex:closed-meet}
Index a diagram by the seven proper subsets of $\{1,2,3\}$, including
the empty set. Its meet is intersection. Put a singleton $\{*\}$ at
the bottom and the discrete set $\{*,0,1\}$ at each other index.
Bottom maps preserve $*$; all other transition maps are identities
except $\{1\}\to\{1,2\}$, which exchanges $0$ and $1$.
Every composable pair of strict arrows begins at the bottom, so this
is a functor of closed embeddings. On the disjoint union of its members,
order two points by the transition maps and the index order. This is
reflexive and transitive by functoriality, and antisymmetric because
indices cannot strictly increase in a cycle and each same-index map
is the identity. The space is finite discrete, so the order relation
is closed; the projection to the index poset is increasing and each
fiber has the equality order. Thus the associated ordered-space
qualification does not repair this example. The six nonbottom indices form a
hexagon, and transport around it exchanges $0,1$. Thus its colimit
has only two points, the $*$ orbit and the $\{0,1\}$ orbit. None of
the three-point members embeds. The cartesian condition fails, even
though the poset is a meet-semilattice.
\end{counterexample}

\begin{proposition}[The precise intersection condition]\label{prop:meet-intersection}
Suppose all strata $S_i$ in \eqref{eq:closed-family} are nonempty.
Then $i\mapsto X_i$ is an order embedding. Binary meets for pairs with
a common lower bound exist exactly when every nonempty intersection
$X_i\cap X_j$ is one member $X_k$. A literal all-pairs meet condition
also requires an empty bottom whenever disjoint closed members occur.
\end{proposition}
\begin{proof}
If $X_i\subset X_j$, choose $x\in S_i$. Its least containing index
is $i$, hence $i\leq j$. The converse is an assumption. An index
representing the intersection is therefore the greatest lower bound.
Conversely, if the meet exists, the intersection axiom identifies the
union of all lower members with that meet's member. Nonempty strata
make a nonempty intersection equivalent to existence of a common
lower index. A formal empty bottom handles disjoint intersections. Every finite meet-semilattice is realized by its upper Alexandrov space with closed members $X_i=\mathord\downarrow i$: intersections are the downsets of meets, and the stratum at $i$ is $\{i\}$. Thus this is a closure-system condition with no implicit smoothness requirement.
\end{proof}

The distinction persists under restriction to unions of strata.
Let $I''\subset I'\subset I$ be lower ideals, and set
$X'=\bigcup_{i\in I'}X_i$, $Z=\bigcup_{i\in I''}X_i$ and
$U=X'\setminus Z$. The remaining strata are precisely those indexed
by $I'\setminus I''$: membership in a removed closed member forces
the least containing index to belong to $I''$. Closed restriction to
$X'$ retains every lower member relevant to a surviving index.
Open restriction need not preserve its tube's homotopy type. For
$\{0\}\subset\R\subset\C$, neighborhoods of $\R$, with $0$
deleted, have a cofinal family of variable-width strips with a puncture
and hence constant pro-homotopy type $S^1$. In $\C^*$, neighborhoods
of $\R^*$ have a cofinal family of two disjoint strips of width less
than $|x|/2$, and hence the type of two points. Deleting the real center
gives respectively two and four contractible components. Positive
continuous minorants of the distance to the complement give cofinality.
Thus these discrepancies survive every admissible shrinking.
Restriction of constructible sheaves to $U$ is nevertheless essentially
surjective: extension by the empty sheaf is a fully faithful left
adjoint. It is an equivalence exactly when $Z$ is empty. If $Z$ is
nonempty, the closed direct images of the constant one- and two-point
sheaves on $Z$ are constructible and nonisomorphic, but both restrict
to the terminal sheaf on $U$.

\subsection{Stable graphs}
The graph conventions and complete classification are in Theorem~\ref{thm:graph-meets}.

\begin{proof}[Proof of Theorem~\ref{thm:graph-meets}]
In genus zero every graph is a tree. Its edges give distinct stable
bipartitions of the labelled legs. A common refinement contains both
sets of bipartitions; contract all edges outside their union. The
result is independent of the common refinement. To see uniqueness,
root at one labelled leg and orient each split by its side avoiding
that leg. These sides form a nested-or-disjoint family, and its
inclusion hierarchy reconstructs the rooted tree. Every common
refinement contracts to that tree.

For the other positive cases the finite certificate below lists every
graph and every single-edge contraction. Here is a fully bounded
way to verify exhaustiveness. Summing stability gives
$|V|\leq2g-2+n$ and $|E|\leq3g-3+n$. For each such vertex count,
list all tuples of vertex genera with total at most $g$. The required
edge count is $g-\sum_vg_v+|V|-1$. Distribute it among the
$|V|(|V|+1)/2$ unordered vertex pairs, including loops, and distribute
the labelled legs among the vertices. Retain the connected stable
matrices and identify vertex permutations. These ranges are finite
and include every graph by its defining data. Contracting the edges
in the displayed representatives gives exactly the displayed
contraction columns. Taking transitive closures and intersecting the
sets of refinements gives the meet entries printed below. Every
nonempty such intersection has the indicated greatest member.
This verifies every pair for the four positive cases.

There are uniform counterexamples in all other stable types. For
$g\geq3$, let $G$ be a bridge with vertex genera $1,g-1$, and $H$
a single vertex of genus $g-1$ with a loop. Put all legs, if present,
on the second vertex of $G$. One common refinement has a loop on a
genus-zero vertex joined to a vertex of genus $g-1$; another has a
bridge with genera $1,g-2$ and a loop on its second vertex. Put all
legs on the second vertex in both. Loop contraction gives $G$ and
bridge contraction gives $H$. Both refinements are stable and their
only one-edge contractions are $G,H$. They are nonisomorphic, since
only one has a genus-zero vertex, so the meet does not exist.

For $g=2,n\geq1$, let $G$ have two genus-one vertices joined by a
bridge, all legs on one vertex; let $H$ have one genus-one vertex,
all legs and a loop. Refine to a genus-zero loop vertex joined to a
genus-one vertex. Put all legs on the loop vertex in one refinement
and on the other vertex in the other. These are two distinct minimal
common refinements by the same two-edge argument.

For $g=1,n\geq4$, let $G$ and $H$ be two-edge cycles with leg
partitions $\{1\}\mid\{2,\ldots,n\}$ and
$\{2\}\mid\{1,3,\ldots,n\}$ respectively. One common refinement is
a triangle with leg sets $\{1\},\{2\},\{3,\ldots,n\}$.
Another is a square whose opposite vertices carry $\{1\},\{2\}$
and whose remaining vertices carry $\{3\},\{4,\ldots,n\}$.
All vertices are rational and stable. Both contract to $G,H$.
In a common refinement obtained from the square, vertices carrying
legs $1$ and $2$ must each remain separate from every other marked
vertex. Every square edge meets one of these two vertices, so no
edge can be contracted while retaining both requirements. The square
is therefore minimal. The triangle is minimal because a two-edge
common refinement would equal both distinct graphs $G,H$.
These cases exhaust the stable types outside the list.
\end{proof}

Stable-curve deformation parameters smooth the individual nodes and
therefore contract the corresponding graph edges. Every stable graph
occurs by gluing smooth pointed components. Thus the combinatorial
closure order and its common-refinement intersections are the actual
ones for the stable-graph stratification of the moduli space; see
\cite[Proposition~4.14; Sections~6.1 and~7.3]{Schmitt}.
Proposition~\ref{prop:meet-intersection} gives the precise connection
with the single-member intersection condition.

\subsection{Complete finite meet certificates}
The following tables are part of the finite proof. In a graph row,
$g$ is its tuple of vertex genera; $\ell$ lists the vertex carrying
leg $1$, then leg $2$, and so on; $E$ lists edges $ij$, with $ii$
a loop. Vertices start at $1$. The final column gives all distinct
single-edge contractions. In the meet table, an entry is the greatest
common refinement; a dash means no common refinement. The diagonal
is implicit. Graph numbers are local to their $(g,n)$ table.

\paragraph{Type $(1,1)$.} There are 2 graphs.
\begin{center}\small
\begin{tabular}{rllll}\toprule
No.&$g$&$\ell$&$E$&Contractions\\\midrule
1 & (0) & 1 & 11 & 2\\
2 & (1) & 1 & -- & --\\
\bottomrule\end{tabular}\end{center}
\begin{center}\scriptsize\setlength{\tabcolsep}{2pt}
\begin{tabular}{r|rr}
$\wedge$&1&2\\\hline
1&1&\\
2&1&2\\
\end{tabular}\end{center}

\paragraph{Type $(1,2)$.} There are 5 graphs.
\begin{center}\small
\begin{tabular}{rllll}\toprule
No.&$g$&$\ell$&$E$&Contractions\\\midrule
1 & (0) & 11 & 11 & 5\\
2 & (0,0) & 11 & 12,22 & 1,4\\
3 & (0,0) & 12 & 12,12 & 1\\
4 & (0,1) & 11 & 12 & 5\\
5 & (1) & 11 & -- & --\\
\bottomrule\end{tabular}\end{center}
\begin{center}\scriptsize\setlength{\tabcolsep}{2pt}
\begin{tabular}{r|rrrrr}
$\wedge$&1&2&3&4&5\\\hline
1&1&&&&\\
2&2&2&&&\\
3&3&--&3&&\\
4&2&2&--&4&\\
5&1&2&3&4&5\\
\end{tabular}\end{center}

\begin{samepage}
\paragraph{Type $(1,3)$.} There are 23 graphs.
\begin{center}\small
\begin{tabular}{rllll}\toprule
No.&$g$&$\ell$&$E$&Contractions\\\midrule
1 & (0) & 111 & 11 & 23\\
2 & (0,0) & 111 & 12,22 & 1,19\\
3 & (0,0) & 112 & 12,22 & 1,20\\
4 & (0,0) & 121 & 12,22 & 1,21\\
5 & (0,0) & 211 & 12,22 & 1,22\\
6 & (0,0) & 112 & 12,12 & 1\\
7 & (0,0) & 121 & 12,12 & 1\\
8 & (0,0) & 122 & 12,12 & 1\\
9 & (0,0,0) & 112 & 13,23,23 & 3,6\\
10 & (0,0,0) & 121 & 13,23,23 & 4,7\\
11 & (0,0,0) & 211 & 13,23,23 & 5,8\\
12 & (0,0,0) & 113 & 13,22,23 & 2,3,16\\
13 & (0,0,0) & 131 & 13,22,23 & 2,4,17\\
14 & (0,0,0) & 311 & 13,22,23 & 2,5,18\\
15 & (0,0,0) & 123 & 12,13,23 & 6,7,8\\
16 & (0,0,1) & 112 & 12,23 & 19,20\\
17 & (0,0,1) & 121 & 12,23 & 19,21\\
18 & (0,0,1) & 211 & 12,23 & 19,22\\
19 & (0,1) & 111 & 12 & 23\\
20 & (0,1) & 112 & 12 & 23\\
21 & (0,1) & 121 & 12 & 23\\
22 & (0,1) & 211 & 12 & 23\\
23 & (1) & 111 & -- & --\\
\bottomrule\end{tabular}\end{center}
\end{samepage}
\begin{center}\scriptsize\setlength{\tabcolsep}{2pt}
\begin{tabular}{r|rrrrrrrrrrrrrrrrrrrrrrr}
$\wedge$&1&2&3&4&5&6&7&8&9&10&11&12&13&14&15&16&17&18&19&20&21&22&23\\\hline
1&1&&&&&&&&&&&&&&&&&&&&&&\\
2&2&2&&&&&&&&&&&&&&&&&&&&&\\
3&3&12&3&&&&&&&&&&&&&&&&&&&&\\
4&4&13&--&4&&&&&&&&&&&&&&&&&&&\\
5&5&14&--&--&5&&&&&&&&&&&&&&&&&&\\
6&6&--&9&--&--&6&&&&&&&&&&&&&&&&&\\
7&7&--&--&10&--&15&7&&&&&&&&&&&&&&&&\\
8&8&--&--&--&11&15&15&8&&&&&&&&&&&&&&&\\
9&9&--&9&--&--&9&--&--&9&&&&&&&&&&&&&&\\
10&10&--&--&10&--&--&10&--&--&10&&&&&&&&&&&&&\\
11&11&--&--&--&11&--&--&11&--&--&11&&&&&&&&&&&&\\
12&12&12&12&--&--&--&--&--&--&--&--&12&&&&&&&&&&&\\
13&13&13&--&13&--&--&--&--&--&--&--&--&13&&&&&&&&&&\\
14&14&14&--&--&14&--&--&--&--&--&--&--&--&14&&&&&&&&&\\
15&15&--&--&--&--&15&15&15&--&--&--&--&--&--&15&&&&&&&&\\
16&12&12&12&--&--&--&--&--&--&--&--&12&--&--&--&16&&&&&&&\\
17&13&13&--&13&--&--&--&--&--&--&--&--&13&--&--&--&17&&&&&&\\
18&14&14&--&--&14&--&--&--&--&--&--&--&--&14&--&--&--&18&&&&&\\
19&2&2&12&13&14&--&--&--&--&--&--&12&13&14&--&16&17&18&19&&&&\\
20&3&12&3&--&--&9&--&--&9&--&--&12&--&--&--&16&--&--&16&20&&&\\
21&4&13&--&4&--&--&10&--&--&10&--&--&13&--&--&--&17&--&17&--&21&&\\
22&5&14&--&--&5&--&--&11&--&--&11&--&--&14&--&--&--&18&18&--&--&22&\\
23&1&2&3&4&5&6&7&8&9&10&11&12&13&14&15&16&17&18&19&20&21&22&23\\
\end{tabular}\end{center}

\begin{samepage}
\paragraph{Type $(2,0)$.} There are 7 graphs.
\begin{center}\small
\begin{tabular}{rllll}\toprule
No.&$g$&$\ell$&$E$&Contractions\\\midrule
1 & (0) & -- & 11,11 & 5\\
2 & (0,0) & -- & 12,12,12 & 1\\
3 & (0,0) & -- & 11,12,22 & 1,4\\
4 & (0,1) & -- & 11,12 & 5,6\\
5 & (1) & -- & 11 & 7\\
6 & (1,1) & -- & 12 & 7\\
7 & (2) & -- & -- & --\\
\bottomrule\end{tabular}\end{center}
\end{samepage}
\begin{center}\scriptsize\setlength{\tabcolsep}{2pt}
\begin{tabular}{r|rrrrrrr}
$\wedge$&1&2&3&4&5&6&7\\\hline
1&1&&&&&&\\
2&2&2&&&&&\\
3&3&--&3&&&&\\
4&3&--&3&4&&&\\
5&1&2&3&4&5&&\\
6&3&--&3&4&4&6&\\
7&1&2&3&4&5&6&7\\
\end{tabular}\end{center}

\section{Closed--open gluing and uniform tube data}\label{sec:recollement}
The overlap obstruction has a useful positive formulation. Write
$\operatorname{Loc}(Z)$ for all locally constant sheaves of sets on
$Z$; stalks need not be finite.

\begin{proposition}[The exact gluing datum]\label{prop:recollement}
For a closed inclusion $i:Y\hookrightarrow X$ and its open complement
$j:U\hookrightarrow X$, sheaves on $X$ are equivalent to triples
\[
 (L,E,\alpha),\qquad L\in\Sh(U),\quad E\in\Sh(Y),\quad
 \alpha:E\longrightarrow i^{-1}j_*L.
\]
The subcategory constructible for this two-part decomposition is
obtained by requiring $L,E$ to be locally constant.
\end{proposition}
\begin{proof}
Put $T=i^{-1}j_*L$ and form
$F=j_*L\times_{i_*T}i_*E$, using the adjunction unit and $i_*\alpha$.
At a point of $U$, the last two sheaves have terminal stalks, so the
pullback stalk is $L_x$. At a point of $Y$, it is
$T_y\times_{T_y}E_y=E_y$. Its induced gluing map is $\alpha$.
Conversely, adjunction maps a sheaf to this pullback formed from its
two restrictions; this map is an isomorphism on every stalk, hence
an isomorphism. The pullback universal property identifies morphisms
with commuting pairs. This proves the equivalence, including empty
parts and arbitrary stalk cardinalities. Compare the standard
recollement formulation in \cite[Example~A.8.4]{Lurie}.
\end{proof}

Suppose a neighborhood $N$ retracts to $Y$ by $r$, and put
$p=r|_{N\setminus Y}$. Taking germs gives
\[
 \theta_L:p_*(L|_{N\setminus Y})\longrightarrow i^{-1}j_*L.
\]
The description by $L\in\operatorname{Loc}(U)$,
$E\in\operatorname{Loc}(Y)$ and an arrow
$p^{-1}E\to L|_{N\setminus Y}$ is valid over each fixed pair exactly
when $\operatorname{Hom}(E,\theta_L)$ is bijective. This follows
immediately from adjunction and Proposition~\ref{prop:recollement}.
Requiring it on every open subset of $Y$ makes $\theta_L$ itself an
isomorphism: test the constant one-point sheaf on each such open.
A conical bundle with compact locally contractible links over a
paracompact base with contractible local charts satisfies this condition.
On a chart the punctured tube is $W\times K\times(0,\epsilon)$;
local-system sections are unchanged by shrinking the contractible
factors. Compactness of $K$ makes these radial neighborhoods cofinal
at each point of the base. For the usual monodromy formulation,
local path connectedness and semilocal simple connectivity are also
required.

\begin{proposition}[The exact pro-tube criterion]\label{prop:pro-recollement}
Let $\mathcal T=2\colim_{N\supset Y}\operatorname{Loc}(N)$ and
$\mathcal T^*=2\colim_{N\supset Y}\operatorname{Loc}(N\setminus Y)$,
where arrows shrink neighborhoods. Let
$r:\mathcal T\to\operatorname{Loc}(Y)$ and
$a:\mathcal T\to\mathcal T^*$ be restrictions, and let
$b:\operatorname{Loc}(U)\to\mathcal T^*$ be the other restriction.
Systems $(L,E,\varphi:aE\to bL)$ give exactly the constructible
sheaves on $X$, compatibly with the two restriction functors, if and
only if $r$ is an equivalence and every natural map
\[
 \operatorname{Hom}_{\mathcal T^*}(aE,bL)
 \longrightarrow\operatorname{Hom}_Y(rE,i^{-1}j_*L)
\]
is bijective.
\end{proposition}
\begin{proof}
A representative of $\varphi$ on one neighborhood gives the map on
the right by adjunction and restriction to $Y$; further shrinking
changes nothing. If $r$ is an equivalence, normalize the second
entry by a quasi-inverse. Both descriptions are then categories over
$\operatorname{Loc}(U)\times\operatorname{Loc}(Y)$ with fibers the
displayed sets and arrow conditions given by equality of composites.
The natural bijections identify these categories. Conversely take $L$
terminal; then the systems form $\mathcal T$ and the corresponding
sheaves form $\operatorname{Loc}(Y)$, forcing $r$ to be an equivalence.
After normalization, the fibers over a fixed pair, with identity
base arrow, are discrete. Isomorphisms of base pairs act by conjugating
the gluing arrow in both categories. An equivalence over the base
therefore induces equivalences on these fibers, which are exactly
the asserted bijections.
\end{proof}

\begin{counterexample}[An infinite stalk need not have one tube representative]\label{ex:earring}
Let $X$ be the Hawaiian earring, with circles of center $(1/n,0)$
and radius $1/n$, and $Y=\{0\}$. Its complement is a disjoint union
of open intervals. Let $L$ be constant with two elements on that
complement, and let $E=\mathbb N$ at $Y$. For each $n$, choose a
punctured neighborhood excluding the point opposite $0$ on the
$n$th circle. There is a locally constant function equal to $1$ on
its upper arm, $0$ on its lower arm, and $0$ on every other circle.
Its germ $a_n$ defines a map $E\to(i^{-1}j_*L)_0$, hence a sheaf by
Proposition~\ref{prop:recollement}.
A tube representative would make the $\mathbb N$-local system
constant on one neighborhood of $0$ and represent all the $a_n$
there. Every such neighborhood contains an entire sufficiently
small circle. On its connected punctured circle the image of the
corresponding element $n$ must be constant, contradicting its two
arm values. Relabelling $\mathbb N$ cannot remove this contradiction.
Thus uniform representability is a substantive hypothesis.
\end{counterexample}
A separate monodromy qualification is necessary even for one stratum:
on a nondiscrete Cantor space the fundamental groupoid has only
identity arrows. Its representations can assign one element to a
single nonisolated point and two elements elsewhere, whereas no
locally constant sheaf has those stalk cardinalities.

If the relevant spaces are connected and $p_*:D=\pi_1(N^*)\to
B=\pi_1(Y)$ is surjective, write $I=\ker p_*$. A $D$-equivariant
gluing map from a $B$-set $E$ to the restricted $\pi_1(U)$-set $L$
is exactly a $B$-map $E\to L^I$: its image must be fixed by $I$,
and the action on fixed points factors through $D/I$. For a fibration
with connected fiber $K$, the exact sequence
\[
 \pi_2(Y)\longrightarrow\pi_1(K)\longrightarrow D
 \longrightarrow B\longrightarrow1
\]
shows why $I$ is a quotient of $\pi_1(K)$, not always that group.
If base and fiber are aspherical, the remaining exact sequence makes
the total space aspherical. For disconnected fibers one retains the
fundamental groupoids and does not assume surjectivity of $p_*$.

\section{Shape, exit paths, and the information in each truncation}\label{sec:shape-exit}
Shape belongs to the category of infinity-topoi. Its behavior cannot
be inferred from an arbitrary bicategorical colimit of ordinary topoi.

\begin{proposition}[Shape of a higher gluing]\label{prop:shape}
If $\mathcal X=\colim_a\mathcal X_a$ is a small colimit of
infinity-topoi, then
$\Pi_\infty\mathcal X\simeq\colim_a\Pi_\infty\mathcal X_a$
in pro-spaces. In particular a genuine open cover gives its full
higher descent formula. For synchronized tube choices this holds
naturally at every stage; it does not exchange an infinite colimit
with the neighborhood pro-system.
\end{proposition}
\begin{proof}
For a space $K$, the defining formula for shape and the universal
property of the slice topos give
\[
 \Map(\Pi_\infty\mathcal X,K)
 \simeq\Gamma_{\mathcal X}\Delta K
 \simeq\Map_{\operatorname{Top}_\infty}(\mathcal X,\mathcal S_{/K}).
\]
Mapping a colimit into this slice turns it into a limit. Evaluation
on every constant target $K$ detects equivalences of pro-spaces,
proving the result. This is the shape adjunction of
\cite[Remark~7.1.6.15]{HTT}. For an open cover, descent for sheaves
of spaces supplies the required infinity-topos colimit. Its
augmentation is natural under shrinking. Keeping these equivalences
stage by stage proves the synchronized assertion without interchanging
the two indexing operations.
\end{proof}
Finite homotopy colimits indexed by finite simplicial sets may be
computed after common cofiltered reindexing. A category with finitely
many objects and arrows need not have finite nerve. Accordingly the
untruncated higher statement uses the full Cech diagram, not just
degrees zero through two. Coherent hypercovers model the shape of
the hypercomplete sheaf infinity-topos \cite[Proposition~5.1]{Hoyois};
the distinction is harmless for the locally contractible examples
here but cannot be discarded for arbitrary topoi.

\begin{counterexample}[An ordinary-topos pushout loses higher shape]\label{ex:ordinary-shape}
Let $G=C_2$ and let $BG$ be its ordinary topos of sets with $G$-action.
The bicategorical pushout of $\mathbf{Set}\leftarrow BG\to
\mathbf{Set}$ is $\mathbf{Set}$. Indeed the inverse image of
$BG\to\mathbf{Set}$ is the fully faithful trivial-action functor.
For any target topos, a cone is therefore a pair of morphisms from
$\mathbf{Set}$ with an isomorphism between them, a category equivalent
to the category of one such morphism. The shapes of the two point
topoi are points and the shape of $BG$ is the classifying space of
$G$. Their homotopy pushout is $\Sigma BC_2$, whose reduced second
integral homology is $H_1(BC_2;\mathbb Z)=\mathbb Z/2$. The shape of
the ordinary-topos pushout is a point, so these cannot agree.
\end{counterexample}

\begin{proof}[Proof of Theorem~\ref{thm:truncation}]
The untruncated equivalence is Lurie's exodromy theorem
\cite[Theorem~A.9.3, p.~1457]{Lurie}. For finite depth its proof can
be organized as follows. Conical open charts form a basis whose
pointwise membership categories are filtered after taking opposites.
Exit excision and sheaf descent reduce the comparison to one chart
$V=Z\times C(K)$, with complement $V^*=Z\times K\times(0,\infty)$.
Induct on its number of labels. In the exit-functor category distinguish
objects extended by the initial value from $V^*$ and local systems
pulled back from $Z$; in the constructible-sheaf category take the
corresponding subcategories $j_!\operatorname{Cons}(V^*)$ and the local
systems from $Z$. The first comparison is an equivalence by induction;
the second is monodromy on $Z$, because radial contraction identifies
the local systems on the cone chart with those on $Z$.
The complementary colocalization of a local system is restriction to
$V^*$ followed by initial extension in both categories. The comparison
preserves this operation and pushouts. The recollement comparison
criterion therefore glues the two equivalences. These are exactly
Lurie's Proposition~A.8.14 and Lemmas~A.9.11--A.9.15, A.9.18.
Taking the limit over the basis glues the chart equivalences. For an
infinite ascending-chain-condition poset the cited theorem uses rank
induction; the finite-depth argument alone is not asserted to prove
that extension. It identifies stalks with
evaluations \cite[Proposition~A.9.16]{Lurie}. Constructible objects
under these hypotheses are hypercomplete, so truncation is detected
on stalks; alternatively use the local conical models. This restricts
the equivalence to the displayed full subcategories, as in
\cite[Remark~A.9.7]{Lurie}. Mapping spaces between $n$-truncated
spaces are $n$-truncated, being limits of such spaces. The universal
property of $h_{n+1}$ therefore gives the second equivalence on both
objects and mapping spaces
\cite[Propositions~2.3.4.12(4) and~2.3.4.18, pp.~110--113]{HTT}.
\end{proof}
The regularity assumptions exclude the one-stratum Cantor example.
For ordinary finite Thom--Mather or Whitney stratifications with
manifold strata they have the usual local interpretation. At $n=1$
the result agrees with Treumann's exit-path description of stacks
\cite[Theorems~7.13--7.14]{Treumann}. For a cone chart with compact
link $K$, the nearby value is the limit of the open-stratum functor
over $\operatorname{Exit}(K)$, the right Kan extension to the cone
point. This remains $n$-truncated when its values are $n$-truncated.

\begin{example}[A sphere and its gerbes]\label{ex:gerbes}
For $(\C,\{0\})$ an ordinary constructible sheaf is a set $E$,
a $\mathbb Z$-set $M$, and a map $E\to M^{\mathbb Z}$. The associated
ordinary exit category has an initial closed point, an open point with
automorphism group $\mathbb Z$, and one arrow from closed to open;
its classifying space is contractible. For $(S^2,\{p\})$ the open
stratum is contractible, so ordinary constructible sheaves reduce to
one map $E\to M$ and again have a contractible classifying category.
The genuine tube span instead has shapes $*\leftarrow S^1\to *$
and reconstructs $S^2$. Its exit mapping space from closed to open
has type $S^1$: $\pi_0$ loses it, while its fundamental groupoid
retains $\mathbb Z$.

For a discrete abelian group $A$, an $A$-banded gerbe on $S^2$ is
neutral on each of a disk and its complementary plane. An equivalence
of their neutral gerbes on the annular overlap is an $A$-torsor,
classified by $H^1(S^1;A)=\operatorname{Hom}(\mathbb Z,A)=A$.
Changing a neutralization acts by torsors on the two contractible
pieces, whose $H^1$ vanishes. Thus distinct overlap classes remain
distinct, and effective open-cover two-descent realizes them all.
This proves $H^2(S^2;A)\cong A$ in the gerbe description. The
standard correspondence with abelian degree-two cohomology is
\cite[Theorem~2.7 and Remark~3.1]{Deopurkar}. A general constructible
stack has a possibly noninvertible functor to nearby data; it need
not be a gerbe.
\end{example}

\subsection{When groupoids retain the full homotopy type}
If $X\simeq Y\cup^h_H U$ and every component of $Y,H,U$ is an
aspherical CW homotopy type, then
\[
 X\simeq B\Pi_1Y\cup^h_{B\Pi_1H}B\Pi_1U.
\]
Indeed the natural maps $\operatorname{Sing}(T)\to N\Pi_1T$ are
weak equivalences for those three spaces; applying homotopy colimit
to this natural equivalence of spans proves the formula. For the
sphere the span is $1\leftarrow\mathbb Z\to1$, and its homotopy
pushout is $S^2$. The full groupoid span retains information that the
ordinary constructible-sheaf category in Example~\ref{ex:gerbes}
forgets.

For more strata retain every nonempty strict chain
$\sigma=(p_0<\cdots<p_m)$ and its space
$D(\sigma)=\Map_{/P}(\sigma,\operatorname{Exit}(X))$.
Restriction along subchains gives a coherent contravariant diagram.
Its Segal equivalence is
\[
 D(\sigma)\simeq D(p_0<p_1)\times^h_{D(p_1)}\cdots
       \times^h_{D(p_{m-1})}D(p_{m-1}<p_m).
\]
If all stratum and all incident-pair spaces are $1$-types, then all
$D(\sigma)$ are $1$-types, since homotopy limits preserve this
property. One sufficient geometric condition is that every incident
holink source is a fibration over an aspherical stratum with aspherical
fiber. The fibration long exact sequence then makes its total incidence
space a $1$-type; for disconnected fibers this assertion is applied
componentwise, without assuming surjectivity on fundamental groups. Their fundamental groupoids, with every restriction functor
and comparison isomorphism, determine $X$ by
\[
 \operatorname{Sing}(X)\simeq
 \colim_{\sigma\in\operatorname{sd}(P)^{\mathrm{op}}}D(\sigma)
 \simeq\colim_{\sigma\in\operatorname{sd}(P)^{\mathrm{op}}}
 B\Pi_1D(\sigma).
\]
For the first comparison, extend the strict-chain diagram to all
simplices by inserting equivalences at repeated indices. This is
the complete Segal nerve; its realization is groupoid completion.
This description is the subdivision comparison of
\cite[Lemma~2.7.3 and Theorem~2.7.4]{BGH}. Groupoid completion of
exit paths is the singular homotopy type by
\cite[Corollary~A.9.4]{Lurie}. The second comparison is the natural
objectwise $1$-truncation equivalence. Without the incident-pair
hypothesis one keeps the spaces $D(\sigma)$, or their full Postnikov
towers and coherent maps. For example, at a boundary point of a
closed four-ball, the whole normal link is a three-ball, but the
incidence link to the remaining boundary stratum is $S^2$.
Asphericity of whole normal links alone would miss that difference.

Explicitly the groupoid data include functors $F_{\sigma\tau}$
and comparison isomorphisms
$\theta_{\sigma\tau\upsilon}:F_{\sigma\tau}F_{\tau\upsilon}
\Rightarrow F_{\sigma\upsilon}$, satisfying
\[
 \theta_{\sigma\tau\nu}(F_{\sigma\tau}\theta_{\tau\upsilon\nu})
 =\theta_{\sigma\upsilon\nu}
   (\theta_{\sigma\tau\upsilon}F_{\upsilon\nu}).
\]
Identity comparisons are identities. Taking one object per connected
groupoid turns these isomorphisms into conjugating elements, which
must be retained. There are no additional layers of modifications
between natural transformations of groupoids.

The trivial circle bundle over $S^2$ and the Hopf bundle show why.
Over a good cover, every nonempty intersection has total space of
type $S^1$, and every inclusion induces the identity of $\mathbb Z$
after orienting the fiber. Thus their diagrams of individual groups
and homomorphisms in the homotopy category coincide. If circle-valued
transition functions have real lifts $\theta_{ij}$, their triple
comparison is the constant integer
$n_{ijk}=\theta_{ij}+\theta_{jk}-\theta_{ik}$. Changing lifts adds a
coboundary. The trivial bundle gives zero; the Hopf bundle gives a
generator, since its equatorial clutching map has winding number one.
Their total spaces have fundamental groups $\mathbb Z$ and $0$.
The missing coherence therefore changes even the fundamental group
of the assembled space.


\section{Comparison axioms and homotopy descent}\label{b.5.-comparison-axioms-and-homotopy-descent}

Let \(A\) be a small ordinary category, \(\mathcal X\) a Grothendieck
topos, and \(\varphi:A\to\mathcal X\) a functor. Put \(U_a=\varphi(a)\)
and \(\widehat A=\operatorname{Fun}(A^{\mathrm{op}},\mathbf{Set})\). Its
cocontinuous extension is

\[L=\varphi_!: \widehat A\to\mathcal X,\qquad L(F)=\int^{a\in A}F(a)\cdot U_a.\]

Here the dot denotes copower. For a finite diagram \(F:J\to\widehat A\),
write \(c_F:L(\lim_JF)\to\lim_JL(F)\) for the canonical comparison. An
object \(U\) of \(\mathcal X\) represents the slice topos
\(\mathcal X/U\), not its set of global sections. This distinction is
essential for interpreting the homotopy-colimit comparison.

\subsection{The forest statement and the full
condition}\label{b5.1.-the-forest-statement-and-the-full-condition}

Hyp 1, 2 and \(2'\) for the comparison scheme are
respectively the epimorphicity of

\[L(1)\to1,\qquad L(ya\times yb)\to U_a\times U_b,\]

and

\[L(ya\times_{yc}yb)\to U_a\times_{U_c}U_b\]

for every cospan \(a\to c\leftarrow b\) in \(A\). These are equivalent
to \(c_F\) being epic for all diagrams indexed by the free category on a
finite oriented forest, including the empty forest. This is an
equivalence with a collection of tests; it does not assert that every
indexing category which is not a forest fails. For example, a category
with an initial object has its limit given by evaluation there.

\textbf{Proof.} The category of elements gives
\(L(F)=\operatorname{colim}_{(a,s)\in A/F}U_a\). The coproduct of its
terms maps epimorphically onto that colimit. Thus the image of each
representable comparison is exactly the image of the covering family of
cones written in the proposed scheme. Hyp \(1'\) asks that arrows into an
object cover it; its identity suffices.

For completeness we prove the extension to arbitrary presheaves. Work
locally in \(\mathcal X\). An element of \(L(F)\) is represented by
\((a,x,s)\) with \(x\in U_a\) and \(s\in F(a)\), and equality is locally
witnessed by a finite zigzag of coend relations. In the internal
category \(E_\varphi\), an object is \((a,x)\) and an arrow
\((a,x)\to(b,z)\) is \(f:a\to b\) such that \(\varphi(f)x=z\). Hyp 2
says any two objects locally have a common predecessor; Hyp \(2'\) says
a cospan locally has a commutative completing square. Every finite
zigzag locally admits a commuting cone: induct on its edges, composing a
cone leg at a forward step and using a completing square at a backward
step. For disconnected diagrams combine the cones by Hyp 2. Only
finitely many covering refinements are used.

Given an element of \(L(G)\times_{L(F)}L(H)\), represent its two entries
locally. Their images are connected by such a zigzag. The cone replaces
them by representatives at a common object with equal images in \(F\),
hence by an element of \(G\times_FH\). This proves the pullback
comparison is epic for arbitrary presheaves; products use the same
argument without the equality zigzag. A tree limit is obtained by
deleting leaves and successively taking pullbacks: a leaf which is a
target contributes no new variable, and a leaf which is a source gives a
pullback over the adjacent vertex. The comparison factors through a
pullback comparison and a base change of the earlier epimorphism.
Products treat components; Hyp 1 treats the empty forest. The converse
uses the empty diagram, the discrete pair, and the cospan. The argument
uses the local meaning of epimorphisms, not enough points. \(\square\)

\textbf{Proof of the finite-comparison clause of Theorem~\ref{thm:higher-comparisons}.} Flatness means left exactness of \(L\).
Let \(M=\lim_JF\) and \(c:L(M)\to N=\lim_JL(F)\). Add
two cone vertices to \(J\), both carrying \(M\) and its canonical cone.
The enlarged diagram has limit \(M\times_M M=M\). After \(L\) its limit
is \(L(M)\times_NL(M)\), so its comparison is the diagonal of \(c\).
That diagonal is epic by hypothesis and monic automatically, hence
invertible. Therefore \(c\) is monic and, by the original hypothesis,
epic; it is an isomorphism. This also treats empty \(J\). The converse
is immediate. \(\square\)

\subsection{A finite
counterexample}\label{b5.2.-a-finite-counterexample}

Hyp 1, 2 and \(2'\) neither
imply full (Hyp) nor suffice for homotopy descent over a point.

\textbf{Proof.} Take \(A=BC_2\) and \(\varphi(*)=\{*\}\). Then \(L\)
takes a \(C_2\)-set to its orbit set. Products compare surjectively by
choosing representatives. For a pullback of equivariant maps,
representatives of two orbits with the same image orbit can be adjusted
by a group element so their images agree; therefore all pullback
comparisons are surjective. On the regular \(C_2\)-set, however, the
identity and translation by the nonidentity element have empty
equalizer, while their maps on the one-point orbit set coincide. Its
equalizer comparison is \(\varnothing\to\{*\}\). Finally,
\(\operatorname{hocolim}_{BC_2}\{*\}=BC_2\). Its normalized simplicial
chain complex over \(\mathbf F_2\) has one generator in degrees 1 and 2
and zero boundary between them, so \(H_1=\mathbf F_2\) and it is not
contractible. \(\square\)

\subsection{The symmetric group nerve: distinguish three
realizations}\label{b5.3.-the-symmetric-group-nerve-distinguish-three-realizations}

Let \(\mathsf F\) be nonempty finite sets and all maps,
\(A=\mathsf F^{\mathrm{op}}\), and \(G\) a discrete group. Put
\(U(S)=G^S/G\) for simultaneous left multiplication. Hyp 1, 2 and \(2'\)
hold, and

\[\operatorname{hocolim}_{\mathsf F^{\mathrm{op}}}U\simeq BG.\]

For \(G\ne1\), full (Hyp) fails. The ordinary simplicial restriction
also realizes to \(BG\). In contrast, the literal symmetric coend
realization with affine simplices is contractible for finite \(G\).

\textbf{Proof.} We have \(U(1)=*\). Concatenation of representatives
proves the product condition. A cospan in \(A\) is opposite to
\(P\leftarrow R\to Q\) in \(\mathsf F\), whose pushout represents its
pullback in \(A\). If tuple orbits on \(P,Q\) have the same restricted
orbit on \(R\), translate one representative so the restrictions agree.
They descend to a tuple on the pushout. This proves the condition for
every cospan, including noninjective maps. For the two coordinate maps
\(2\rightrightarrows1\) in \(A\), the equalizer is \(1\to2\). Its image
under \(U\) is \(*\to G\), but the equalizer after \(U\) is all of
\(G\).

Put \(T(S)=G^S\). Its category of elements is cofiltered: it is
nonempty, concatenation gives common predecessors, and parallel tuple
maps are equalized by quotienting their finite coordinate set by the
required identifications. The tuple values agree on those
identifications precisely because the two arrows have the same target
tuple. The quotient remains nonempty. A cofiltered category has
contractible nerve: give a common predecessor to the vertices of any
finite subcomplex, then successively equalize the finitely many edge
compatibility equations to obtain a cone. This contracts maps from
finite simplicial spheres.

The \(G\)-action on this element category is free on objects, and on
nerve simplices, since tuples are nonempty. Choosing a representative of
the initial tuple uniquely lifts a whole chain. Thus
\(N(E_U)=N(E_T)/G\), a free contractible simplicial \(G\)-space modulo
\(G\), hence \(BG\). The category-of-elements nerve is the simplicial
replacement computing a homotopy colimit. Restricting to
\(\Delta^{\mathrm{op}}\) gives the homogeneous bar construction
\(G^{n+1}/G\).

For finite \(G\), the symmetric realization of \(T\) is instead the
simplex \(\Delta(G)\): equal labels merge their barycentric weights,
leaving a probability vector on \(G\). Coends commute with orbit
quotients, so the symmetric realization of \(U\) is \(\Delta(G)/G\). The
affine contraction to the invariant barycenter is equivariant, hence
contracts the quotient. For \(C_2\) this is an interval modulo
reflection. \(\square\)

\subsection{Classification of flat symmetric
systems}\label{b5.4.-classification-of-flat-symmetric-systems}

Flat functors
\(\mathsf F^{\mathrm{op}}\to\mathcal X\) are exactly the systems
\(S\mapsto E^S\) with \(E\to1_{\mathcal X}\) epic. In sets, \(E\) is a
nonempty set.

\textbf{Proof.} Let \(E=\varphi(1)\). The object \(n\) is the product of
\(n\) copies of \(1\) in \(A\). Yoneda and left exactness of \(L\) give
the canonical isomorphism \(\varphi(n)\cong E^n\), naturally for all
coordinate maps. Hyp 1 gives an epimorphism \(\coprod_{n\ge1}E^n\to1\),
which factors through \(E\to1\), proving inhabitedness. Conversely, the
internal tuple category of an inhabited \(E\) is locally nonempty, has
common predecessors by concatenation, and equalizes parallel arrows by
the coordinate quotient used above. It is locally cofiltered. The coend
is internally a filtered colimit of evaluations: finitely many
representatives can be moved to one index, and finitely many equalities
made simultaneously true at a further index. This proves both
surjectivity and injectivity of every finite-limit comparison. All
arguments use finitely many local refinements and apply to any
Grothendieck topos. \(\square\)

\subsection{Flatness, forest comparisons, and Verdier
hypercovers}\label{b5.5.-flatness-forest-comparisons-and-verdier-hypercovers}

A hypercover condition agreeing with Verdier's cannot in
general lie between the proposed scheme's forest conditions and flatness.

For \(A=\Delta^{\mathrm{op}}\), the ordinary Cech nerve
\(U_n=\{0,1\}^{n+1}\) of a cover of a point is a hypercover, and its
realization is the contractible nerve of the indiscrete groupoid on two
objects. Yet Hyp \(2'\) fails. Take in \(\Delta\) the span
\([0]\leftarrow[1]\to[2]\), with the first map constant and the second
selecting endpoints 0,2. The corresponding fiber product contains
\((0,(0,1,0))\). A cocone must have a monotone map \([2]\to[m]\) with
equal endpoint values, hence a constant map. Every tuple lifted from a
completing cone is therefore constant in those three positions. The
displayed element has no lift. In particular full (Hyp) fails as well.

Conversely, \(\Delta[1]\) is a flat functor on \(\Delta^{\mathrm{op}}\):
its extension is evaluation at \([1]\). It is not a hypercover of a
point, because its degree-one matching map omits the ordered pair of
vertices \((1,0)\). Thus even flatness and Verdier's condition are
incomparable for ordinary simplicial indexing.

The failure of ``between'' persists for symmetric indexing. Put
\(U(S)=\operatorname{Cont}(\Delta(S),\mathbf R)\), with maps by affine
precomposition. Every continuous boundary map to \(\mathbf R\) extends
to the simplex by coning from the barycenter and interpolating to zero;
compactness gives continuity at the barycenter. Hence the ordinary
simplicial restriction is a hypercover of a point. For the finite-set
span \(1\leftarrow2\to3\), with the right map selecting vertices 0,2,
take zero on \(\Delta^0\) and

\[b(t_0,t_1,t_2)=t_1(t_0-t_2).\]

They agree on edge 02. Any amalgamating finite-set map identifies
vertices 0,2, so its affine map takes \((1/2,1/2,0)\) and
\((0,1/2,1/2)\) to the same point. A lifted function must have equal
values there; \(b\) has values \(1/4,-1/4\). Thus Hyp \(2'\) fails.
The displayed argument works for every target cardinality.

There is nevertheless a one-way implication specific to symmetric
indexing: full (Hyp) implies the hypercover condition. By Section~\ref{b5.4.-classification-of-flat-symmetric-systems} its system
is \(U(S)=E^S\) with \(E\to1\) epic. The simplicial restriction is the
Cech nerve of this epimorphism; every positive matching map is an
isomorphism. The symmetric continuous-simplex example just given proves
the converse false. Thus ordinary flatness and hypercovers are
incomparable, whereas symmetric flatness is strictly stronger than
hypercovers; in either case hypercovers need not satisfy the forest
conditions.

Over sets there is also an explicit ordinary-index classification: flat
functors \(\Delta^{\mathrm{op}}\to\mathbf{Set}\) are precisely nerves of
nonempty linear orders. Indeed a flat functor is a filtered colimit of
the covariant representables
\(\operatorname{Hom}_{\Delta^{\mathrm{op}}}([n],-)=\Delta[n]\). A
filtered colimit of finite ordinals and monotone maps is a linear order:
comparisons and antisymmetry are witnessed at a common stage. Nerves
commute with filtered colimits, since finite strings use finite limits
of sets. Conversely a linear order is the filtered union of its nonempty
finite suborders. The filtered-colimit characterization follows from the
cofiltered element-category argument in Section~\ref{b5.4.-classification-of-flat-symmetric-systems}.~\(\square\)

\subsection{Exact general homotopy-colimit
criterion}\label{b5.6.-exact-general-homotopy-colimit-criterion}

Let
\(\mathscr X\) denote either the associated 1-localic infinity-topos or
its hypercompletion; fix this choice in each formula. Regard \(U_a\) as
0-truncated objects of \(\mathscr X\). Form

\[B_\varphi{}_n=\coprod_{a_0\to\cdots\to a_n}U_{a_0},\qquad
H_\varphi=|B_\varphi|=\operatorname{colim}_{a\in A}U_a.\]

The first face uses \(U_{a_0}\to U_{a_1}\), inner faces compose arrows,
the last face drops the last arrow, and degeneracies insert identities.
The augmentation \(\epsilon:H_\varphi\to1\) is the terminal comparison
of the derived colimit extension
\(\mathbf L\varphi_!:\mathcal P(A)\to\mathscr X\), since
\(1_{\mathcal P(A)}=\operatorname{colim}_{a\in A}y(a)\).

We prove the shape clause of Theorem~\ref{thm:higher-comparisons}, including the sufficient condition supplied by finite-limit preservation before hypercompletion.

\textbf{Proof.} A constant object in the slice over \(U\) is
\(U\times\Delta K\to U\), and therefore

\[\operatorname{Map}_{\operatorname{Pro}(\mathcal S)}(\Pi_\infty(\mathscr X/U),K)
=\operatorname{Map}_{\mathscr X}(U,\Delta K).\]

Maps from the colimit of slice shapes to \(K\) give the limit on the
right of (HS2), which is
\(\operatorname{Map}_{\mathscr X}(H_\varphi,\Delta K)\). Shape is the
functor \(K\mapsto\Gamma\Delta K\), and pro-spaces are fully faithfully
represented by these functors. This proves the exact criterion and its
sufficient condition.

Under full (Hyp), Section~\ref{b5.1.-the-forest-statement-and-the-full-condition} makes \(L\) left exact; it is cocontinuous and
has right adjoint \(V\mapsto[a\mapsto\operatorname{Hom}(U_a,V)]\). Thus
it is the inverse image of a geometric morphism
\(\mathcal X\to\widehat A\). The universal property of the associated
1-localic infinity-topos
\cite[Proposition~6.4.5.7]{HTT} first extends this to
\(\mathcal X_\infty\to\mathcal P(A)\). Its inverse image agrees with
\(\varphi\) on representables and hence is the derived colimit extension
into \(\mathcal X_\infty\). It preserves the terminal object. If
\(\mathscr X=\widehat{\mathcal X_\infty}\), compose this inverse image
with hypercompletion, which preserves colimits and finite limits. This
proves \(H_\varphi=1\) for either choice without asserting that the
hypercompletion is itself 1-localic. \(\square\)

In the hypercomplete choice, \(H_\varphi=1\) is equivalent to the
augmented simplicial replacement being a local weak equivalence. After
local fibrant replacement this can be tested by local relative homotopy
fillings of all finite simplicial boundaries (Dugger--Hollander--Isaksen
Proposition 3.1 \cite[Propositions 3.1--3.2]{DHI}); strict local
fillings characterize a local acyclic fibration (Proposition 3.2). These
are higher coherence conditions, not the ordinary forest comparisons.

\textbf{Exact Artin--Mazur interpretation.} Require \(\mathcal X\)
locally connected, and let
\(q:\operatorname{Pro}(\mathcal S)\to
\operatorname{Pro}(\operatorname{Ho}\mathcal S)\) be the canonical
functor. The hypercover pro-homotopy type is
\[
\operatorname{AM}(\mathcal X)\cong
q\bigl(\Pi_\infty(\widehat{\mathcal X_\infty})\bigr).
\tag{HS3}
\]
Indeed Barnea--Harpaz--Horel Theorem 6.0.4, printed pp.~625--626,
identifies the derived topological realization with the shape of the
\emph{hypercompleted} infinity-topos \cite{BHH}.
Barnea--Schlank Definition 8.2, Proposition 8.4 and Remark 8.5(2),
pp.~62--65 of the cited preprint, identify its image under
\(q\) with the Artin--Mazur construction for every locally connected
topos \cite{BS}. Their Remark 8.5(3) separately records
the further localization used by Artin--Mazur; applying that localization
preserves all isomorphisms proved here.

Thus (HS1), formed in pro-spaces using the hypercomplete choice and
then sent through \(q\), is the precise single-diagram comparison with
\(\operatorname{AM}(\mathcal X)\). It is an isomorphism in
\(\operatorname{Pro}(\operatorname{Ho}\mathcal S)\) if and only if, for
every space \(K\), the map
\[
\pi_0\operatorname{Map}_{\widehat{\mathcal X_\infty}}(1,\Delta K)
\longrightarrow
\pi_0\!\left(\lim_{a\in A^{\mathrm{op}}}
\operatorname{Map}_{\widehat{\mathcal X_\infty}}(U_a,\Delta K)\right)
\tag{HS4}
\]
is a bijection. To prove this, a pro-space \(P=(P_i)\) satisfies
\(\operatorname{Hom}(qP,K)=\operatorname{colim}_i[P_i,K]
=\pi_0\operatorname{Map}(P,K)\), since filtered colimits commute with
\(\pi_0\). Pro-objects of \(\operatorname{Ho}\mathcal S\) are fully
faithfully detected by their set-valued functors of maps to constant
objects. Apply this to the already formed comparison (HS1).
The limit in (HS4) is a limit of spaces taken \emph{before} \(\pi_0\).
We neither interchange these operations nor claim that \(q\) preserves
the colimit in (HS1). The mapping-space equivalences (HS2) imply
(HS4), but the latter alone is not asserted to recover them.

For the full, possibly nonhypercomplete shape, use (HS2) with
\(\mathscr X=\mathcal X_\infty\). A local weak equivalence or a
hypercover alone supplies the hypercomplete version, and does not by
itself justify that full-shape conclusion. It also suffices for the
protruncated comparison, because truncated constant objects are
hypercomplete. Full ordinary (Hyp) supplies the stronger terminal-object
descent before hypercompletion, as proved above. These are distinct
targets, rather than implicit identifications of pro-homotopy objects
and pro-spaces; see also HTT Example 7.1.6.9 \cite{HTT}.

\subsection{Verdier matching as comparisons of
representables}\label{b5.7.-verdier-matching-as-comparisons-of-representables}

Define

\[M_nU=\lim_{\varnothing\ne S\subsetneq\{0,\ldots,n\}}U_{|S|-1},\qquad M_0U=1,\]

with face restriction arrows. The conditions \(U_n\to M_nU\) epic for
all \(n\ge0\) are the hypercover conditions. The realization is
infinity-connected in \(\mathcal X_\infty\) and terminal after
hypercompletion. This is HTT Definition 6.5.3.2, Lemma 6.5.3.11 and
Theorem 6.5.3.12 \cite{HTT} (or DHI Definition 4.2 and Proposition
3.2).

For \(A=\Delta^{\mathrm{op}}\) and \(n\ge2\), the limit in
\(\widehat A\) of the representables of all proper nonempty faces of
\([n]\) is \(y[n]\). Consequently its comparison \(c_F\) is precisely
\(U_n\to M_nU\). Degrees zero and one must be added separately. For
\(A=\mathsf F^{\mathrm{op}}\) the same statement holds for all
\(n\ge1\), so all positive-degree matching conditions are the
corresponding boundary comparisons, with degree zero \(U(1)\to1\).

\textbf{Proof of the representable identity.} Compatible maps from all
proper faces into an ordinal determine the vertex images. For \(n\ge2\),
every pair of vertices lies in a proper face; the images are therefore
monotone, determining a unique map from \([n]\). At \(n=1\) the two
vertex images need not be monotone. For finite sets, no monotonicity
restriction occurs, and compatible vertex maps determine a unique
function also in the two-vertex case. Yoneda turns these colimit
properties into the asserted limits. \(\square\)

The inclusion \(j:\Delta^{\mathrm{op}}\to\mathsf F^{\mathrm{op}}\) is
homotopy final. For any nonempty \(S\), the comma category
\((S\downarrow j)\) is the category of nonempty ordered tuples in \(S\)
and monotone restriction maps. Choose \(s\in S\). Prepending \(s\)
defines an endofunctor with natural transformations to the identity
(drop the first entry) and to the constant tuple \((s)\) (select the
first entry). Its nerve is contractible. The cofinality criterion HTT
Theorem 4.1.3.1 applies. Thus symmetric matching conditions compute the
homotopy colimit over all finite-set maps, not only the ordinary
simplicial subcategory.

\subsection{Every matching degree is independently
needed}\label{b5.8.-every-matching-degree-is-independently-needed}

For every \(m\ge0\) there is a simplicial set
with all matching maps except the \(m\)th surjective and with
noncontractible realization. The example extends to symmetric indexing
by all nonempty finite sets.

\textbf{Proof.} For \(m=0\), take the empty simplicial set: every
positive matching target is empty. For \(m=1\), take the constant
two-point set: only the diagonal map to its square fails, since all
higher simplex boundaries are connected.

For \(m\ge2\), let \(n=m-1\) and define
\(K_n(k)=Z^n(\Delta^k;\mathbf F_2)\) with cochain pullback maps. A
simplicial map \(T\to K_n\) is precisely a normalized \(n\)-cocycle on
\(T\): its values on \(n\)-simplices form the cochain and the
\((n+1)\)-simplex relation is the cocycle equation. Hence the \(k\)th
matching map is

\[Z^n(\Delta^k;\mathbf F_2)\to Z^n(\partial\Delta^k;\mathbf F_2).\]

For \(k<n\) both groups vanish; for \(k=n\) the target vanishes. At
\(k=n+1\), the target consists of arbitrary values on the \(n+2\)
boundary \(n\)-faces; the image is the codimension-one subspace whose
sum is zero. For \(k>n+1\), every \(n\)- and \((n+1)\)-simplex already
belongs to the boundary, so restriction is an isomorphism. Thus only
degree \(m\) fails.

The normalized Moore complex of \(K_n\) is \(\mathbf F_2\) in degree
\(n\) and zero otherwise: a normalized element vanishes on all specified
faces, and in degree \(n+1\) the cocycle equation forces the remaining
face to vanish; above that degree all \(n\)-faces already lie in
vanishing faces. Its homotopy group \(\pi_n\) is therefore
\(\mathbf F_2\), by Curtis, Theorem 3.7 (printed p.~131)
\cite{Curtis}. This is also the cocycle model of the Dold--Kan
object of \(\mathbf F_2[n]\); normalization is the equivalence of Stacks
Theorem 14.24.3 (019G) \cite{StacksDK}. For the symmetric extension,
define the same cocycles on the abstract simplex on a finite set; an
arbitrary vertex map induces a normalized simplicial chain map, sending
a face with repeated image vertices to zero. Over \(\mathbf F_2\) this
construction is independent of ordering. It gives the same ordinary
simplicial restriction, hence the same homotopy colimit by Section~\ref{b5.7.-verdier-matching-as-comparisons-of-representables}. \(\square\)

\subsection{Component sets alone cannot reconstruct the
homotopy
type}\label{b5.9.-component-sets-alone-cannot-reconstruct-the-homotopy-type}

No condition requiring even all ordinary comparisons to be
isomorphisms can ensure \(\operatorname{hocolim}_A\pi_0(U_a)\) models
the Artin--Mazur type without a condition on the pieces.

\textbf{Proof.} Take \(A\) terminal,
\(\mathcal X=\operatorname{Sh}(S^2)\) and \(U_*=1_{\mathcal X}\). Then
\(L:\mathbf{Set}\to\operatorname{Sh}(S^2)\) is constant-sheaf inverse
image, so it preserves finite limits. But \(\pi_0(U_*)=*\), and its
homotopy colimit is a point; the sphere has
\(H^2(S^2;\mathbf Z)=\mathbf Z\). The correct single-diagram formula
uses slice shapes. Component sets can replace them when each canonical
map \(\Pi_\infty(\mathscr X/U_a)\to\pi_0(U_a)\) is an equivalence and
the diagram is homotopy effective: replace each object by that
equivalent pro-space in Section~\ref{b5.6.-exact-general-homotopy-colimit-criterion}. In general Artin--Mazur retains the
cofiltered pro-system of all hypercovers and their component simplicial
sets. Even the constant terminal hypercover is not generally a single
model of the shape. \(\square\)

\subsection{An explicit all-index local filling
criterion}\label{b5.10.-an-explicit-all-index-local-filling-criterion}

The
exact stronger condition \(H_\varphi=1\) in the hypercomplete
infinity-topos has an explicit test using the original \(U_a\), finite
limits, coproducts and a sequential colimit. It does not require one to
guess boundary objects in the arbitrary category \(A\).

Let \(N=B_\varphi\) be the simplicial replacement from Section~\ref{b5.6.-exact-general-homotopy-colimit-criterion};
equivalently it is the nerve of the internal element category, whose
objects are \(\coprod_aU_a\) and whose arrow object is
\(\coprod_{a\to b}U_a\). For a finite simplicial set \(K\), let
\([K,N]\) be the object of simplicial maps, expressed as a finite limit
using the nondegenerate simplices of \(K\). Put

\[R_n=\operatorname*{colim}_{r\ge0}[\operatorname{Sd}^{r}\Delta^n,N],\]

with the transition maps induced by the standard last-vertex subdivision
maps. Here colimits and finite limits are in \(\mathcal X\), and
\(\operatorname{Sd}\) is barycentric subdivision. The local-filling clause of Theorem~\ref{thm:higher-comparisons} is the epimorphicity of all matching maps of this derived simplicial object. It is equivalent to \(N\to1\) being a local weak equivalence.

For \(n=0\) the right side is terminal. Since filtered colimits commute
with finite limits in a Grothendieck topos, the right side is precisely
\(M_nR\). Thus (HF) is an ordinary epimorphic matching-map scheme on the
explicitly derived simplicial object \(R\). It says that every finite
boundary admits a filling locally after sufficiently many subdivisions.
The required subdivision number may vary over the covering family; no
assertion of algorithmic decidability is made.

\textbf{Proof.} Present \(\mathcal X\) by a small site, and regard \(N\)
as a simplicial sheaf there. Apply \(\operatorname{Ex}^{\infty}\)
sectionwise to obtain a simplicial presheaf \(F\). For every object of
the site, \(F\) is a Kan complex and \(N\to F\) is a weak homotopy
equivalence, by Kerodon Propositions 3.3.6.7 and 3.3.6.9 (Tags 00ZH and
00ZK) \cite{Kerodon}. Therefore \(N\to F\) is a local weak
equivalence, and \(F\) is locally fibrant. Dugger--Hollander--Isaksen
Proposition 3.2, reformulated as their equation (3.4) \cite{DHI},
says that \(F\to1\) is a local weak equivalence if and only if every map
\(F_n\to[\partial\Delta^n,F]\) is a local epimorphism: its local
fibration hypothesis is automatic here. Sheafification preserves finite
limits and colimits. Sheafifying these maps gives precisely (HF),
by the subdivision--extension adjunction and finite presentation of the
boundaries. This proves equivalence of the local weak equivalence and matching-map conditions, independently of the
chosen site. The local model category presents the hypercompletion, so
the hypercomplete augmentation condition is equivalent to the local weak equivalence condition. This last restriction matters: local weak
equivalence gives an infinity-connected augmentation in a
nonhypercomplete topos, not necessarily an equivalence there.
\(\square\)

This supplies an effective mathematical description for every small
\(A\), in addition to the direct sufficient condition on the
proposed \(c_F\) in Section~\ref{b5.1.-the-forest-statement-and-the-full-condition} and the sharper boundary schemes of Section~\ref{b5.7.-verdier-matching-as-comparisons-of-representables}.
The weaker forest scheme fails by Section~\ref{b5.2.-a-finite-counterexample}--Section~\ref{b5.3.-the-symmetric-group-nerve-distinguish-three-realizations}. Omitting any dimension
from the strict matching scheme fails by Section~\ref{b5.8.-every-matching-degree-is-independently-needed}, in ordinary and symmetric
indexing. Conditions for shape equivalence alone can be weaker than
universal local descent; their precise criteria are HS2 for pro-spaces
and HS4 for the original pro-homotopy category. In particular,
(HF) is necessary and sufficient for the stronger local descent
assertion, not a claimed necessary condition for equality only of shapes.

\subsection{Scope of the comparison
result}\label{scope-of-the-comparison-result}

The answer separates the literal ordinary comparisons, the universal
local descent condition, and the Artin--Mazur target. All finite
ordinary comparisons epic is exactly flatness and suffices even before
hypercompletion; the forest subcollection fails. The boundary
comparisons in ordinary and symmetric indexing give Verdier hypercovers,
with every matching degree indispensable. For arbitrary indexing,
(HF) is an explicit necessary-and-sufficient higher local descent
criterion after hypercompletion. For equality only of shapes, (HS2)
is the exact criterion; for the original Artin--Mazur pro-homotopy type,
(HS3)--(HS4) give the exact locally connected interpretation.
A universal minimal subset of
the ordinary finite diagram shapes for every possible \(A\) is not
asserted or needed for these results.

\section{Presentations of nonempty finite sets}\label{b.6.-presentations-of-the-category-of-nonempty-finite-sets}

Composition is written right to left. Let \(\mathsf F\) have skeleton
\(n=\{1,\ldots,n\}\), \(n\ge1\). The generators considered here
are the permutations \(\mathfrak S_n\), the inclusion
\(\delta_n:n\to n+1\) of the first \(n\) elements, and \(K_n:n+1\to n\)
which fixes \(1,\ldots,n\) and sends \(n+1\) to \(n\). Let \(i_n\)
extend a permutation by fixing the last element, and let
\(t_n=(n\ n+1)\in\mathfrak S_{n+1}\).

\subsection{The proposed global presentation is
incomplete}\label{b6.1.-the-proposed-global-presentation-is-incomplete}

The four displayed relation families for finite sets do not
present \(\mathsf F\). Adding both additional global invariance families, even generalized to all degrees, still does not suffice.

The four families, with the inclusions made explicit to fix types, are

\[K_n\delta_n=\mathrm{id},\qquad
\delta_n\sigma=i_n(\sigma)\delta_n\quad(\sigma\in\mathfrak S_n),\]
\[i_{n-1}(\sigma)K_n=K_n i_n i_{n-1}(\sigma)\quad(\sigma\in\mathfrak S_{n-1}),\qquad K_n t_n=K_n.\]

The additional families are

\[(n+1\ n+2)\delta_{n+1}\delta_n=\delta_{n+1}\delta_n,
\qquad K_nK_{n+1}(n\ n+1)=K_nK_{n+1}.\]

\textbf{Countermodel.} Interpret these contravariantly in sets. Put
\(E_1=\{0\}\) and \(E_n=\{0,1\}\) for \(n\ge2\), and let every
permutation act trivially. Put \(p_1=\delta_1^*:E_2\to E_1\) equal to
the unique map and \(d_1=K_1^*:E_1\to E_2\) equal to the inclusion of
\(0\). For every \(n\ge2\), put \(p_n=d_n=\mathrm{id}_{\{0,1\}}\). The
retraction equation is \(p_nd_n=\mathrm{id}_{E_n}\), which holds. All
equivariance and invariance equations hold because every permutation
acts trivially.

In actual finite sets, however, one has the typed equation between maps
\(2\to2\)

\[\boxed{(1\ 2)K_2(1\ 3)\delta_2=\delta_1K_1.}\tag{FS1}\]

Both sides are constant with value \(1\): the middle composite on the
left sends both elements to \(2\), and the leading transposition sends
that value to \(1\). In the contravariant model, the left side is
\(p_2d_2=\mathrm{id}_{E_2}\) and the right side is \(d_1p_1\), constant
with value \(0\). They differ at \(1\). Thus (FS1) does not follow
from the proposed presentation, even after both invariance families are
added. Only the three objects \(1,2,3\) are needed to witness the
failure; the model is defined in every degree. \(\square\)

\subsection{A complete replacement is available in the primary
literature}\label{b6.2.-a-complete-replacement-is-available-in-the-primary-literature}

Grandis, \emph{Finite sets and symmetric simplicial sets},
§3 and Theorem 4.2 \cite{Grandis}, gives a complete presentation
of finite cardinals by all elementary cofaces, codegeneracies and
adjacent transpositions, subject to the cosimplicial, symmetric-group
and mixed interchange relations. His §5 restricts the presentation to
nonempty finite cardinals. Thus a complete corrected presentation is
obtained by adding all these translated relations to the proposed scheme
generators.

Here are explicit substitution words, including their arities.  Write
\(c_i^{(m)}=(i\ i+1\ \cdots\ m)\in\mathfrak S_m\), with
\(c_m^{(m)}=1\).  For the coface \(d_i^{(n)}:n\to n+1\) omitting
\(i\), and the codegeneracy \(e_i^{(n)}:n+1\to n\) identifying
\(i,i+1\), set
\[
 d_i^{(n)}=c_i^{(n+1)}\delta_n\quad(1\le i\le n+1),\qquad
 e_i^{(n)}=c_i^{(n)}K_n\bigl(c_i^{(n+1)}\bigr)^{-2}
       \quad(1\le i\le n).
 \tag{FS2}
\]
Adjacent transpositions are \(s_i=(i\ i+1)\), already present in the
symmetric groups.  The first word fixes positions before \(i\) and
moves every later old position one step to the right.  In the second
word, the source permutation fixes \(k<i\), sends \(i,i+1\) to
\(n,n+1\), and sends \(k>i+1\) to \(k-2\); \(K_n\) merges the last
two, and the output cycle puts the merged value back in position
\(i\).  Thus these are exactly the indicated finite-set maps.
Crucially, the last words recover the old generators literally:
\(d_{n+1}^{(n)}=\delta_n\) and \(e_n^{(n)}=K_n\).

For a completely explicit enlarged relation scheme, use (FS2) in
all the following formulas, whenever their arities are positive and
well typed; dimension superscripts are suppressed:
\[
 d_i d_j=d_{j+1}d_i\ (i\le j),\qquad
 e_j e_i=e_i e_{j+1}\ (i\le j),
\]
\[
 e_jd_i=
 \begin{cases}
 d_i e_{j-1},&i<j,\\
 1,&i=j\text{ or }i=j+1,\\
 d_{i-1}e_j,&i>j+1;
 \end{cases}
\]
\[
 s_i^2=1,\quad s_i s_{i+1}s_i=s_{i+1}s_i s_{i+1},\quad
 s_i s_j=s_j s_i\ (i<j-1),
\]
\[
 \begin{array}{lll}
 s_i d_j=d_j s_i &(i<j-1),&s_i e_j=e_j s_i\quad(i<j-1),\\
 s_i d_i=d_{i+1},&&s_i e_i=e_{i+1}s_i s_{i+1},\\
 s_i d_j=d_j s_{i-1}&(i>j),&s_i e_j=e_j s_{i+1}\quad(i>j),
 \end{array}
 \qquad e_i s_i=e_i.
 \tag{FS3}
\]
These are Grandis's (3), (10), (12), (13), with indices shifted from
zero-based to one-based notation.  The omitted mixed cases follow by
\(s_i^2=1\), exactly as in Grandis.  The group equations are redundant
if the complete groups \(\mathfrak S_n\) and their multiplication laws
are retained as the proposed scheme prescribes.

We verify carefully that no unspoken defining equations are missing.
Let \(P\) be Grandis's presented category on positive cardinals and
let \(Q\) be the category on the proposed scheme generators with their group
laws, the four valid displayed equations, and all of (FS3) after the
substitutions (FS2).  Substitution defines a functor
\(F:P\to Q\), because every defining relation of \(P\) has explicitly
been imposed in \(Q\).  This functor is surjective on every hom-set:
its image contains \(\delta_n\) and \(K_n\) literally and contains all
adjacent transpositions, which generate each of the already imposed
symmetric groups.  The finite-set interpretation gives
\(G:Q\to\mathsf F\), because every imposed relation holds for these
functions.  The composite \(GF\) is the standard interpretation of
Grandis's generators and is an isomorphism by Theorem 4.2 and §5.

For faithfulness of \(G\), choose lifts \(u,v\) in \(P\) of two
parallel arrows of \(Q\).  Equality of their \(G\)-images gives
\(GF(u)=GF(v)\), hence \(u=v\), and hence equality of the arrows in
\(Q\).  Fullness follows from the same composite.  Thus
\(Q\cong\mathsf F\), with exactly the indicated original generators.
This proof does not assume that generator substitutions are inverses
without checking their defining equations.  The explicit coordinate argument proves the formulas in every arity.  We make no
minimality claim for this enlarged relation list.

\subsection{The corrected low-degree presentation is
complete}\label{b6.3.-the-corrected-low-degree-presentation-is-complete}

The low-degree relations below
present the full subcategory \(\mathsf F_{\le3}^{\mathrm{op}}\),
with the specified subscript \(1\). Replacing it by \(3\), or replacing the stabilizer relation as indicated below, gives false equations.

Write the three objects \(E_0,E_1,E_2\) for cardinalities \(1,2,3\),
respectively, and abbreviate

\[p=p_0:E_1\to E_0,\quad d=\delta_0:E_0\to E_1,
\quad q=p_{12}:E_2\to E_1,\quad e=\delta_1:E_1\to E_2.\]

The groups are \(\mathfrak S_2\) and \(\mathfrak S_3\). In the
coordinate interpretation, let

\[\sigma(x,y)=(y,x),\quad \rho(x,y,z)=(y,z,x),
\quad\sigma_1(x,y,z)=(x,z,y),\quad\sigma_3(x,y,z)=(y,x,z).\]

Thus \(\sigma_i\) fixes coordinate \(i\). The corrected list is

\[pd=1,\quad \sigma d=d,\quad qe=\sigma,\quad q\rho e=dp,
\quad q\rho^2e=1,\]
\[\sigma_1e=e,\quad q\sigma_3=\sigma q,\quad ged=ed\ (g\in\mathfrak S_3),
\quad pq\sigma_1=pq,\]

together with the group equations. The coordinate maps

\[p(x,y)=x,\quad d(x)=(x,x),\quad q(x,y,z)=(x,y),\quad e(x,y)=(y,x,x)\]

satisfy them. These maps and the permutations generate all maps between
sets of cardinalities at most three, by elementary inclusion and
retraction factorization. Hence the presented category \(Q\) maps
surjectively on hom-sets to \(\mathsf F_{\le3}^{\mathrm{op}}\).

We now prove injectivity by a terminating reduction and exact hom-set
bounds. First combine adjacent permutations. A path which rises from
\(E_1\) to \(E_2\) and immediately returns has a segment \(qge\). Every
\(g\in\mathfrak S_3\) is \(\rho^i h\) with
\(h\in\langle\sigma_1\rangle\). Since \(he=e\), the three relations
\(q\rho^ie\in\{\sigma,dp,1\}\) eliminate this segment. A path rising
from \(E_0\) to \(E_1\) and returning has a segment
\(p\sigma^\epsilon d=p d=1\). These replacements terminate: order words
first by the number of \(q,e\) transitions and then by the number of
\(p,d\) transitions. A first replacement lowers the first count; a
second lowers the second without increasing the first.

Every surviving path descends in height and then ascends, with at most
one valley. Paths through \(E_0\) are constant maps: on the output side
\(\sigma d=d\) and \(ged=ed\) erase all output permutations. On the
input side, there is one form from \(E_0\), at most two forms
\(p\sigma^\epsilon\) from \(E_1\), and at most three forms \(pqh\) from
\(E_2\), because \(pq\sigma_1=pq\). A path with valley \(E_1\) from
\(E_2\) to \(E_2\) is \(ge\sigma^\epsilon qh\). The stabilizer relation
\(\sigma_1e=e\) leaves three choices for the coset of \(g\); the
equation \(q\sigma_3=\sigma q\) moves the two-element left stabilizer of
\(h\) into the central \(\sigma^\epsilon\), leaving three choices for
its coset. Thus there are at most \(3\cdot2\cdot3=18\) forms of rank
two. The same calculation gives at most six rank-two maps each from
\(E_1\) to \(E_2\) or from \(E_2\) to \(E_1\). Finally there are at most
two invertible forms at \(E_1\) and six at \(E_2\).

Collecting these bounds, with source cardinality indexing rows and
target cardinality columns, gives

\[\begin{pmatrix}1&1&1\\2&4&8\\3&9&27\end{pmatrix}.\]

The entries are exactly
\(a^b=|\operatorname{Hom}_{\mathsf F^{\mathrm{op}}}(a,b)|\). Since the
functor was already surjective on each hom-set, each bound is attained
and it is bijective on each hom-set. It is the identity on objects,
hence an isomorphism of categories. The reduction proof supplies completeness, which checking the
relations in the target alone would not establish. \(\square\)

We transcribe the final relation as
\(p_0p_{12}\sigma_1=p_0p_{12}\). The subscript in the boxed list on
archive PDF page~5 (printed footer~4) is densely inked; the same
relation is written clearly with subscript~1 on PDF page~7 (printed
footer~6)~\cite{Groth}. Replacing~1 by~3 would assert that the first-
and second-coordinate projections agree, which is false already on
\((0,1,0)\). The alternative \(\sigma_3e=e\) is also false: the two
maps have coordinate values \((x,y,x)\) and \((y,x,x)\). The displayed
list instead has \(\sigma_1e=e\), which is true. Retaining both
incompatible alternatives therefore gives a false presentation.

\section{Controlled deployments and their exit categories}\label{sec:deployment}
We use the finite controlled deployment defined in Section~1, including its source, target, flag and collar data. The following examples separate the different control requirements.

\begin{proposition}[Three distinct control requirements]\label{prop:control}
Subset incidence alone does not imply the second equation in
\eqref{eq:control}. Neither a conical mapping-cylinder attachment nor
properness of its source implies a source fibration. Agreement at one
radius level does not imply preservation of the radius function.
\end{proposition}
\begin{proof}
For the first claim take all three object and arrow spaces to be
$\{0,1\}$, all target inclusions to be identities, $s_{01}=s_{12}=1$,
and $s_{02}$ the transposition. All subset conditions hold and
composition fails. A compact smooth version uses
$\Sigma_0=S^1$, $\Sigma_1=S^1\times[0,1]$ and
$\Sigma_2=S^1\times[0,1]^2$. The incoming faces are the two opposite
horizontal and two opposite vertical faces; use coordinate projections
for $s_{12},s_{01}$ and rotate $S^1$ through $\pi$ for $s_{02}$.
These are proper bundles on transverse faces, but the projections do
not commute.

For the second claim attach the mapping cylinder of the constant map
$p:S^1\to\R$, $p=0$, to a disc along $S^1\times\{1\}$.
The result is a sphere attached to a line at one point. It has cut-open
pieces $\R,D^2$ and incoming face $S^1$. The source is proper but the
path $t\mapsto t$ in $\R$ cannot lift starting at any point of $S^1$.
Thus the weaker mapping-cylinder data are insufficient. The example is
not locally conically trivial along the line.

For the third claim use coordinates $r>0,y\ge0$ and fix $\epsilon>0$.
The two retractions
\[
 \pi(r,y)=(r,0),\quad
 \pi'(r,y)=\left(r+(r-\epsilon)h(y),0\right),\qquad
 h(y)=-\frac{y}{4(1+y)}
\]
agree on $r=\epsilon$ and have the same inverse image of that level.
Indeed $r\pi'-\epsilon=(r-\epsilon)(1+h(y))$ and $1+h>3/4$.
Their first coordinates are positive; together with $y$, each is a
valid tubular coordinate after shrinking. Off that level $r\pi'\ne r$.
If one fixed retraction preserves every sufficiently small level, the
radius equation does follow, by taking the level through each point.
\end{proof}

\phantomsection\label{proof:projection-control}
With \eqref{eq:control}, the relation defined by the deployment is
transitive: successive arrows compose by the displayed source equation.
The lower interval of each point is linear even when the label poset
branches: two incomparable predecessor labels would place that point
in two incoming faces required to be disjoint. For comparable labels
$i<j<k$, the inverse-image identity places the predecessor of label
$j$ in $\Sigma_{ij}$, and projection control identifies its predecessor
of label $i$ with the given one. Each label has at most one predecessor
because targets are embeddings. Conversely suppose the relation is transitive and its pointwise lower
intervals have the order induced by the labels. For a point in
$\Sigma_{ik}\cap\Sigma_{jk}$ with $i<j$, its label-$i$ predecessor
must precede its label-$j$ predecessor. The definition of that arrow
places the latter in $\Sigma_{ij}$ and gives projection control.
This proves one subset inclusion in the inverse-image identity. For
the reverse inclusion, a point of $\Sigma_{jk}$ whose predecessor
lies in $\Sigma_{ij}$ gives a composable two-arrow chain; transitivity
places it in $\Sigma_{ik}$. Thus both the subset and projection
identities follow. These are exactly the required algebraic facts. Higher
associativity follows by composing these equations; it is not an extra choice.

\begin{proposition}[The categorical nerve]\label{prop:seg}
If all strict source maps of a finite graded topological category are
Hurewicz fibrations and its arrows within each label are identities,
its ordinary nerve is a complete Segal object in spaces.
\end{proposition}
\begin{proof}
The Segal maps identify the nerve with iterated ordinary pullbacks.
Each is a homotopy pullback because the source maps are fibrations,
and that property is preserved by base change and composition.
An equivalence maps to an identity in the label poset, so lies within
one label, where all arrows are identities. Hence the equivalence
subspace of degree one is precisely the object space. This is the
completeness condition. A Reedy fibrant model may still be needed:
the identity category on $S^1$ has constant nerve and matching map
$S^1\to S^1\times S^1$, which is not a fibration.
\end{proof}

\subsection{Simultaneous mapping cylinders}
For a flag $e=\{i_0<\cdots<i_n\}$ let $\Sigma_e$ be its iterated
arrow space. Write $e_{\le a}=\{i\in e:i\le a\}$.
For $d\subset e\setminus\{\max e\}$, write
$d\cup\{\max e\}=\{a_0<\cdots<a_p\}$. Define
\[
 C(d,e)=\operatorname{hocolim}
 (\Sigma_{e_{\le a_p}}\to\cdots\to\Sigma_{e_{\le a_0}})
\]
by the thick iterated mapping-cylinder model: its blocks are
$\Sigma_{e_{\le a_k}}\times\Delta^k$, with a face deleting the final
vertex attached by the corresponding source map. Other faces remain
boundary faces. There are inclusions when $d$ is enlarged, restriction
maps to the selected subflag $d\cup\{\max e\}$, and, crucially,
inclusions of each initial-prefix cylinder as a face of the larger
one. All relations are those induced by flag inclusion and source
composition.

\begin{proposition}[Reconstruction by cylinders]\label{prop:cylinders}
The colimit of this complete cylinder diagram is the realization of
the nerve. If a finite geometric deployment is built recursively by
compatible conical collars, this realization is homeomorphic to the
original space, relative to every retained outer face.
\end{proposition}
\begin{proof}
A map from the nerve realization to a space $T$ is a family
$\Sigma_e\times\Delta^{|e|-1}\to T$ compatible with all faces.
Faces deleting the final vertex occur inside a cylinder and are
identified with their standalone prefix blocks by the prefix maps.
Every other face contains the final vertex and is accounted for by
enlargement and selected-subflag restriction. Thus cocones out of the
cylinder diagram are exactly these compatible families, naturally in
$T$. The quotient topologies therefore agree. If prefix maps were
omitted, even the constant diagram on $0<1$ would retain an extra
isolated vertex.

For the geometric claim remove a minimal stratum $A=\Sigma_i$.
Let $W$ be its cut complement and $H$ its retained link boundary.
Recursive compatibility means
$X=W\cup_H C(H\to A)$, and that $W,H$ inherit the same kind of data,
including all earlier outer faces. Let $D$ be the category on the
remaining labels and let $M$ have flag spaces
$\Sigma_{i,j_0,\ldots,j_n}$. Every simplex either lies in $D$ or starts
with $i$. In the latter simplex put $t_i=1-r$ and $t_j=ru_j$.
These coordinates give mutually inverse continuous maps
\[
 |C|\cong |D|\cup_{|M|}C(|M|\to A).
\]
Induction identifies this pushout with the prescribed geometric one,
including its boundary maps. Fibrancy is unnecessary for this
homeomorphism; it enters the exit-category statement next.
\end{proof}

\begin{proof}[Proof of Theorem~\ref{thm:exit-deployment}]
We first record the categorical gluing formula. For a category
$\mathcal C$ conservative over a finite poset and a minimal label $i$,
let $A=\mathcal C_i$, let $\mathcal D$ be the remaining full category,
and let $\mathcal M$ be its crossing-arrow category, with source
$q:\mathcal M\to A$ and target $p:\mathcal M\to\mathcal D$.
Then
\begin{equation}\label{eq:collage}
 \mathcal C\simeq A\coprod_{\mathcal M\times\{0\}}
 (\mathcal M\times[1])
 \coprod_{\mathcal M\times\{1\}}\mathcal D.
\end{equation}
Indeed a functor from the right side to $K$ is a pair of functors on
$A,\mathcal D$ and a natural transformation $F_Aq\to F_Dp$.
Every composable string crosses from $A$ to $\mathcal D$ at most once.
Its initial part transports the crossing arrow by equivalences in
$A$ and its remaining part composes in $\mathcal D$. Naturality on
$\mathcal M$ supplies exactly those compositions and all their
simplicial coherences. The Segal maps have contractible inverse choices,
so this construction and restriction are inverse on mapping spaces.
This proves \eqref{eq:collage} by Yoneda.

In a deployment the category $M$ of flags beginning with $i$ presents
$\mathcal M$. Its degree-$n$ comparison is
$M_n\simeq M_0\times^h_{D_0}D_n$: the strict flag identity
\eqref{eq:control} gives this pullback, and source fibrations make it
a homotopy pullback. These are the left-fibration comparisons for
$p$, while the projection equation defines $q$. Thus this is the
actual crossing-arrow category, with source-side paths included.

Induct on the labels. For one label the statement is the fundamental
infinity-groupoid of that stratum. At a minimal label use the previous
paragraph and boundary-compatible induction to identify
$\mathcal D\simeq\operatorname{Exit}(W)$ and
$\mathcal M\simeq\operatorname{Exit}(H)$.
Retained outer faces can be removed coherently: in each collar
coordinate $u\in[0,1)$ replace $u$ by $(1-t)u+tf(u)$, where $f(0)>0$,
$f(u)\ge u$ and $f(u)=u$ near the outer collar edge. The finitely many
pushes commute. They preserve the remaining strata and commute with
source maps; hence they give exit equivalences. Old faces of $H$ are
the inverse images of old faces of $A$, because the removed labels form
a lower ideal. Consequently restricting $H\to A$ to these interiors
is a base change and remains proper. Its comparison to the complementary
open piece is a positive-radius collar section, with the stated collar
homotopy; no deleted-face inclusion is falsely required to commute
strictly with an inward push.

The proper fiberwise-cone formula of Volpe
\cite[Lemma~2.13]{Volpe} identifies the exit category of the cone
collar with $A\coprod_{\mathcal M\times\{0\}}(\mathcal M\times[1])$.
Exit excision for the open collar and its enlarged complement
\cite[Corollary~2.17]{Volpe} glues this to $\operatorname{Exit}(W)$
along $\operatorname{Exit}(H)$. This is precisely
\eqref{eq:collage}. All equivalences arise from these same source,
incidence and collar maps, so their naturality supplies the
boundary-compatible induction hypothesis at the next step.
\end{proof}

\subsection{Resolution and the exact corner condition}
The usual smooth resolution theorem concerns manifolds with iterated
boundary fibrations, not bare incidence spaces. The incoming collective
hypersurfaces must cover the boundary, their source bundles must have
compact manifold-with-corners fibers, horizontal faces must be pullbacks
of base faces, vertical faces must be exactly the boundaries of fibers,
and compatible radial collar germs must be retained. On a maximal
piece $\Sigma_m$ the face labelled $i$ is $\Sigma_{im}\to\Sigma_i$;
\eqref{eq:control} is exactly its commuting boundary-fibration diagram.
Iterating it recovers every higher face. Conversely an iterated boundary
fibration supplies these data. Collective hypersurfaces retain a common
base with possibly disconnected links; separating them would normalize
the singularity.

Albin's resolution/blowdown theorem \cite[Theorem~6.3, pp.~48--49]{Albin}
identifies the compact smooth objects so specified with Thom--Mather
spaces. In a collar $(h,r_i)$, blowdown defines
$\pi_i(h,r_i)=s_i(h)$ and $\rho_i(h,r_i)=r_i$.
Compatibility gives both $\pi_k\pi_i=\pi_k$ and
$\rho_k\pi_i=\rho_k$ for $k<i$; the bundle charts give the required
stratumwise submersions. Conversely cutting controlled tubes at positive
radius recovers faces, sources and commuting radial collars. This
explains the hypotheses of the known theorem and the information that
one isolated level forgets. Without fixed collars blowdown is defined
up to isomorphism, not equality of radius functions. For noncompact
finite-depth smooth atlas spaces the corresponding object theorem is
\cite[Theorems~2.10 and~2.13]{ABP}; its category of maps has additional
liftability conditions.

\begin{proposition}[Corners with specified faces]\label{prop:corners}
Recursive compatible cuts give transverse collars and Boolean flag
intervals. The output is a manifold with its specified corners in TOP,
PL or smooth precisely when its open face bases are manifolds in that
category and the collar transitions have that regularity, with the
usual Hausdorff and second-countability conditions.
\end{proposition}
\begin{proof}
A new cut is a product with every retained collar coordinate.
Simultaneously minimal cuts are disjoint; by recursion all later cuts
preserve the earlier products. Incomparable incoming labels cannot meet:
a common link point would lie in both original strata outside their
already removed lower intersection. Thus nonempty face intersections
are flags. If $e\subset d$ have the same maximal label, intermediate
flags are exactly $e\cup A$ for $A\subset d\setminus e$; this is a
Boolean interval. At $x$, intersect exactly the boundary germs through
$x$, remove smaller faces, and call the resulting open face $F_x^\circ$.
The collar germ is
\[
 (X,x)\cong(F_x^\circ,x)\times([0,\infty)^{k(x)},0).
\]
A manifold chart on $F_x^\circ$ makes an orthant chart. Conversely a
corner chart preserving the specified faces restricts to a chart on
$F_x^\circ$. Dimension must satisfy $\dim F_x^\circ+k(x)=n$ when a
single dimension is prescribed. This proves the equivalence.

For a counterexample to omitting the base condition, take the Cantor
set $K$ and $K\times[0,\infty)$. Its evident collar and one-step
Boolean interval satisfy the combinatorial requirements, while no
neighborhood is locally connected. Even a one-label Cantor deployment
shows that bare incidence data do not characterize Thom--Mather spaces.
\end{proof}

The same local proof applies to locally finite finite-depth recursive
data when the cut charts and retained collar family are locally finite.
Every point then has a neighborhood meeting only finitely many active
cuts and collar coordinates; the preceding finite induction supplies
its chart. On overlapping neighborhoods these charts use restrictions
of the same specified collars, so the local conclusions agree. The
usual Hausdorff and second-countability requirements remain separate.
An arbitrary accumulating family of faces is not covered by this
argument and need not give a finite-dimensional orthant chart.

Ordinary manifold strata still do not suffice for an arbitrary TOP
choice of cone data. With $F=\Sigma H$ as in
Example~\ref{ex:cannon}, the open cone $C(F)$ is $\R^5$:
$\Sigma F\cong S^5$ by double suspension, and deleting one suspension
point identifies its complement with $C(F)$. Its apex is a point
stratum, and its complement $F\times(0,\infty)$ is a five-manifold.
Cutting this chosen cone nevertheless gives $F\times[0,\infty)$,
which is not a manifold with boundary by the same endpoint local-homology
argument in that example. The condition on the resolved face bases
therefore has mathematical content even for these manifold strata.

For conically smooth spaces with manifold strata, the unzip theorem
\cite[Theorem~7.3.8]{AFT} supplies the required resolved links and
corner charts. Locally its induction replaces $B\times C(L)$ by
$B\times[0,\infty)\times\widetilde L$. For normal crossings, branch
labels must be retained: monodromy may permute branches, so an incidence
category, rather than its orbit poset, is the correct global object.

\subsection{The moduli qualification}\label{sec:orbifold-exit}
We specify the category in which the moduli comparison holds. In the
smooth controlled case cut lower strata before higher ones, retaining
each new boundary. The cut-open stratum is $\Sigma_i$; its incoming
face in $\Sigma_j$ is $\Sigma_{ij}$. For a flag
$d=(i_0<\cdots<i_n)$, its space is
\[
 \Sigma_d=\bigcap_{a<n}\Sigma_{i_ai_n}.
\]
Deleting a nonfinal entry is an inclusion. Deleting the last entry is
the corresponding source projection. The two identities
\eqref{eq:control} make their compositions agree and their Segal
squares cartesian. Proper source bundles and compatible transverse
collars supply the hypotheses of Theorem~\ref{thm:exit-deployment}.
Cuts and face intersections inherit precisely these same bundle and
collar structures by the resolution theorem just stated. Thus this
construction has the recursive closure required there; it is not an
inference from arbitrary mapping-cone attachments.

For stable curves, take the finite incidence category of stable graphs
of type $(g,n)$, retaining individually labelled legs and vertex genera.
An arrow is a specified edge contraction together with an isomorphism
from the contracted graph to its target; all graph isomorphisms are
retained. Composites contract the union of the specified edge sets.
A loop contraction raises the vertex genus, and an edge between distinct
vertices merges them and adds their genera. The corresponding open
stratum is the orbifold
\[
 \mathcal S_\Gamma=
 \left[\prod_{v\in V(\Gamma)}\mathcal M_{g(v),n(v)}
       /\operatorname{Aut}(\Gamma)\right].
\]
The factors themselves retain their curve automorphisms. A graph
orbit poset would forget both the contraction incidences and isotropy.

A stable-curve universal deformation chart has smoothing coordinates
$q_e$ for its actual nodes, with a finite automorphism action
\cite[Theorems~8.11, 9.13; Proposition~9.16]{Looijenga}.
The replacements $q_e=r_e e^{i\theta_e}$ give an orbifold with corners,
and a manifold with corners on an actual free level cover. To obtain
a face for every strict stratum incidence, further blow up its corner
faces in increasing dimension. In normal coordinates this replaces
the face lattice by its flag lattice. For an $r$-node degeneration
the original coordinate-stratified normal link is $S^{2r-1}$ and the
full deployment uses its recursively resolved link. If $r=1$ it is a
circle. If $r>1$, radial ratios and their faces are also retained;
the torus part alone is not the complete link.

Choose the compatible boundary fibrations and collars on this orbifold
resolution, then pull them back to the charts. Such choices can be made
equivariantly: the normal-crossings blowups preserve incidences, finite
group averaging gives invariant metrics and connections, and the
iterated collar construction preserves the previously chosen boundary
fibrations. At an intersection its horizontal collar is pulled back
from the base and its vertical collars are constructed in the fibers;
this induction makes them commute. Orbifold partitions of unity make
the choices global. Changes of charts preserve these pulled-back data.
Independent arbitrary collar choices are not asserted to agree.

For every chain of specified contractions take its flag-face space and
the groupoid of chart identifications and graph isomorphisms. In degree
$k$ take the disjoint union over contraction chains of length $k$ with
these groupoids, and then their classifying spaces. Source maps are
representable proper locally trivial bundles. The flag identities
therefore identify the Segal comparisons with homotopy fiber products.
Crucially, isomorphisms act at every vertex of a chain and compatibly
on every arrow. For a single object with automorphism group $G$ the
degree-$k$ groupoid is
\[
 G^k\mathbin{//}G^{k+1}.
\]
Changing vertex trivializations acts transitively, with diagonal
stabilizer $G$, so its classifying space is $BG$ in every degree.
In general the groupoid of equivalence arrows is equivalent by source
evaluation to the object groupoid: its fiber over an object is the
contractible groupoid of isomorphisms out of it. This is the completeness
condition. Equivalently take the complete-Segal localization of the raw
incidence/action-category nerve. Replacing only object sets by
classifying spaces, or deleting automorphisms, would not give this
construction.

Here is the atlas descent argument, including its independence of
choices. Define $\operatorname{Exit}_{\rm orb}$ as the colimit in
infinity-categories of the exit categories of a simplicial orbifold
atlas. A surjective stratum-preserving local diffeomorphism $V\to U$
has local stratum-preserving sections. On an open set with a section,
inserting it gives an extra degeneracy in the augmented Cech diagram;
applying exit paths leaves a split augmented simplicial diagram, whose
colimit is the exit category of that open. Exit open-cover excision
\cite[Theorem~A.7.1]{Lurie} and interchange of colimits prove the same
statement over $U$. Two atlases have a common refinement by their fiber
product. Applying this argument in both directions of the resulting
bisimplicial diagram identifies their colimits. On a finite quotient
chart $[U/G]$ the atlas nerve is $U\times G^\bullet$, so the result is
$\operatorname{Exit}(U)\mathbin{//}G$, including ineffective isotropy.

The comparison in Theorem~\ref{thm:exit-deployment} is natural for
the actual source, boundary and chart-isomorphism maps. It gives an
equivalence on every chart and every atlas intersection. Taking the
preceding colimit therefore proves
\[
 \widehat{N\boldsymbol\Sigma}\simeq
 \operatorname{Exit}_{\rm orb}(\overline{\mathcal M}_{g,n}),
\]
where the hat denotes complete-Segal localization of the full chain
construction. This is descent of categorical equivalences, not a
deduction from equality of realizations. Groupoid completion gives the
orbifold homotopy type. An ineffective stabilizer still contributes
its classifying space; taking its effective quotient changes this
answer.

On a free level cover the construction gives the ordinary space and
its exit category. On coarse spaces finite orbit quotients commute
with the explicit topological cylinder colimits, so the realization
still recovers the coarse space. The ordinary exit theorem needs
additional care: coarse quotients can destroy the source-bundle
property and the boundary-only strata can fail to be manifolds, as
Example~\ref{ex:moduli} proves. One can instead refine the finite
quotient charts by orbit type and their boundary incidences. Averaged
metrics give invariant normal slices; the linear slice quotients are
cones on their sphere quotients. Induction on slice depth gives the
conical local models for this refinement, and the corresponding
resolved incidence diagram satisfies the qualified ordinary theorem.
Specified graph contractions alone do not make that refinement.

In $\overline{\mathcal M}_{0,4}\cong\mathbb P^1$, the cut-open pair
of pants, its three boundary circles and the three point bases give
an explicit example. Each source sends a boundary circle to its point;
the target is the boundary inclusion. Attaching the three mapping
cylinders fills the holes by discs and reconstructs $S^2$.

The empty deployment gives the empty category; one stratum gives its
fundamental infinity-groupoid. Empty labels may be removed and
disconnected strata retained as disjoint unions. A strictly increasing
grading may merge incomparable components, provided their actual
incidences and disjoint incoming branches remain in the diagram: no
strict arrow then becomes a same-label arrow. No dense-order or
divisibility assumption is used. Finally these homotopy and exit
comparisons do not give arbitrary-sheaf descent from only the named
pair and triple stones. For that purpose Theorem~\ref{thm:tubes}
retains the genuine neighborhood intersections and their descent
maps; applying homotopy type forgets different information.

\begin{counterexample}[An interior coarse quotient]\label{ex:moduli}
The coarse space of genus-four curves need not be a topological
manifold, even away from the boundary. Consequently boundary blowups
and specified contractions do not remove every obstruction.
\end{counterexample}
\begin{proof}
Let $C$ be the smooth completion of
\[
 y^2=x(x-1)(x-3)(x-7)(x-12)(x-20)(x-31)(x-45)(x-64).
\]
The ten branch points, including infinity, give genus four.
A fractional linear transformation preserving them is determined by
the images of $(\infty,0,1)$. The exact finite calculation described
below checks all $10\cdot9\cdot8=720$ possibilities; only the identity
preserves the set. This also exhausts complex transformations, since
three rational source and target points determine a rational map.
Two inequivalent degree-two maps $C\to\mathbb P^1$ would give a
birational curve of bidegree $(2,2)$ in
$\mathbb P^1\times\mathbb P^1$, of arithmetic genus one. If the joint
map has degree two instead, its image has bidegree $(1,1)$, making the
maps equivalent. Thus the hyperelliptic quotient is unique, and
$\operatorname{Aut}(C)=\{1,\iota\}$, $\iota(x,y)=(x,-y)$.

The invariant quadratic differentials are
$x^k(dx/y)^2$, $0\le k\le6$; the anti-invariant ones are
$x^k dx^2/y$, $0\le k\le1$. At a finite branch point use
$x-a=t^2$; at infinity $x=t^{-2}$ and $y=t^{-9}$ times a unit.
Their respective orders are $12-2k$ and $3-2k$, proving regularity and
the bounds. They are independent and their dimensions sum to
$9=3g-3$. Serre duality therefore gives the deformation quotient
$\C^7\times(\C^2/\{\pm1\})$. Finite holomorphic actions are locally
linearized by averaging a coordinate map after applying inverse linear
actions; its derivative is the identity. This proves that the tangent
quotient is the actual analytic germ.

The link is $S^{13}*\mathbb{RP}^3=\Sigma^{14}\mathbb{RP}^3$.
The cellular boundary maps of $\mathbb{RP}^3$ are $d_2=2$ and
$d_1=d_3=0$, so local homology has $\Z/2$ in degree 16 in addition
to $\Z$ in degree 18. This is impossible at an interior or boundary
point of an 18-manifold. The curve is smooth, so modifications of the
Deligne--Mumford boundary do not affect this germ.

For completeness the finite test uses homogeneous vectors $a,b,c$ for
three distinct branch points. The matrix with columns
$\det(c,b)a$ and $\det(a,c)b$ sends $(\infty,0,1)$ to $(a,b,c)$.
Normalize each image pair by its greatest common divisor and sign, and
compare the resulting ten-element set with the original set. This
integer-arithmetic recipe checks all triples without approximation or
any search over unbounded coefficients.
\end{proof}


\section{Pro-neighborhoods, homotopy links, and cohomology}\label{c.4.-pro-neighbourhoods-holinks-and-cohomology}

\subsection{The germ comparisons}\label{c4.01-proved-for-finite-metric-hss-the-germ-comparisons}

Use a finite stratification satisfying the frontier condition, with closed lower unions and locally closed strata \(S_i=X_i^*\). Assume the metric space is homotopically stratified in Quinn's sense. For every incident pair this requires forward tameness (a neighborhood deformation fixing the lower stratum and staying in the upper stratum until its final time) and the Hurewicz fibration property for initial evaluation of the homotopy link. Whitney and Thom--Mather examples fall in this scope. Put \(O_i=X\setminus\bigcup_{k<i}S_k\), so \(S_i\) is closed and pure in the open subspace \(O_i\). Let \(V_i\) be the prospace of open neighborhoods of \(S_i\) in \(O_i\), ordered by shrinking. These neighborhoods are also open in \(X\). Finite intersections define \(V_d\).

All prospaces below are objects of $\operatorname{Pro}(\mathcal S)$,
where $\mathcal S$ is the infinity-category of spaces. Every union and
intersection uses one tuple of stratum neighborhoods, ordered by
simultaneous shrinking. Comparisons over a constant space are taken in
the corresponding infinity-categorical slice.

Quinn, \cite[Proposition~3.2, p.~469]{Quinn}, gives one fixed nearly strict and nearly stratum-preserving neighborhood deformation \(h_i:N_i\times[0,1]\to O_i\), with \(h_i(x,0)=x\), \(h_i(x,1)\in S_i\), and \(h_i(s,t)=s\) for \(s\in S_i\). Until the final time it stays in the starting stratum. Reversing Quinn's interval convention gives the notation used here. The theorem applies to \(O_i\), whose induced stratification is again HSS. More generally, choose a linear extension of the finite label poset. Its prefixes are lower ideals, hence give closed skeleta. On any union of retained strata, their intersections remain closed relative skeleta, and each retained two-stratum subspace is unchanged, so its pairwise tameness and holink fibration remain valid. For a nonincident pair the two strata are clopen in their union by the frontier condition, so the lower one is trivially tame and the holink is empty. Thus the induced union is a finite metric HSS in precisely Quinn's sense. No assumption that an arbitrary restricted deformation preserves a newly chosen neighborhood is used.

\textbf{Single-stratum germs.} Given an open \(U\supset S_i\), compactness of the time interval and continuity of \(h_i\) imply that \[
 W_U=\{x\in N_i:h_i(x,[0,1])\subset U\}
\] is open and contains \(S_i\). The fixed endpoint retraction \(r_i=h_i(-,1)\), restricted to \(W_U\), and the restricted homotopy prove that \(S_i\to V_i\) is a pro-homotopy equivalence. The same endpoint map is used at all levels, so these are compatible pro maps, not unrelated levelwise equivalences. If \(Z\) is a locally closed union of strata containing \(S_i\), the homotopy on \(W_U\cap Z\) stays in \(Z\), including its endpoint, proving the corresponding two germ comparisons.

\textbf{Flag germs.} Let \(d=(i_0<\cdots<i_n=j)\), and fix target neighborhoods \(U_{i_k}\) for every flag entry. First shrink \(N_j\) so that its whole fixed homotopy stays in \(U_j\). For \(k<n\), metric normality and closedness of \(S_{i_k}\) in \(O_{i_k}\) give open neighborhoods \(U'_{i_k}\) with \(\overline{U'_{i_k}}^{\,O_{i_k}}\subset U_{i_k}\). Set \(A=\bigcap_{k<n}U_{i_k}\) and \(A'=\bigcap_{k<n}U'_{i_k}\). Because \(O_j\subset O_{i_k}\), one has \(\overline{A'\cap O_j}^{\,O_j}\subset A\). Let \[
 W=\{x\in N_j:h_j(x,[0,1])\subset A\}.
\] It is open and contains \(S_j\cap A\). Therefore \(F=\overline{A'\cap O_j}^{\,O_j}\setminus W\) is closed in \(O_j\) and disjoint from \(S_j\). Choose \(U'_j=N_j\setminus F\). Then \(A'\cap U'_j\subset W\), so its whole homotopy lies in \(A\cap U_j\). These restrictions of the one fixed \(h_j\) define a pro-homotopy retraction \[
 V_d\longrightarrow V_d\cap S_j.
\] It is the identity on \(S_j\), giving both inverse homotopies. If \(Z\) is a locally closed union of strata containing \(S_j\), the same homotopy restricted to \(Z\) stays there. The triangle through \(V_d\cap S_j\) proves the comparison whenever \(S_j\subset Z\).

These are homotopy equivalences in prospaces, represented by maps and homotopies after cofinal shrinking; in particular they are isomorphisms in \(\operatorname{Pro}(\mathrm{HoTop})\). No assertion is made that taking a limit in ordinary spaces or topoi preserves them. A coherent pro-space formulation uses the same fixed homotopies, so all refinements come from restriction.

For precision, index a refinement argument by pairs $(\mathbf V,\mathbf U)$
with $\mathbf V\subset\mathbf U$ for which the fixed homotopy sends the
source intersection into the target intersection. This poset and the
comma categories for either projection are nonempty and cofiltered:
intersect finitely many target constraints, then apply the preceding
construction inside all prescribed source neighborhoods. Both projections
are therefore coinitial. Endpoint restrictions and the one fixed homotopy
are strictly natural on this poset and give inverse maps in
$\operatorname{Pro}(\mathcal S)$.

\paragraph{Bands.} Let \(J,L\subset I\) be convex bands satisfying \(J\prec L\): every member of \(J\) is below some member of \(L\), and every member of \(L\) is above some member of \(J\). The interval hull \[
 I'=\{i:\exists j\in J,\ell\in L,\ j\le i\le\ell\}
\] is convex; \(J\) is a lower set and \(L\) an upper set in \(I'\). For example, if \(j_0\le i\le j\) and \(j_0,j\in J\), convexity puts \(i\) in \(J\); the upper statement is dual. The same formula gives the unique smallest interval hull.

For a finite metric stratification choose a cofinal system of neighborhoods with \(U_i\cap U_k=\varnothing\) when \(i,k\) are incomparable. Indeed the two strata are disjoint closed subsets after removing both lower frontiers; normality separates them, and finitely many such shrinkings can be intersected. Thus any nonempty finite neighborhood intersection has flag labels. The space \(V_J\cap V_L\) is covered by the finitely many \(V_j\cap V_\ell\), \(j\in J,\ell\in L\). If \(j\le\ell\), the maximum label is in \(L\). If \(\ell<j\), then convexity and domination imply both labels belong to \(J\cap L\); again the maximum is in \(L\). The same remains true for every finite intersection of these pieces.

For a convex band $L$, both $\downarrow L$ and
$(\downarrow L)\setminus L$ are lower ideals. Hence
$X_L^*=X_{\downarrow L}^*\setminus X_{(\downarrow L)\setminus L}^*$
is locally closed, as required for the flag restriction comparison.

Apply the flag result to each such intersection, with \(Z=X_L^*\), which contains its largest stratum. It gives a pro-homotopy equivalence on every intersection after restricting to \(X_L^*\). Finite open-cover homotopy descent then gives \[
 V_J^L=V_J\cap X_L^*\ \longrightarrow\ V_J\cap V_L
\] as a pro-homotopy equivalence. To justify descent concretely, use the finite poset of nonempty subsets of the finite cover-piece label set, with opposite inclusion for its intersection diagram. Its nerve has finitely many nondegenerate simplices. A partition of unity gives a homotopy inverse from the union to this finite homotopy colimit. This is the numerable-cover comparison of \cite[Proposition~4G.2, p.~459]{HatcherCover}; metric spaces are paracompact. The natural restriction map of intersection diagrams is an equivalence in prospaces at every object by the fixed-homotopy proof above; its inverses need not be chosen as separate natural maps. Finite homotopy colimits preserve this equivalence and commute with a synchronized cofiltered presentation: on mapping spaces this follows from filtered colimits commuting with finite limits \cite[Theorem~9.1.5.1]{KerodonFiltered}. More explicitly, if $Q$ is that finite intersection poset, $F_q(\mathbf U)$ its diagram, and $K$ a constant target space, then
\[
\begin{aligned}
 \operatorname{Map}_{\operatorname{Pro}(\mathcal S)}
  (\{\operatorname*{colim}_{q\in Q}F_q(\mathbf U)\}_{\mathbf U},K)
 &\simeq\operatorname*{colim}_{\mathbf U}
    \lim_{q\in Q^{\rm op}}
       \operatorname{Map}_{\mathcal S}(F_q(\mathbf U),K)\\
 &\simeq\lim_{q\in Q^{\rm op}}
    \operatorname*{colim}_{\mathbf U}
       \operatorname{Map}_{\mathcal S}(F_q(\mathbf U),K).
\end{aligned}
\]
The limits are finite because the nerve of $Q$ has only finitely many nondegenerate simplices. Constant target tests detect pro-equivalences by the pro-Yoneda description \cite[Section~3.2]{BHH}. Only finitely many source refinements are required. This argument does not truncate the unbounded ordinary \v{C}ech nerve, and would require a separate justification for an infinite index family.

This argument does not require a strictly commuting system of Mather retractions. 

\subsection{The relative holink comparison}\label{c4.02-proved-the-relative-holink-comparison}

For possibly overlapping dominated convex bands put
\(Y=X_J^*\), \(L_0=X_L^*\), and \(Z=Y\cup L_0\). Define the relative path space
\[
 P_{J,L}=\{\gamma\in C([0,1],Z):\gamma(0)\in Y,
                    \ \gamma(t)\in L_0\text{ for every }t>0\}.
\]
It contains \emph{all} such paths: paths starting in \(Y\cap L_0\) need not be constant. In particular it contains all constant paths on that intersection. Give it the compact-open topology. The band order lemma above makes \(J\) a lower ideal in \(I'\), hence in the retained label set \(J\cup L\). The finite frontier condition therefore makes \(Y\) closed in \(Z\). It is a union of whole strata, hence pure. The preceding induced-filtration argument proves that \(Z\) is a finite metric HSS. Consequently \cite[Proposition 3.2]{Quinn} supplies a fixed neighborhood \(N\) of \(Y\) and a nearly strict and nearly stratum-preserving deformation
\(h:N\times[0,1]\to Z\), with \(h_0=\mathrm{id}\), \(h_1=r:N\to Y\), and \(h_t|_Y=\mathrm{id}\). For \(t<1\), a point outside \(Y\) remains in its original stratum.

For each open neighborhood \(U\) of \(Y\) in \(Z\), let \(P_U\) be the paths in \(P_{J,L}\) whose whole image lies in \(U\). There is a natural zigzag
\[
 \operatorname{const} P_{J,L}
 \ \longleftarrow\ \{P_U\}_{U\supset Y}
 \ \xrightarrow{e_1}\ \{U\cap L_0\}_{U\supset Y}=V_J^L.
\]
The equality on the right means equality of pro-germs: restrictions of the unions of the finitely many individual stratum neighborhoods form a cofinal system of neighborhoods of \(Y\) in \(Z\). Indeed each such union is open and contains \(Y\); conversely an open neighborhood of \(Y\) in \(Z\) is the restriction of an ambient open set containing \(Y\), and intersecting that ambient open set with the prescribed open complements of the lower frontiers gives the required individual neighborhoods.

\textbf{The left arrow.} The path space is metrizable, using the uniform metric of any bounded compatible metric on \(Z\), and hence is paracompact. The open subsets
\[
 C_n=\{\gamma:\gamma([0,1/n])\subset U\},\qquad n\ge1,
\]
cover \(P_{J,L}\). Choose a locally finite partition of unity subordinate to a refinement of this cover; for each member of the refinement retain its integer \(n\). If its partition functions are \(\lambda_a\) and integers are \(n_a\), put
\(\epsilon(\gamma)=\sum_a\lambda_a(\gamma)/n_a\).
This is positive and continuous. Every active \(1/n_a\) is an admissible initial interval for \(\gamma\); since their finite maximum is also admissible, \(\gamma([0,\epsilon(\gamma)])\subset U\). Thus
\[
 S\gamma(u)=\gamma(\epsilon(\gamma)u)
\]
is a continuous map \(P_{J,L}\to P_U\). The homotopy
\[
 K_s\gamma(u)=\gamma\bigl((1-s+s\epsilon(\gamma))u\bigr)
\]
joins the identity to inclusion followed by \(S\), remains in \(P_{J,L}\), and restricts to a homotopy in \(P_U\) when \(\gamma\in P_U\). It fixes the initial point. Consequently every inclusion \(P_U\hookrightarrow P_{J,L}\) is a homotopy equivalence. These inclusions are a natural transformation, so the left arrow is a pro-homotopy equivalence. The endpoints are not fixed by this shortening homotopy; their coherent comparison in \(L_0\) is the explicitly displayed terminal segment of \(\gamma\).

Each inclusion $P_U\to P_{J,L}$ is a map over $L_0$ whose underlying
map is an equivalence. It is therefore an equivalence in
$\mathcal S_{/L_0}$. These inclusions are natural in $U$, so they form
an objectwise equivalence of diagrams and induce an equivalence over
$\operatorname{const}(L_0)$ in prospaces. No compatible choice of the
shortening inverses is needed. The structural endpoint homotopy for a
chosen inverse is
$s\mapsto\gamma(\epsilon(\gamma)+(1-\epsilon(\gamma))s)$.

\textbf{The right arrow.} For a target neighborhood \(U\), put
\[
 W_U=\{x\in N:h(x,[0,1])\subset U\}.
\]
Compactness of \([0,1]\) makes \(W_U\) an open neighborhood of \(Y\).

The assignment $U\mapsto W_U$ is monotone, satisfies $W_U\subset U$,
and preserves finite intersections; it is consequently a coinitial
refinement of the neighborhood system. The maps $s_U$ and the homotopies
below are strictly natural after this reindexing. Define
\[
 s_U:W_U\cap L_0\longrightarrow P_U,
 \qquad s_U(x)(u)=h(x,1-u).
\]
For \(x\in L_0\setminus Y\), the positive-time points of this reversed path remain in the stratum of \(x\), hence in \(L_0\). For \(x\in Y\cap L_0\) the path is constant. Thus it belongs to the relative path space in either case, and \(e_1s_U=\mathrm{id}\). All the maps \(s_U\) are restrictions of the same formula using the fixed \(h\); for finitely many targets a common smaller neighborhood gives literal agreement. They therefore define one pro map.

For the other composite, take \(\gamma\in P_{W_U}\) and define, for \(t>0\),
\[
 H_t(\gamma)(u)=
 \begin{cases}
 h\bigl(\gamma(t),1-u/t\bigr),&0\le u\le t,\\
 \gamma(u),&t\le u\le1,
 \end{cases}
 \qquad H_0(\gamma)=\gamma.
\]
The formulas agree at \(u=t\). Every positive-time point lies in \(L_0\), including when \(\gamma(t)\in Y\cap L_0\), because then \(h\) fixes that point. Every path lies in \(U\), and its endpoint is always \(\gamma(1)\). At \(t=1\) the path is \(s_U(e_1\gamma)\). Continuity away from \((t,u)=(0,0)\) is immediate. At that corner, \(\gamma(t)\to\gamma(0)\in Y\), and continuity of \(h\), compactness of its time interval, and \(h(\gamma(0),v)=\gamma(0)\) uniformly in \(v\) prove continuity. The same argument with \(\gamma\) varying gives continuity of the adjoint map on \(P_{W_U}\times[0,1]^2\), hence continuity into the compact-open path space. This proves the second inverse homotopy after shrinking. All refinements restrict this one fixed formula. The right arrow is consequently a pro-homotopy equivalence, in fact over \(L_0\).

The zigzag thus gives an equivalence in $\operatorname{Pro}(\mathcal S)_{/\operatorname{const}(L_0)}$. Its arrows literally commute with the endpoint maps to \(L_0\); the start map on the middle term is evaluation at zero. Under the inverse supplied by \(h\), the induced map to \(Y\) is \(r\), and the displayed \(H_t\) supplies its coherent comparison with the initial point of any representative path.

When \(J\cap L=\varnothing\), this relative path space is the conventional strict holink of \((Z,Y)\), recovering Quinn's usual formulation. When \(J=L\ne\varnothing\), that conventional strict holink is empty, whereas \(V_J^J=X_J^*\). The relative path space instead becomes \(C([0,1],X_J^*)\), which contracts by endpoint evaluation to the constant paths. Thus the literal strict-holink claim needs disjointness; the all-paths relative formulation above is valid with overlap.

The domination conditions force $J=L=\varnothing$ if either label set
is empty; all spaces and prospaces in that case are empty. An empty
stratum has the empty neighborhood as a cofinal choice. If a nonempty
label band has empty geometric union, the same empty-space argument
applies to the relative germ and path space.

\subsection{The finite cover spectral sequence}\label{c4.03-proved-the-finite-cover-spectral-sequence}

Fix finite dominated bands, the interval hull \(I'\), and \[
 D=\{d\subset J\cup L:d\text{ a nonempty flag},\
             \min d\in J,\ \max d\in L\}.
\] Give \(D\) arrows \(d\to e\) when \(d\subset e\). Thus \(V_e\subset V_d\), and cohomological restriction is covariant in \(D\). For an abelian sheaf \(F\) on the neighborhood germ \(V_J^L\), interpret its cohomology as the filtered colimit of cohomology of representatives. There is a first-quadrant, strongly convergent spectral sequence \[
 E_2^{p,q}=R^p\!\lim_{d\in D}H^q(V_d^{X'},F)
 \ \Longrightarrow\ H^{p+q}(V_J^L,F).
\] Here the derived limit is taken in abelian groups over the displayed flag category.

First choose cofinal representatives for which every upper-band neighbourhood \(U_l\) lies in the open upper union \(X_L^*\subset X'\). Their union is then exactly \(X_L^*\). Consequently the union of the pairwise cover pieces \(U_j\cap U_l\) is exactly \(V_J\cap X_L^*=V_J^L\) at that stage. All flags in \(D\) have maximum in \(L\), so their intersections are actual opens of this union. Thus no homotopy invariance of arbitrary sheaf cohomology is being invoked. Choose these representatives for the finite open cover and a flabby resolution of \(F\). Form the double complex whose degree \((p,q)\) is the product of \(\Gamma(V_{d_p},\mathcal I^q)\) over chains \(d_0\subset\cdots\subset d_p\), with the usual alternating restriction differential horizontally. Its horizontal augmentation computes sections on the union. Exactness can be checked locally inside a covering piece: adjoining that piece to an intersection gives the common refinement and contracts the augmented \v{C}ech complex; empty intersections contribute zero. More precisely use the augmented complex of sheaves obtained by direct image from the open intersections. On a covering open set, contraction is by adjoining its flag; a nonflag union has empty intersection and contributes zero. This local contraction proves exactness without identifying the stalk of an open direct image outside its open set with zero. With a flabby resolution, all the open direct images are flabby. Flabbiness makes all restriction terms acyclic for sections, so the total complex computes sheaf cohomology on the union.

Taking vertical cohomology first yields the derived-limit complex of the coefficient diagram \(d\mapsto H^q(V_d,F)\), proving the displayed \(E_2\)-page. Taking filtered colimits over representative neighborhoods is exact in abelian groups and commutes with the finite products here. Finally \(D\) has bounded chain length, so the horizontal filtration is finite; there is no convergence or product-completion ambiguity. Empty \(D\) gives the empty union and zero groups.

\subsection{Homotopical stratification need not give conical collars}\label{c4.04-refuted-hss-alone-gives-conical-collars}

Quinn's example \cite[p.452]{Quinn} is the shrinking hedgehog \[
 X=\bigvee_{n\ge1}[0,1/n],\qquad Y=\{0\},
\] with the metric \(|s-t|\) on one arm and \(s+t\) on distinct arms. It is compact metric. The radial contraction \(h((n,s),t)=(n,(1-t)s)\) makes \(Y\) tame, strictly off \(Y\) until the final time; the holink map to a point is a fibration. Thus the two-step filtration is homotopically stratified in Quinn's general definition.

It has no mapping-cylinder neighborhood of \(Y\). Every neighborhood contains all sufficiently short arms, including their terminal endpoints, and these endpoints converge to \(Y\). In a mapping-cylinder neighborhood with radial coordinate \(r\), a point with \(0<r<1\) lies on an embedded interval through itself, obtained by varying \(r\) in both directions. A terminal endpoint of an arm cannot be an interior point of any embedded interval in \(X\). Hence every such endpoint would have to lie on the outer level \(r=1\). Continuity is contradicted by their convergence to the apex \(r=0\). The same argument applies after making the cylinder smaller.

Therefore \(\Sigma_d\) in the actual collar/mapping-cylinder sense need not exist for a general HSS. The germ comparisons, homotopy-link comparison, and finite spectral sequence have the separate proofs above. An equality \(V_d\simeq\Sigma_d\) requires an actual compatible conical deployment, or must define a replacement \(\Sigma_d\) using the full iterated holink diagram. The example uses Quinn's general strata, not manifold strata; no claim is made here that it satisfies an additional manifold-strata hypothesis.

\subsection{Germ and face}\label{c4.05-proved-with-proper-controlled-collars-germ-and-face}

Use the actual finite proper conical deployment with compatible iterated fibre collars of the iterated boundary-fibration hypotheses above, in the metric setting of Section~\ref{c4.01-proved-for-finite-metric-hss-the-germ-comparisons}. We spell out the radial data used in this assertion. Each stratum \(S_i\) has a fixed control tube \(N_i\), a retraction \(\pi_i:N_i\to S_i\), and a radius \(\rho_i:N_i\to[0,1)\) with zero set \(S_i\). Each closed subtube \(0\le\rho_i\le a<1\) is proper over \(S_i\times[0,a]\) under \((\pi_i,\rho_i)\). This follows from the proper compact link bundle in its mapping-cone presentation: the product of that bundle with \([0,a]\), followed by the cone quotient, is proper.

On the resolution \(\beta_j:\Sigma_j\to\overline{S_j}\), the interior is identified with \(S_j\). If the original construction retains a cut core instead, its collar stretching gives the same identification: append an outer collar with coordinate in \((-1,0]\), combine it with an inner collar \([0,a)\), and choose a monotone homeomorphism \((-1,a)\to(0,a)\) equal to the identity near \(a\). Extend by the identity outside the collar. At corners apply these reparametrizations in the commuting product coordinates. Retained outer faces in intermediate induction diagrams are handled by the inward-collar argument in Theorem~\ref{thm:exit-deployment}. The compatible iterated fibre collars give the following precise simultaneous coordinates. For
\(d=(i_0<\cdots<i_{n-1}<j)\), put
\[
 F_d=\bigcap_{k<n}\Sigma_{i_kj}=\Sigma_d,
 \qquad F_d^\circ=F_d\setminus
       \bigcup_{\ell\notin\{i_0,\ldots,i_{n-1},j\}}
                  (F_d\cap\Sigma_{\ell j}).
\]
Thus all remaining incoming boundary faces, including intermediate ones not selected by \(d\), are removed in \(F_d^\circ\). After fixing compatible collar widths, the intersection of the selected collar neighborhoods in the interior is identified with
\[
 c_d:F_d^\circ\times(0,a)^n\xrightarrow{\cong}
           S_j\cap\bigcap_{k<n}N_{i_k}.
\]
Here and below \(c_d\) includes the blowdown on the interior. The widths may first be normalized by the given collar coordinates. Their compatibility with the boundary fibrations gives continuous maps
\[
 p_k:F_d^\circ\times(0,a)^k\longrightarrow S_{i_k}
\]
with
\[
 \rho_{i_k}c_d(z,r)=r_k,
 \qquad
 \pi_{i_k}c_d(z,r)=p_k(z,r_0,\ldots,r_{k-1}).
 \tag{RC1}
\]
The absence of later radii in the second formula is essential: the later collars are vertical for the earlier source fibration. These are the simultaneous product and fibre-compatibility axioms of the actual resolved collar system, equivalently the compatible radial control constructed in the iterated boundary-fibration hypotheses above. Pairwise commutativity of arbitrary retractions alone would not supply (RC1).

Here is the translation from the resolved data of the iterated boundary-fibration hypotheses above. On the face \(\Sigma_{i_kj}\), the faces with labels \(i_\ell<i_k\) are the inverse images of the corresponding boundary faces of the base \(\Sigma_{i_k}\); this is the flag inverse-image identity. Their collars are carried to the base collars. The faces with labels \(i_\ell>i_k\) are vertical faces of the source fibre, and their compatible collars leave that source map fixed. Applying the commuting collar maps successively therefore carries the source only through the earlier coordinates, exactly as in (RC1). The transverse product-collar axiom identifies intersections of the chosen collar neighborhoods with the simultaneous collar of their face intersection, after shrinking the finite collar system. It also preserves every nonselected face. Removing all boundary faces to pass to \(S_j\) consequently removes the nonselected faces from the base \(F_d\), giving precisely \(F_d^\circ\), and requires each selected radius to be positive. This proves the displayed coordinate description from the full iterated fibre-collar data, rather than assuming it for arbitrary retractions.

\textbf{A cofinal family consisting of actual neighborhoods.} For every positive continuous function \(\epsilon_i:S_i\to(0,a)\), let
\[
 T_i(\epsilon_i)=
 \{x\in N_i:\rho_i(x)<\epsilon_i(\pi_i(x))\}.
\]
These are genuine open neighborhoods of \(S_i\) in the required open complement of its lower frontier. They form a cofinal family. Indeed, for a prescribed neighborhood \(U_i\), choose a fixed smaller closed subtube \(K_i=\{\rho_i\le a'\}\), \(0<a'<a\). The set \(K_i\setminus U_i\) is closed; properness makes its image under \((\pi_i,\rho_i)\) a closed subset of \(S_i\times[0,a']\) disjoint from the zero section. There is a positive continuous function \(\epsilon_i<a'\) whose subgraph misses that image. To construct it explicitly, choose a cover of \(S_i\) by open sets \(B_\alpha\) and numbers \(b_\alpha>0\) for which \(B_\alpha\times[0,b_\alpha]\) misses the closed image. A locally finite subordinate partition of unity gives \(\epsilon_i(y)=\frac12\sum_\alpha\lambda_\alpha(y)b_\alpha\), capped at \(a'/2\). Each active bound is admissible at \(y\), hence so is their weighted average. Then \(T_i(\epsilon_i)\subset U_i\).

Consequently the germ \(V_d\cap S_j\) is represented by the actual intersections
\[
 H_\epsilon=S_j\cap\bigcap_{k<n}T_{i_k}(\epsilon_{i_k}).
\]
The last neighborhood, that of \(S_j\) itself, imposes no restriction after intersecting with \(S_j\). In the fixed simultaneous coordinates, (RC1) identifies \(H_\epsilon\) with
\[
 \left\{(z,r):
 0<r_k<\epsilon_{i_k}\bigl(p_k(z,r_0,\ldots,r_{k-1})\bigr)
 \quad(0\le k<n)\right\}.\tag{RC2}
\]
These are recursively varying radius regions. The family is cofinal in both required senses: every member is exactly an actual intersection of the indicated neighborhood representatives, and the preceding properness argument puts some such member inside every prescribed tuple of neighborhoods. We do \emph{not} claim that arbitrary positive functions on the entire resolved face or on fibre directions produce cofinal ambient neighborhood systems.

\textbf{The homotopy equivalence and its compatibility.} Each region (RC2) is homeomorphic to \(F_d^\circ\times(0,1)^n\). The homeomorphism is given successively by
\[
 r_k=\epsilon_{i_k}\bigl(p_k(z,r_0,\ldots,r_{k-1})\bigr)u_k,
 \qquad 0<u_k<1.
\]
Its inverse divides by the same strictly positive continuous functions, successively in increasing \(k\). Projection to \(z\) is therefore a homotopy equivalence. All the transition inclusions for smaller control functions commute with this one fixed projection. We obtain a natural objectwise equivalence from the cofinal diagram of the \(H_\epsilon\) to the constant space \(F_d^\circ\), and hence an equivalence of prospaces, with full coherent compatibility. There is no inference merely from unrelated isomorphisms in \(\mathrm{HoTop}\).

Finally \(F_d^\circ\hookrightarrow F_d\) is a homotopy equivalence. The finitely many remaining commuting collars can be pushed inward, using in each coordinate a function positive at zero and equal to that coordinate outside a smaller collar. They preserve the selected faces because their collars commute. The resulting global homotopy has image in \(F_d^\circ\) at its final time and preserves \(F_d^\circ\) at every time. Thus
\[
 V_d\cap S_j\simeq F_d^\circ\simeq F_d=\Sigma_d.
\]
Section~\ref{c4.01-proved-for-finite-metric-hss-the-germ-comparisons} supplies \(V_d\simeq V_d\cap S_j\). Empty face intersections have an empty cofinal neighborhood intersection; the same statement then holds with the empty space. For a one-element flag the assertion reduces to the inward-collar equivalence between the resolved piece and its stratum.

The rectangular version must be qualified. An arbitrary neighborhood of an intermediate stratum can impose \(r_1<r_0^2\) near a lower stratum, so no product rectangle with two fixed positive radii is contained in it. Formula (RC2), with radius functions pulled back from the actual stratum bases, is what makes the comparison valid. The resulting equivalence is canonical in the homotopy category for the fixed compatible deployed collar germs; the contractible choices of smaller radius functions and the displayed natural projection specify its compatible homotopies.

\subsection{The canonical normal-crossings comparison}\label{c4.06-proved-the-canonical-normal-crossings-comparison}

Let \(B=D_{d_*}^*\) denote a chosen open intersection stratum of a smooth normal-crossings arrangement, with local branches retained. Consider a filtered normal bundle over \(B\), with two entries \(d_0''\le d_0'\) of a flag and an exact sequence \[
 0\longrightarrow A=\check N_{d_*,d_0'}
 \longrightarrow E=\check N_{d_*,d_0''}
 \xrightarrow{q}Q=\check N_{d_*,(d_0',d_0'')}
 \longrightarrow0.
\] The normal-crossings branches give \(Q\) a local decomposition into complex lines, with transitions allowed to permute those lines. Write \(Q^*\) for the complement of all coordinate hyperplanes. Thus its fibre is \((\mathbb C^*)^r\), not \(\mathbb C^r\setminus\{0\}\). For rank zero, the empty coordinate arrangement has complement the whole zero bundle; the torus fibre is a point. The same argument then applies without a puncture, and an empty base gives the empty equivalence. Define \(E^*=q^{-1}(Q^*)\); then \(E^*\to Q^*\) is a torsor for the pulled-back additive bundle \(A\).

The quotient map \(q:E^*\to Q^*\) is a canonical homotopy equivalence. Choose a Hermitian splitting \(E=A\oplus Q\); then \((a,v)\mapsto((1-t)a,v)\) contracts the kernel and proves the claim. Splittings form an affine space, so their homotopies identify all inverse choices. The map \(q\) itself does not depend on a splitting.

The sector tube \(V^{d_0'}_{d_*,d_0''}\) is, by definition of the normal-crossings incidence sector, the neighborhood germ of \(B\) in the \(E\)-directions, with precisely the coordinates belonging to \(Q\) required to be nonzero. Its pro-homotopy type has a canonical comparison with \(E^*\). Here is a construction and a check of its ambiguity. For a simple normal crossings divisor, use its divisor line bundles with defining sections. Choose an ordinary smooth tubular projection onto the intersection \(B\) and connections on those line bundles. Parallel transport identifies their restrictions on a small neighborhood with their pullbacks from \(B\); the tuple of defining sections then has differential the normal-coordinate isomorphism along \(B\). The inverse function theorem, after variable-radius shrinking, gives an adapted tubular germ from the direct sum of the normal lines to the original space. It takes every coordinate hyperplane and every flag subbundle to the corresponding divisor incidence. This constructs the required sector comparison.

Normalize each adapted tubular germ to have identity derivative on the intrinsic normal bundle, by composing with the inverse of its fibrewise normal derivative. This preserves the coordinate arrangement and the flag subbundles. Arbitrary normal-bundle automorphisms are not being declared homotopic to the identity. If a kernel block has no chosen line decomposition, retain its coordinates from an ordinary tubular chart and use the defining sections in the quotient block; the same inverse-function argument respects every flag subbundle.

If \(\Phi,\Psi\) are two adapted tubular germs with the same normal derivative, their transition \(g=\Psi^{-1}\Phi\) fixes the zero section and has identity normal derivative. Fibre dilations give \[
 g_t=\delta_{1/t}\,g\,\delta_t\quad(t>0),\qquad g_0=\mathrm{id}.
\] Taylor's formula in normal coordinates shows this is a continuous, indeed smooth, family of germs through \(t=0\); its tangential displacement tends to zero and its normal derivative tends to the identity. Every \(g_t\) preserves the coordinate-hyperplane arrangement. Restricting to sufficiently small variable-radius neighborhoods makes the family defined for all \(t\in[0,1]\). The same construction works for parameter families. Hence the induced comparison on the sector has a unique homotopy class, compatible with refinements and changes of tubular choices. This proves independence of the tubular choices at the level of homotopy classes.

The small-sector pro-germ also compares explicitly with the entire punctured bundle. For a bundle norm and a positive continuous radius function \(\eta\) on the base, the inclusion \(\{v\in E^*:\|v\|<\eta(\pi v)\}\hookrightarrow E^*\) is a homotopy equivalence. An inverse up to homotopy is \[
 v\longmapsto \frac{\eta(\pi v)}{\eta(\pi v)+\|v\|}v.
\] Interpolation of its strictly positive scalar factor with \(1\) supplies both inverse homotopies and preserves every coordinate nonvanishing condition. Such variable-radius disk neighborhoods are cofinal by the tubular chart and paracompactness of the base. Their natural inclusions therefore identify the sector prospace with the constant space \(E^*\), coherently under shrinking.

Combining it with \(q\) gives the precise statement \[
 V^{d_0'}_{d_*,d_0''}
 \ \simeq\
 \left.\check N_{d_*,(d_0',d_0'')}^{\,*}\right|_{D_{d_*}^*}.
\] The ambiguity is homotopy, not a distinguished homeomorphism of neighborhoods. A Hermitian norm further retracts each punctured line to its unit-circle bundle, giving the corresponding torus bundle; metrics form a convex space of choices.

For normal crossings with branch monodromy, perform the construction on the finite local branch-label cover. The normal line bundles and the quotient map descend equivariantly, as does the homotopy-class comparison, because the coordinate arrangement and the dilation are invariant under branch permutations. On a Deligne--Mumford orbifold chart finite-group averaging supplies invariant auxiliary metrics and connections. Thus the result is canonical in the incidence/orbifold category. Passing to an unlabelled orbit-poset and discarding stabilizers would lose data.


\begin{proof}[Proof of Corollary~\ref{cor:holink-pushout}]
For every open $N\supset Y$, the numerable open cover $X=N\cup U$
gives $N\cup^h_{N\setminus Y}U\simeq X$. A fixed tame deformation
identifies the neighborhood prospace with the constant $Y$. The path
shortening and terminal evaluation proof above, applied to this pair,
identifies the punctured-neighborhood prospace with the constant $H$,
including both endpoint maps and the path homotopy between their
composites to $X$. Finite span colimits commute with a synchronized
pro-presentation because filtered colimits of spaces commute with finite
limits. Hence $cX\simeq c(Y\cup^h_HU)$; the constant embedding is fully
faithful. The resulting map of ordinary double cylinders is path
evaluation. This proof does not identify Quinn's controlled-cylinder
topology with the ordinary quotient-cylinder topology.
\end{proof}
The endpoint fibration condition alone does not give tameness. For the
Hawaiian earring and its common point the endpoint map is a fibration
over a point, but a tame neighborhood contraction would kill every
sufficiently small circle. Retraction onto any chosen circle shows that
its loop is nonzero, a contradiction.


\section{What changes in TOP, PL, and smooth}\label{sec:categories}
The preceding constructions use several notions of regularity that
must be distinguished. A Hurewicz fibration need not be a bundle, a
bundle factor of a topological manifold need not be a manifold, and a
smooth map that is a Hurewicz fibration need not be a submersion.
For general Hausdorff spaces here, a proper map means a closed map
with compact fibers. For locally compact Hausdorff spaces this agrees
with the compact-inverse-image convention. Closedness is used explicitly
in the finite-fiber bound below; no conclusion is claimed from a weaker
compact-inverse-image condition on an arbitrary Hausdorff target.

\begin{theorem}[Components of a topological order]\label{thm:components}
Let $P$ be Hausdorff and locally path connected, with a closed partial
order $R\subset P\times P$. Suppose its target $b:R\to P$ is a proper
local embedding and its source $s:R\to P$ is a proper Hurewicz fibration.
Distinct comparable points lie in different path components, and the
induced relation on components is a partial order. If every lower
interval is linear, the target on arrows between two fixed components
is a closed embedding.
\end{theorem}
\begin{proof}
The diagonal is closed by Hausdorffness and open by local injectivity
of $b$ near identities. Thus strict arrows form a clopen subspace.
The fibers of $b$ are compact discrete sets, hence finite. Their
cardinality $N(y)$ is locally bounded: choose injective neighborhoods
of the finitely many points over $y$ and remove the closed image of
their complement. Near $y$ every fiber meets only those neighborhoods.
For fixed $x$, the upper interval $b(s^{-1}(x))$ is compact, so $N$ is
bounded on it.

If $x<y$ lie in one path component $C$, lift any path in $C$ starting
at $x$ through $s$, beginning at that strict arrow. The lift stays
strict and its target stays in $C$. Every point of $C$ therefore has
a strict successor in $C$. Arbitrarily long chains starting at $x$
would have endpoints in its compact upper interval, with arbitrarily
large $N$, a contradiction. To compose component arrows, lift a path
in the intermediate component to align their representatives.
Transitivity follows; a component cycle gives a comparable pair in
one component and then antisymmetry in $P$ makes the cycle trivial.
Finally two arrows from $C$ to one target have comparable sources in
its linear lower interval, hence identical sources. Thinness makes
the arrows identical. Components are clopen, so the restricted target
remains proper; a proper injection into a Hausdorff space is a closed
embedding.
\end{proof}

\begin{counterexample}[Automatic product transversality]\label{ex:transverse}
Projection control alone does not give the facewise product charts needed
in a deployment. Smooth Hurewicz fibrancy does not imply a smooth
normal-product condition either.
\end{counterexample}
\begin{proof}
Put $\Sigma_0=*$, $\Sigma_1=\R\times[0,\infty)$ and
$\Sigma_2=\R\times[0,\infty)^2$. Take the two coordinate boundary
faces of $\Sigma_2$ and define on the incoming face
\[
 s_{12}(x,t)=(xt,t)
\]
in the smooth example, or
$s_{12}(x,t)=(\max(-t,\min(x,t)),t)$ in the PL example.
The other sources are the unique maps to the point. Subset incidence
and projection control hold. At $t=0$ the whole boundary axis maps
to one point; at positive $t$ its image has two-dimensional interior.
A product collar through the one-point boundary image would have
one-dimensional image, a contradiction. Differential transversality
to the target boundary is weaker here: the last coordinate is exactly
$t$. For the proper smooth refinement take $\Sigma_0=*$,
$\Sigma_1=[0,1]$, $\Sigma_2=[0,1]^2$, and all four boundary edges,
with $\Sigma_{01}=\{0,1\}$,
$\Sigma_{12}=[0,1]\times\{0,1\}$ and
$\Sigma_{02}=\{0,1\}\times[0,1]$. Put
$s_{12}(u,v)=u^2$ and let the other sources be constant. On each
component this source is a homeomorphism, hence a proper Hurewicz
fibration. All target collars, intersections and projection identities
hold. Its normal derivative at $u=0$ is zero; smooth collar changes
multiply that derivative by nonzero factors and cannot make it the
nonzero derivative of a product collar. Thus the smooth normal-product
conclusion also requires differential regularity. In every category,
local triviality of the entire face diagram over a base collar is a
sufficient correction: pull the whole diagram back along that product
coordinate. Mere source fibrancy is not that condition.
\end{proof}

\begin{counterexample}[Topological manifold factors]\label{ex:cannon}
There are proper topological source bundles between manifold object
spaces whose fibers are not manifolds. There are transverse collective
faces in a topological manifold whose specified individual face spaces
are not manifolds with boundary.
\end{counterexample}
\begin{proof}
Let $H$ be a homology three-sphere with nontrivial fundamental group
and put $F=\Sigma H$. The suspension vertices are not manifold points:
punctured conical neighborhoods have fundamental group $\pi_1(H)$,
with unchanged groups under shrinking. Cannon's double-suspension
theorem \cite[Theorem~1.5, p.~833]{Cannon} gives
$\Sigma^2H\cong S^5$. Deleting the two outer suspension points shows
$F\times\R\cong S^4\times\R$, a five-manifold. Therefore
$E=F\times S^1$ is a closed topological five-manifold.
Take $\Sigma_0=S^1$, $\Sigma_1=E\times[0,1]$, and its entire two-ended
boundary as incoming face, with source the projection to $S^1$.
It is a proper bundle with fiber $F\amalg F$, which is not a manifold.

For the face assertion use $\Sigma_0=*$, $\Sigma_1=[0,1]$ and
$\Sigma_2=F\times[0,1]^2$. Since the square is a disc, its product
with $F$ is a topological six-manifold with boundary
$F\times S^1$. Its two collective coordinate faces have commuting
product collars. The sources are coordinate projections, so the subset
and projection conditions and proper fibrancy hold. An individual face
$F\times[0,1]$ is not a manifold with boundary at all. At an interior
interval coordinate it is a five-manifold. At an endpoint, all local
relative homology groups vanish: a small half-interval neighborhood
and its puncture are both contractible, and the product relative
complex therefore has zero homology. If the product were a manifold
with boundary, its intrinsic boundary would consequently be exactly
$F\times\{0,1\}$. The boundary of a manifold with boundary is a
manifold, contrary to the singular suspension vertices of $F$.
This proves why ambient TOP manifoldness cannot replace the exact
specified-face criterion of Proposition~\ref{prop:corners}.
\end{proof}

\begin{proposition}[Positive regularity hypotheses]\label{prop:positive-regular}
For PL source bundles between PL manifold pieces, fibers are PL
manifolds. For proper smooth source systems with product-transverse
base faces, projection control, and submersions on every full source
face, the fiber faces are smooth manifolds with corners and the sources
are locally trivial bundles respecting those faces.
\end{proposition}
\begin{proof}
In PL a local bundle chart gives $F\times\R^m$ as a PL manifold.
Choose a product triangulation and take the link at a vertex of $F$
times an interior point of an $m$-simplex. Its transverse link is the
link of that vertex in $F$; the manifold link condition says precisely
that it is a sphere or ball. This gives the required PL charts in $F$.

In smooth geometry, a horizontal boundary is a pullback of a base
boundary and is regular by the assumed product charts. A vertical
face for $i<k<j$ satisfies
$s_{ij}|_{\Sigma_{ikj}}=s_{ik}s_{kj}|_{\Sigma_{ikj}}$.
The first factor is a submersion by hypothesis; horizontal base change
for the second preserves surjectivity. Iteration treats all flag faces.
Thus restrictions to every vertical face are submersions. Choose local
horizontal lifts of base vector fields tangent to all these faces,
and combine them with a partition of unity. Properness makes the
required flows exist uniformly over small base neighborhoods.
They give the usual Ehresmann product charts, simultaneously respecting
the fiber faces. The full boundary of an arrow space is
\[
 \partial\Sigma_{ij}
 =\bigcup_{k<i}\Sigma_{kij}\ \cup\!
   \bigcup_{i<k<j}\Sigma_{ikj}.
\]
The first terms are horizontal pullbacks of base boundaries. The second
are vertical, and precisely these intermediate-label terms form the
boundary of a source fiber. Intersections of the displayed terms give
all its corners; no other faces occur by the full boundary decomposition.

For manifold incoming faces alone no source bundle is needed in PL or
smooth. A transverse PL collar product has normal-coordinate simplices
whose transverse links are the links of the face; the same PL link
criterion makes each face intersection a manifold with corners. In
smooth collar charts, setting its independent normal coordinates to
zero restricts to smooth face charts. Deleting the other incoming faces
through a point gives exactly its open face and boundary. These
arguments also cover empty faces and dimensions zero, one and two.
They justify the positive face rows of the comparison table without
using TOP factor cancellation.
\end{proof}

\subsection{Proper PL and smooth Hurewicz maps with singular fibers}\label{sec:singular-fibers}
We prove Theorem~\ref{thm:singular-fibers}. Its PL and smooth realizations are constructed separately below.

\paragraph{A contractible five-manifold with nontrivial boundary group.}
Take a smooth compact contractible Mazur four-manifold $A$ such that
$G=\pi_1(\partial A)\ne1$.  A primary construction is Zeeman,
\emph{On the dunce hat}, printed p.~347, the passage preceding Theorem~(5),
which describes the handle construction, records the nontrivial boundary
group, and gives the original locator Mazur~1961, p.~227
\cite[p.~347]{Zeeman}.  Here Zeeman's dotted-manifold notation denotes
the boundary.  The handle attachment can be made smooth and corners rounded.
Choose a smoothly embedded closed three-disc $D^3\subset\partial A$ and set
\[
 W=(A\times S^1)\mathop{\cup}_{D^3\times S^1}(D^3\times D^2),
 \qquad H=\partial W.
 \tag{Q1}
\]
Round the corners of this attachment.  Van Kampen gives
$\pi_1(W)=\mathbb Z*_{\mathbb Z}1=1$, since the attaching circle generates
$\pi_1(A\times S^1)$.  In the reduced Mayer--Vietoris sequence the only
possibly nonzero positive homology map is the isomorphism
$H_1(D^3\times S^1)\to H_1(A\times S^1)$.  Thus $W$ is acyclic.
A simply connected acyclic finite CW complex is contractible, by the
Hurewicz and Whitehead theorems, so $W$ is contractible.

The boundary is
\[
 H=\bigl((\partial A\setminus\operatorname{int}D^3)\times S^1\bigr)
       \mathop{\cup}_{S^2\times S^1}(S^2\times D^2).
 \tag{Q2}
\]
Removing a three-disc does not change the fundamental group of the closed
three-manifold $\partial A$.  The second van Kampen calculation consequently
reads
\[
 \pi_1(H)=(G\times\mathbb Z)*_{\mathbb Z}1\cong G\ne1.
\]
This construction is a boundary spin, not stabilization by an interval.
In particular it does not make the invalid inference that
$\partial(A\times I)$ retains $\pi_1(\partial A)$.

Write $M=DW$ for the double of $W$.  This is a simply connected smooth
homology five-sphere.  More specifically, the rounded manifold $W\times I$
is contractible and its boundary is $DW$.  Apply the smooth $h$-cobordism
theorem after deleting an interior six-disc: the resulting six-dimensional
cobordism has simply connected boundary components $S^5$ and $DW$, and their
inclusions are homology equivalences.  It is a product.  Hence
\[
 W\times I\cong D^6,\qquad DW\cong S^5.
 \tag{Q3}
\]
The dimension restriction is satisfied.  We use the implication
``contractible, simply connected boundary $\Rightarrow$ smooth disc'' in
Milnor's Proposition~A, pp.~108--109, following Theorem~9.1, p.~107
\cite[\S9]{Milnor}.

\paragraph{The quotient map and its exceptional fibre.}
Use one copy of $W\subset M=DW$ and fix $0\in S^1$.  Form
\[
 E=(M\times S^1)/(W\times\{0\}),\qquad
 p([x,t])=t,
 \qquad Y=M/W.
 \tag{Q4}
\]
Only the entire set $W\times\{0\}$ is identified to one point $v$; all
other equivalence classes are singletons.  The quotient is compact
metrizable.  All fibres except $p^{-1}(0)$ are copies of $M$, whereas
$p^{-1}(0)=Y$.  An outward collar of $H=\partial W$ identifies a neighbourhood
of the image of $W$ in $Y$ with an open cone on $H$.
The local homology there is that of an interior point of a five-manifold,
because $H$ is a homology four-sphere.  But the punctured conical
neighbourhoods have fundamental group $G\ne1$, and all inclusions between
smaller and larger punctured conical neighbourhoods induce isomorphisms.
This contradicts local simple connectivity at an interior point of a
five-manifold.  The local homology also excludes a boundary point.  Thus
$Y$ is not a topological manifold, with or without boundary.

\paragraph{The Hurewicz lifting property: complete control argument.}
We prove this for the ordinary quotient (Q4), without assuming its
collapse map $f:M\to Y$ is PL.  A small closed collar enlargement $A_\eta$ of
$W$ is a compact contractible manifold, and
\[
 B_\eta=A_\eta/W\cong C(H)
\]
is a compact polyhedron.  Choose a contraction of $W$ within itself.
The pair $(A_\eta,W\cup\partial A_\eta)$ is a cofibration pair, since all
these subspaces have compatible collar triangulations.  Extend the
contraction, together with the stationary homotopy on $\partial A_\eta$,
to a homotopy $H_t:A_\eta\to A_\eta$ starting at the identity.  It satisfies
$H_t(W)\subset W$ for all $t$ and $H_1|W$ is constant.  Therefore $H_1$
factors through a map $s_\eta:B_\eta\to A_\eta$, fixed on the outer
boundary.  Extend $s_\eta$ outside $B_\eta$ by the inverse of $f$, obtaining
$s_\eta:Y\to M$.

The homotopy $H_t$, extended by the identity outside $A_\eta$, gives
$\mathrm{id}_M\simeq s_\eta f$, supported in $A_\eta$.  Because
$H_t(W)\subset W$ at every time, $fH_t$ descends to a homotopy
$\mathrm{id}_Y\simeq f s_\eta$, supported in $B_\eta$.  Both are fixed on
their outer boundaries.  These are the required controlled inverse
homotopies; they use only the homotopy extension property of a finite
polyhedral pair.

Equip $E$ with any compatible metric.  Given $\epsilon>0$, first choose
$\eta$ and then $\delta>0$ so small that
\[
 \operatorname{diam}\bigl(q(A_\eta\times[-\delta,\delta])\bigr)<\epsilon,
 \tag{Q5}
\]
where $q:M\times S^1\to E$ is the quotient and the signed coordinate is
chosen in a small arc about $0$.  Such choices exist because these compact
sets shrink to $q(W\times\{0\})=\{v\}$.  Also arrange
\[
 \sup_{x\in M,\ |t|\le\delta}
    d\bigl(q(x,t),q(x,0)\bigr)<\epsilon,
\]
by uniform continuity of $q$.  The maps between the fibres at $0$ and $t$
are $f$ and $s_\eta$, placed in these respective levels.  Outside the
supports they move points less than $\epsilon$ by the latter inequality;
inside the supports they and both inverse homotopies move points less than
$\epsilon$ by (Q5).  At every nonzero base point, ordinary product
identifications of nearby fibres give the same property, with stationary
inverse homotopies.

Thus $p$ is strongly regular in the precise sense of Ferry's definition on
p.~373.  Its fibres are compact ANRs: $M$ is a manifold and $Y$ is a finite
polyhedron, obtained by attaching the cone on $H$ to the complementary
collared manifold.  The base $S^1$ is complete and finite dimensional.
Ferry's Theorem~1, pp.~373--374, therefore proves that $p$ is a Hurewicz
fibration \cite{Ferry}.  This establishes the full lifting property,
not merely a quasifibration or a statement about homotopy-equivalent fibres.

\paragraph{A PL total space and a PL map.}
Triangulate the smooth pair $(M,W)$ compatibly and choose a PL outward
collar $H\times[0,a]$.  In $M\times(-2a,2a)$ take the polyhedral
neighbourhood
\[
 N=(W\times[-a,a])\ \cup\
   \{((z,s),t):z\in H,\ 0\le s\le a,\ s+|t|\le a\}.
 \tag{Q6}
\]
It is a genuine neighbourhood of $W\times\{0\}$, including an outward
collar; the inadequate set $W\times[-a,a]$ alone would touch $W\times\{0\}$
at its boundary.  The boundary $L=\partial N$ is the double of $W$ with an
extra collar, hence is a PL five-sphere by (Q3).  Let
$h:L\to[-a,a]$ be the restriction of the coordinate $t$, a PL function.
Replace $N/(W\times\{0\})$ by the abstract cone $C(L)$ and put
\[
 \widetilde p(r,\theta)=r h(\theta),\qquad \widetilde p(v)=0.
 \tag{Q7}
\]
This is the linear cone extension of a PL function and is PL.  Attach the
cone to $(M\times S^1)\setminus\operatorname{int}N$ by the identity on $L$,
keeping the original projection outside.  Since $C(L)$ is a PL six-disc,
the resulting space $E_{\mathrm{PL}}$ is a closed PL six-manifold and
$p_{\mathrm{PL}}$ is PL.

Here is an exact height-preserving identification with (Q4).  For a
boundary point $(x,\pm a)$ with $x\in W$, send
\[
 (r,(x,\pm a))\longmapsto q(x,\pm ar).
\]
For a side-boundary point $((z,s),t)$ with $s+|t|=a$, send
\[
 (r,((z,s),t))\longmapsto q((z,rs),rt).
\]
The formulas agree on their common faces and send the whole cone vertex
to $v$.  Every point off $v$ has a unique preimage: use $r=|t|/a$ over $W$,
and $r=(s+|t|)/a$ on its outward collar.  The continuous bijection between
compact Hausdorff spaces is a homeomorphism, identity on $L$, and it
preserves height exactly.  Extending it outside $N$ identifies
$p_{\mathrm{PL}}$ with $p$.  Hurewicz lifting and properness therefore pass
to the PL model.

\paragraph{A smooth total space and a smooth map.}
Choose a smooth outward collar coordinate $s$ and smooth functions
$\alpha,\beta$ which vanish for $s\le0$, are positive for $s>0$, are flat
at zero, and satisfy $\alpha'(s)>0$ for $s>0$.  On $W$ take $\alpha=0$ and
the horizontal vector field to be zero.  For small $c>0$ put
\[
 N_{\mathrm{sm}}=\{\alpha(s)+t^2\le c\},\qquad
 V=\beta(s)\partial_s+t\partial_t,
 \tag{Q8}
\]
with the first summand zero on $W$.  Choose the collar long enough that
$\alpha$ exceeds $c$ before its outer end.  The boundary $L_{\mathrm{sm}}$
is smooth; the derivative of $\alpha(s)+t^2$ along $V$ is
$\alpha'(s)\beta(s)+2t^2$, positive except on $W\times\{0\}$.
Both halves of $L_{\mathrm{sm}}$ are diffeomorphic to $W$, so it is
$DW\cong S^5$.  The function $h=t|L_{\mathrm{sm}}$ is smooth and has zero
as a regular value: at $t=0$, $\partial_t$ is tangent to $L_{\mathrm{sm}}$
and $dh(\partial_t)=1$.

Every forward orbit from $N_{\mathrm{sm}}\setminus(W\times\{0\})$ exits
through $L_{\mathrm{sm}}$ in finite time.  For $t\ne0$ this follows from
$V(t)=t$; for $t=0,s>0$ it follows from $\beta(s)>0$.
Backward orbits converge uniformly towards the compact set
$W\times\{0\}$: $t$ decreases exponentially and the exterior collar
coordinate decreases monotonically to zero.  The flow consequently gives
a diffeomorphism
\[
 L_{\mathrm{sm}}\times(0,1]\longrightarrow
 N_{\mathrm{sm}}\setminus(W\times\{0\}),\qquad
 (\theta,r)\longmapsto\Phi_{\log r}(\theta),
\]
and after collapse a homeomorphism $C(L_{\mathrm{sm}})\cong
N_{\mathrm{sm}}/(W\times\{0\})$.  It is the identity on the boundary and
satisfies the exact formula $p(r,\theta)=rh(\theta)$.
Using a diffeomorphism $L_{\mathrm{sm}}\cong S^5$ gives a smooth disc
structure at the cone point, compatible off that point with the original
smooth structure through the flow chart.  The flow extends across the
boundary, so these charts glue smoothly to the exterior.

Choose a strictly increasing smooth $g:[0,1]\to[0,1]$ with $g(0)=0$,
$g(r)=\exp(-1/r^2)$ near zero, $g(r)=r$ near one, and $g'(r)>0$ for $r>0$.
The radial homeomorphism $(r,\theta)\mapsto(g(r),\theta)$ is the identity
near the boundary.  Precompose $p$ with it.  The resulting local function
is
\[
 p_{\mathrm{sm}}(r,\theta)=g(r)h(\theta).
 \tag{Q9}
\]
In ordinary Euclidean cone coordinates this is smooth and flat at the
origin: derivatives of the angular factor have at most power growth in
$r^{-1}$ and are dominated by the flat radial factor.  Off the origin it
is a submersion: use $g'(r)h(\theta)\ne0$ when $h\ne0$, and $dh\ne0$ when
$h=0$.  The domain homeomorphism preserves the Hurewicz property and
properness, and the exceptional fibre still has local model $C(H)$.
This proves the smooth assertion of the counterexample.

\paragraph{Insertion into two-stratum deployment data.}
For either of the two categories, put
\[
 \Sigma_0=S^1,\qquad \Sigma_1=E\times[0,1],\qquad
 \Sigma_{01}=E\times\{0,1\}.
\]
The target is the whole-boundary embedding
$\Sigma_{01}=\partial\Sigma_1\hookrightarrow\Sigma_1$.
The source is $p$ on each of its two components.  All spaces are compact
manifolds in the chosen category; the source is proper and Hurewicz;
$\partial\Sigma_0=\varnothing$; there are no strictly composable pairs or
higher faces.  Thus boundary-transversality and triple compatibility
conditions impose no extra restriction.  The source fibre over $0$ is
$Y\amalg Y$, which is not a manifold.  Thus proper Hurewicz fibrancy and manifold object spaces do not imply manifold fibers.


\subsection{Connected components without local path connectivity}
\begin{counterexample}\label{ex:pseudoarc}
Local path connectivity cannot simply be omitted from
Theorem~\ref{thm:components}, even for a compact finite-dimensional
Hausdorff topological order with proper source and target.
\end{counterexample}
\begin{proof}
Choose a nondegenerate plane pseudoarc $K$, as constructed by Bing
\cite[Example~1, pp.~43--44]{Bing}. Every path in $K$ is constant:
a nonconstant path has a nondegenerate locally connected continuum as
image, which contains an arc; an arc is a decomposable subcontinuum,
contrary to hereditary indecomposability.
Choose a continuous $q:[1,\infty)\to\R^2$ with cluster set exactly
$K$ and distance to $K$ tending to zero. One construction successively
visits finite $1/n$-nets of $K$ by paths inside its connected
$1/n$-neighborhood, concatenating the finite tours. Define
\[
 C=\{(q(t),1/t):t\ge1\},\qquad P=C\cup(K\times\{0\}).
\]
This compact connected subset of $\R^3$ has finite covering dimension.
The ray endpoint $b=(q(1),1)$ is collared. Fix $a\in K\times\{0\}$.
No path starting at $a$ can enter $C$: at a first passage from height
zero to positive height, the continuous ray parameter would pass
through every sufficiently large $t$, forcing the entire tail of $q$
to converge to $a$, whereas its cluster set is the nondegenerate $K$.
A path staying in $K$ is constant.

Give $P$ the order $R=\Delta_P\cup\{(a,b)\}$.
It is a closed partial order, and the strict arrow is isolated and
clopen in $R$. Its target is a proper local embedding, including at
that isolated arrow; its image endpoint is collared in $P$.
The source has the full Hurewicz lifting property: a lifting problem
splits over the clopen inverse image of the strict-arrow component.
On that part every source path starts at $a$ and is constant, so lift
it by the constant strict arrow; on the diagonal part use the base
homotopy itself. These lifts are continuous and work for arbitrary
parameter spaces. Yet $a<b$ lies within the single connected component
of $P$.
\end{proof}

When connected components are open, the conclusion about those components
is restored if the source is also open. The image of the clopen strict-arrow subspace is then
both open and closed, since the source is proper. Once a component
contains a strict arrow with both ends in that component, the same
open-and-closed argument, with its target component fixed, supplies a
strict successor at every point there. The bounded-chain proof of
Theorem~\ref{thm:components} then excludes such an internal strict arrow.
For transitivity, the source image of arrows with target in a fixed
component is open and closed; if it meets an intermediate component it
contains that entire component. Thus the endpoint of a first arrow can
be the source of the second, and composition gives the required arrow.
A component cycle would give an internal strict comparison, unless all
its arrows were identities, by antisymmetry of the point order. This
proves the component partial order. Linear lower intervals again give
injectivity of the target between fixed components, and its restricted
properness gives a closed embedding. Alternatively one may work with
open path-connected components from the outset.

\begin{proposition}[A categorical dimension argument]\label{prop:category-order}
Let a topological category have Hausdorff object and arrow spaces,
open manifold object components, proper locally embedded target, and
Hurewicz source. Assume its strict-arrow target image is locally nowhere
dense and that its only invertible arrows are identities. Then its component relation
is a partial order. The same conclusion holds for finite-dimensional locally finite PL
polyhedral components, with PL arrow charts and PL locally embedded
target, when the strict-arrow image is locally a rare PL subset.
\end{proposition}
\begin{proof}
The identity section is closed: it is the equalizer of the identity
on the arrow space and the continuous map sending an arrow to the
identity at its source, in a Hausdorff arrow space. It is also open.
Choose a target-injective open arrow neighborhood of an identity;
continuity of the identity section gives an open object neighborhood
whose identities all lie in this arrow neighborhood. Over that object
neighborhood target injectivity leaves exactly those identities.
Thus identities, and hence strict arrows, are clopen.

Suppose a strict arrow starts and ends in the same component.
Source path lifting first moves its source into an interior coordinate
ball, while keeping the target in that component. On a smaller
contractible ball, the fibration supplies a continuous local source
section through this arrow: lift a contraction of the ball from its
center. Shrink so that the target is injective on the resulting
section. The composite embeds an open $n$-ball into an $n$-manifold
chart; if necessary compose a boundary chart with its inclusion in
$\R^n$. Invariance of domain gives an open image, contradicting the
assumed local rarity. Thus no strict arrow lies within one component.
Composition after aligning intermediate points by source path lifting
makes the component relation transitive. A cycle composes to an
identity; its reverse composite does also, so its arrows are invertible
and hence identities. This proves antisymmetry.

For PL polyhedra choose a simplex of maximal dimension in the component
and carry out the same local-section construction on its interior.
A rare PL subset in that dimension is a locally finite union of
lower-dimensional simplices. The embedded ball would be covered by
their closed inverse images; Baire's theorem makes one have interior,
and invariance of domain in the ambient simplex rules out that lower
dimension. The same contradiction proves the claim. This argument uses
the geometric rarity hypothesis explicitly; it makes no dimension
claim for arbitrary finite-dimensional compacta such as
Example~\ref{ex:pseudoarc}.
\end{proof}

\begin{center}\small
\begin{tabular}{p{.38\linewidth}p{.17\linewidth}p{.17\linewidth}p{.17\linewidth}}
\toprule
Conclusion from the stated input & TOP & PL & Smooth\\\midrule
Manifold fibers from proper Hurewicz sources & False & False & False\\
Manifold fibers from regular source bundles & False without a fiber hypothesis & True for PL bundles & True for proper facewise submersions\\
Specified manifold faces from transverse collars in manifold pieces & False & True & True\\
Component order with open path-connected components & True & True & True\\
\bottomrule
\end{tabular}
\end{center}
Here the positive statements include the horizontal and vertical boundary
formulas proved above. If a target boundary embedding is instead required
to be product-collar transverse to its own boundary image, its inverse
image of that image is the whole domain. A nonempty manifold cannot be
its own collared boundary, so that alternative condition forces the
strict arrow space to be empty. It yields only a vacuous mixed assertion.


\section{Scope of the conclusions}
The open-tube theorem concerns ordinary sheaves and a synchronized
pro-object; its minimal arities stop at triple cocycles. Higher shape
requires genuine infinity-topos descent, with hypercompletion stated
when hypercovers are used. The exact comparison conditions, and the
finite countermodels, identify why these are different assertions.
Controlled deployments recover exit categories under the recursive
proper bundle and collar hypotheses stated here. Homotopy-link germ
comparisons use the weaker tame hypotheses, while the identified
manifold faces and fibers require the stronger category-specific
conditions. The counterexamples persist even when the object strata
are well behaved, so those hypotheses cannot be supplied by sheaf
descent alone.

The general implication from a merely weak cone attachment to a
regular deployment is false. Broader equivalences obtained by weakening
the recursive collar conditions, an optimal minimal set of finite
comparison diagrams for every index category, and a minimal-dimensional
version of each manifold counterexample are not established here.

\section*{Method and records}
The mathematics and text were produced by GPT-6 Astra (ChatGPT Pro mode). Separate instances of the same model independently re-proved the open-cover descent, recursive exit, homotopy-link comparison and singular-fiber results in their stated scopes; the runtime did not expose their effort setting. The individual proof checks and computations are recorded separately. Earlier work on these questions was checked in part by Claude Opus 5.5 at effort max; those checks are distinguished from the present checks. No human has checked this paper. The text is dedicated to the public domain under CC0~1.0. The program \emph{Finite stratification models}~\cite{FiniteStratificationModels} and its check records are prepared for the proposed workbench. Permanent work records will use the registry \url{\recordurl}.

\begin{thebibliography}{99}
\bibitem{Groth} A. Grothendieck, \emph{Espaces stratifi\'es}, Montpellier archive, cotes~150--151, especially cote~151, scans~5, 7, 57--61 and~69. \url{https://grothendieck.umontpellier.fr/151.pdf}.
\bibitem{Lurie} J. Lurie, \emph{Higher Algebra}, Appendix~A, Theorem~A.9.3. \url{https://www.math.ias.edu/~lurie/papers/HA.pdf}.
\bibitem{Haine} P. J. Haine, \emph{On the homotopy theory of stratified spaces}, 27 April 2022, Theorem~3.2.4, pp.~25--26. \url{https://peterjhaine.github.io/files/Homotopy_of_stratified_spaces.pdf}.
\bibitem{DW} S. Douteau and L. Waas, \emph{From homotopy links to stratified homotopy theories}, arXiv:2112.02394v2, 2022. \url{https://arxiv.org/abs/2112.02394}.
\bibitem{AFT} D. Ayala, J. Francis and H. L. Tanaka, \emph{Local structures on stratified spaces}, Adv. Math. 307 (2017), 903--1028, Theorem~7.3.8. \url{https://arxiv.org/abs/1409.0501}.
\bibitem{Schmitt} J. Schmitt, \emph{The moduli space of curves}, lecture notes, 30 September 2020. Proposition 4.14; Sections 6.1 and 7.3. \url{https://www.math.uni-bonn.de/~schmitt/ModCurves/Script.pdf}.
\bibitem{HTT} Jacob Lurie, \emph{Higher Topos Theory} (2009). \url{https://www.math.ias.edu/~lurie/papers/HTT.pdf}.
\bibitem{Hoyois} Marc Hoyois, \emph{Higher Galois theory} (2017). \url{https://hoyois.app.uni-regensburg.de/papers/highergalois.pdf}.
\bibitem{Treumann} David Treumann, \emph{Exit paths and constructible stacks}, Compositio Mathematica 145 (2009), 1504--1532. \url{https://arxiv.org/pdf/0708.0659}.
\bibitem{Deopurkar} Ashwin Deopurkar, \emph{Torsors and gerbes} (2023). \url{https://arxiv.org/pdf/2305.15412}.
\bibitem{BGH} Clark Barwick, Saul Glasman, Peter Haine, \emph{Exodromy} (2020). \url{https://www.sas.rochester.edu/mth/sites/doug-ravenel/otherpapers/barwick-glasman-haine.pdf}.
\bibitem{Quinn} Frank Quinn, \emph{Homotopically stratified sets}, Journal of the American Mathematical Society 1 (1988), 441--499. \url{https://webhomes.maths.ed.ac.uk/~v1ranick/papers/quinn5.pdf}.
\bibitem{Grandis} Grandis, Marco, \emph{Finite sets and symmetric simplicial sets}, Theory and Applications of Categories 8 (2001), 244--252. \url{https://tac.mta.ca/tac/volumes/8/n8/n8.pdf}.
\bibitem{DHI} Dugger, Daniel, Hollander, Sharon, Isaksen, Daniel C., \emph{Hypercovers and simplicial presheaves}, Mathematical Proceedings of the Cambridge Philosophical Society 136 (2004), 9--51. \url{https://arxiv.org/pdf/math/0205027}.
\bibitem{BS} Barnea, Ilan, Schlank, Tomer M., \emph{A Projective Model Structure on Pro-Simplicial Sheaves, and the Relative {\'E}tale Homotopy Type} (2015). \url{https://arxiv.org/html/1109.5477v9}.
\bibitem{BHH} Barnea, Ilan, Harpaz, Yonatan, Horel, Geoffroy, \emph{Pro-categories in homotopy theory}, Algebraic and Geometric Topology 17 (2017), 567--643. \url{https://msp.org/agt/2017/17-1/agt-v17-n1-p17-s.pdf}.
\bibitem{Curtis} Curtis, Edward B., \emph{Simplicial Homotopy Theory}, Advances in Mathematics 6 (1971), 107--209. \url{https://www.sas.rochester.edu/mth/sites/doug-ravenel/otherpapers/simplicial-homotopy-theory.pdf}.
\bibitem{StacksDK} {The Stacks Project Authors}, \emph{The {Dold--Kan} correspondence} (2026). \url{https://stacks.math.columbia.edu/tag/019G}.
\bibitem{Kerodon} Lurie, Jacob, \emph{Kerodon: The {Ex-infinity} functor} (2026). \url{https://kerodon.net/tag/00ZA}.
\bibitem{Volpe} Volpe, Marco, \emph{Verdier duality on conically smooth stratified spaces} (2023). \url{https://arxiv.org/abs/2206.02728}.
\bibitem{Albin} Albin, Pierre, \emph{On the Hodge theory of stratified spaces} (2016). \url{https://arxiv.org/abs/1603.04106}.
\bibitem{ABP} Albin, Pierre, Banagl, Markus, Piazza, Paolo, \emph{Smooth atlas stratified spaces, K-homology orientations, and Gysin maps, Part 1} (2026). \url{https://arxiv.org/abs/2505.14952}.
\bibitem{Cannon} Cannon, James W., \emph{The recognition problem: What is a topological manifold?}, Bulletin of the American Mathematical Society 84 (1978), 832--866. \url{https://www.maths.ed.ac.uk/~v1ranick/papers/cannon2.pdf}.
\bibitem{Looijenga} Looijenga, Eduard, \emph{Moduli spaces of Riemann surfaces at Tsinghua}. \url{https://webspace.science.uu.nl/~looij101/tsinghua.pdf}.
\bibitem{Bing} Bing, R. H., \emph{Concerning hereditarily indecomposable continua}, Pacific Journal of Mathematics 1 (1951), 43--51. \url{https://msp.org/pjm/1951/1-1/pjm-v1-n1-p04-p.pdf}.
\bibitem{Zeeman} Zeeman, E. C., \emph{On the dunce hat}, Topology 2 (1964), 341--358. \url{https://webhomes.maths.ed.ac.uk/~v1ranick/papers/dunce001.pdf}.
\bibitem{Milnor} Milnor, John, \emph{Lectures on the h-Cobordism Theorem} (1965). \url{https://webhomes.maths.ed.ac.uk/~v1ranick/papers/hcobord.pdf}.
\bibitem{Ferry} Ferry, Steve, \emph{Strongly regular mappings with compact {ANR} fibers are {Hurewicz} fiberings}, Pacific Journal of Mathematics 75 (1978), 373--382. \url{https://msp.org/pjm/1978/75-2/pjm-v75-n2-p07-s.pdf}.
\bibitem{FarberSchutz} M. Farber and D. Sch\"utz,
\emph{Closed 1-forms with at most one zero}, Topology \textbf{45}
(2006), 465--473. Theorem~1, author version dated 31 August 2004,
pp.~1--2. \url{https://maths.dur.ac.uk/users/dirk.schuetz/d10.pdf}.
\bibitem{FunarPinteaZhang} L. Funar, C. Pintea and P. Zhang,
\emph{Examples of smooth maps with finitely many critical points in
dimensions $(4,3)$, $(8,5)$ and $(16,9)$}, Proc. Amer. Math. Soc.
\textbf{138} (2010), 355--365. Proposition~4.1, p.~363.
\url{https://www-fourier.univ-grenoble-alpes.fr/~funar/2010pams.pdf}.
\bibitem{CacciatoriPigola} S. L. Cacciatori and S. Pigola,
\emph{Hurewicz fibrations, almost submetries and critical points of
smooth maps}, Forum Math. \textbf{29} (2017), 751--760.
Theorems~1.1--1.2, author version arXiv:1601.07763v1, pp.~1--2.
\url{https://arxiv.org/abs/1601.07763v1}.
\bibitem{StacksGluing} The Stacks Project Authors, \emph{The Stacks Project}, Lemmas~6.33.2 and~6.33.4, Tags~00AL and~00AN, accessed 23 November 2026. \url{https://stacks.math.columbia.edu/tag/00AN}.
\bibitem{ACP} D. Abramovich, L. Caporaso and S. Payne, \emph{The tropicalization of the moduli space of curves}, Ann. Sci. \'Ec. Norm. Sup\'er. (4) 48 (2015), 765--809. \url{https://www.numdam.org/item/10.24033/asens.2258.pdf}.
\bibitem{WangZhou} Z. Wang and J. Zhou, \emph{M\"obius Inversion and Duality for Summations of Stable Graphs}, Acta Math. Sinica, English Series 42 (2026), 269--292. The cited page numbers refer to arXiv:2401.11717v1 (2024). \url{https://arxiv.org/abs/2401.11717}.
\bibitem{Tosteson} P. Tosteson, \emph{Stability in the homology of Deligne--Mumford compactifications}, arXiv:1801.03894v5 (15 November 2021), Section~3. \url{https://arxiv.org/abs/1801.03894v5}.
\bibitem{Hua} M. Hua, \emph{\'Enonc\'es des notes de travail de Grothendieck d\'emontr\'es en Lean~4, et questions ouvertes}, version~2 (5 November 2026), pp.~339--340. \url{https://hexagonmath.org/2611.00125v2}.
\bibitem{FiniteStratificationModels} GPT-6 Astra, \emph{Finite stratification models} (2026), proposed workbench program record. \url{\workbenchurl}.
\bibitem{HatcherCover} A. Hatcher, \emph{Algebraic Topology} (2002), Proposition~4G.2, p.~459. \url{https://pi.math.cornell.edu/~Hatcher/AT/ATch4.4.pdf}.
\bibitem{KerodonFiltered} J. Lurie, \emph{Kerodon}, Theorem~9.1.5.1, Filtered Colimits of Spaces, accessed 23 November 2026. \url{https://kerodon.net/tag/05XR}.
\end{thebibliography}
\end{document}
